% 第五讲：效用最大化
% 第五章统一前言包装
\newcommand{\lectureno}{第五章}
\newcommand{\lectureshort}{效用最大化}
% ============================================================
%  微观经济学理论课程讲义 —— 统一课程前言（四讲共用）
%  版式参考 Nolan Miller, Notes on Microeconomic Theory (SWP/notes3 风格)
%  编译：LuaLaTeX（两遍）
%  各讲 preamble.tex 只需设定 \lectureno / \lectureshort 后 \input 本文件。
% ============================================================
\documentclass[11pt,a4paper,fontset=none,oneside,openany]{ctexbook}

% ---------- 版式 ----------
\usepackage{geometry}
\geometry{
  a4paper,
  top=2.8cm, bottom=2.8cm,
  left=2.6cm, right=2.6cm,
  headsep=0.85cm, footskip=1.25cm
}
\setlength{\headheight}{14pt}

% ---------- 字体（全部矢量嵌入，保证 PDF 高清晰度）----------
% 注意：本机 macOS 已移除 Kaiti SC，CJK 斜体统一回退到 Songti SC；
%        Songti SC / Heiti SC 经 OSFONTDIR 指向 Supplemental 后由 luaotfload 解析。
\usepackage{fontspec}
\setmainfont{TeX Gyre Termes}[Ligatures=TeX]
\setsansfont{TeX Gyre Heros}[Ligatures=TeX,Scale=MatchLowercase]
\setmonofont{TeX Gyre Cursor}[Scale=MatchLowercase]
\setCJKmainfont{Songti SC}[BoldFont={Songti SC Bold},ItalicFont={Songti SC}]
\setCJKsansfont{Heiti SC}[BoldFont={Heiti SC}]

% ---------- 数学 ----------
\usepackage{amsmath}
\usepackage{mathtools}
\usepackage{amsthm}
\usepackage{unicode-math}
\setmathfont{TeX Gyre Termes Math}
\allowdisplaybreaks[2]
\setlength{\abovedisplayskip}{9pt plus 3pt minus 3pt}
\setlength{\belowdisplayskip}{9pt plus 3pt minus 3pt}

% ---------- 定理类环境（书稿格式：按章编号，如“定理 1.1”）----------
\theoremstyle{plain}
\newtheorem{theorem}{定理}[chapter]
\newtheorem{proposition}[theorem]{命题}
\newtheorem{lemma}[theorem]{引理}
\newtheorem{corollary}[theorem]{推论}
\theoremstyle{definition}
\newtheorem{definition}[theorem]{定义}
\newtheorem{condition}[theorem]{条件}
\newtheorem{example}[theorem]{例}
\theoremstyle{remark}
\newtheorem{remark}[theorem]{注}
% 公式、图、表按章编号（书稿格式：式 (1.1)、图 1.1、表 1.1）
\numberwithin{equation}{chapter}
\numberwithin{figure}{chapter}
\numberwithin{table}{chapter}

% 定理类环境标题后的标点：amsthm 默认加英文 "."，此处统一改为中文句号。
% 原理：amsthm 的 \@thm 每次开环境都执行 \thm@headpunct{.}，而样式宏在其之后展开，
% 故把 \thm@headpunct{。} 追加到三个样式宏末尾即可覆盖（用 \AtBeginDocument 以免被 ctex 覆盖）。
\makeatletter
\AtBeginDocument{%
  \let\WBth@plain\th@plain
  \let\WBth@definition\th@definition
  \let\WBth@remark\th@remark
  \def\th@plain{\WBth@plain\thm@headpunct{。}}%
  \def\th@definition{\WBth@definition\thm@headpunct{。}}%
  \def\th@remark{\WBth@remark\thm@headpunct{。}}%
}
\makeatother

% 无编号小标题（仿 notes3 的 “A Note on ...”）
\newcommand{\notehead}[1]{\par\medskip\noindent{\bfseries #1}\par\smallskip}

% ---------- 表格 ----------
\usepackage{booktabs}
\usepackage{array}
\renewcommand{\arraystretch}{1.25}
\setlength{\tabcolsep}{8pt}

% ---------- 图形 ----------
\usepackage{graphicx}
\usepackage{tikz}
\usepackage{pgfplots}
\pgfplotsset{compat=1.18}
\usetikzlibrary{arrows.meta,positioning,shapes.geometric,calc,patterns,decorations.pathreplacing,decorations.markings}
\usepackage{float}
\usepackage{caption}
\captionsetup{font=small,labelfont=bf,labelsep=quad,skip=4pt}

% 统一图形配色（黑白模式：黑/灰梯度，供所有 TikZ 图复用）
% —— 全局黑白化：原彩色调色板统一改为黑/灰，保证打印与投影均可读 ——
\definecolor{AxisBlue}{RGB}{0,0,0}        % 轴线 / 主曲线：黑
\definecolor{SoftBlue}{RGB}{130,130,130}  % 节点填充：中灰
\definecolor{LightRed}{RGB}{224,224,224}  % 浅填充：浅灰
\definecolor{DeepBlue}{RGB}{0,0,0}        % 描边：黑
\definecolor{MidBlue}{RGB}{70,70,70}      % 次曲线 / 网格：深灰
\tikzset{
  >=Stealth,
  every picture/.append style={font=\small},
  axisline/.style={thick,color=black,-Stealth},
  axislabel/.style={color=black},
  fillsoft/.style={fill=LightRed,fill opacity=0.6,draw=black,thick},
  nodeblue/.style={circle,fill=black!85,draw=black,thick,inner sep=2pt},
  nodewhite/.style={circle,fill=white,draw=black,thick,inner sep=2pt},
  arrblue/.style={->, thick, color=black},
  boxsoft/.style={draw=black, rounded corners=2pt, align=center, inner sep=5pt, fill=black!6},
  % 多曲线 / 多集合区分：以线型 + 标记代替颜色
  curveA/.style={thick,black,solid},
  curveB/.style={thick,black,dashed},
  curveC/.style={thick,black,dotted},
  markA/.style={mark=*,mark size=2.2pt,color=black},
  markB/.style={mark=square*,mark size=2.0pt,color=black},
  markC/.style={mark=triangle*,mark size=2.4pt,color=black},
  % 图内区域 / 点标签（白底避免与图形重叠）
  figlabel/.style={font=\scriptsize,fill=white,inner sep=1pt,outer sep=1pt,text=black},
}
% 图形说明（黑白模式下对图形区域 / 曲线 / 关键点的解释，置于图形下方）
\newcommand{\figexplan}[1]{%
  \par\smallskip\noindent{\small\itshape 图形说明：#1}\par\smallskip}

% ---------- 页眉页脚（书稿格式：左=书名，右=章次与章名）----------
\usepackage{fancyhdr}
% 章标题：居中、大号（书稿章首样式）
\CTEXsetup[format={\centering\LARGE\bfseries}]{chapter}
\CTEXsetup[beforeskip=16pt,afterskip=14pt]{chapter}
\CTEXsetup[format={\large\bfseries}]{section}
\CTEXsetup[format={\normalsize\bfseries}]{subsection}
\pagestyle{fancy}
\fancyhf{}
\providecommand{\booktitle}{微观经济学理论：可行集、偏好与效用}
\providecommand{\lectureno}{第一章}
\providecommand{\lectureshort}{课程讲义}
\fancyhead[L]{\small\itshape \booktitle}
\fancyhead[R]{\small \lectureno\quad \lectureshort}
\fancyfoot[C]{\small \thepage}
\renewcommand{\headrulewidth}{0.4pt}
\fancypagestyle{plain}{\fancyhf{}\fancyfoot[C]{\small \thepage}\renewcommand{\headrulewidth}{0pt}}

% ---------- 正文细节 ----------
\usepackage{setspace}
\setstretch{1.12}
\usepackage{microtype}
\tolerance=900
\emergencystretch=2em
\raggedbottom

% ---------- 超链接与 PDF 书签 ----------
\usepackage{hyperref}
\hypersetup{
  colorlinks=true,
  linkcolor=black,
  urlcolor=black,
  citecolor=black,
  pdfencoding=auto,
  bookmarks=true,
  pdftitle={\booktitle\quad \lectureno：\lectureshort},
  pdfsubject={微观经济学理论书稿}
}
% 章节标题中常见数学符号的安全书签映射（避免 "Improper alphabetic constant"）
\pdfstringdefDisableCommands{%
  \def\leftrightarrow{<->}%
  \def\star{*}%
  \def\S{S}%
  \def\succeq{\text{≽}}%
  \def\succ{\text{≻}}%
  \def\sim{\text{∼}}%
}

% ---------- 常用记号 ----------
\newcommand{\R}{\mathbb{R}}
\newcommand{\N}{\mathbb{N}}
\newcommand{\Z}{\mathbb{Z}}
\newcommand{\Q}{\mathbb{Q}}
\newcommand{\succeqq}{\succeq}
\newcommand{\argmax}{\operatorname*{arg\,max}}
\newcommand{\argmin}{\operatorname*{arg\,min}}
\newcommand{\st}{\text{s.t.}\ }

\setcounter{chapter}{4}
% 缺字回退：主字体缺 ①②③④ 等字形时回退到 Songti SC
\directlua{luaotfload.add_fallback("wbfallback", {"name:Songti SC:mode=node;"})}
\setmainfont{TeX Gyre Termes}[Ligatures=TeX,RawFeature={fallback=wbfallback}]

\begin{document}
\chapter{效用最大化：从可行集到 Slutsky 分解}
\noindent\textbf{本讲任务。}本章按九个小节展开：
(1) 效用最大化问题的假设与 Debreu 表示条件（含字典序反例）；
(2) 需求概念与上确界（Marshallian 需求、概念家族与存在性）；
(3) 间接效用函数的定义与结构性质；
(4) Marshallian 需求求解与 Kuhn--Tucker 条件；
(5) 边际替代率与可分性；
(6) 包络定理与 Roy identity；
(7) 支出最小化与对偶；
(8) Hicksian 需求与 Slutsky 分解；
(9) 本章小结（完整逻辑链、一页检查清单与常见陷阱）。
全文统一区分 Marshallian 需求 $x(p,w)$、Hicksian 需求 $h(p,\bar u)$、间接效用 $v(p,w)$ 与支出函数 $e(p,\bar u)$。
\tableofcontents
\clearpage
\section{效用最大化问题：假设与 Debreu 表示条件}
% ============================================================
% 第五讲 · Part 01：效用最大化问题的假设与 Debreu 表示定理
% ============================================================

\subsection{效用最大化问题：设定与假设}
\label{sec:ump-assumptions}

消费者理论的核心问题可以概括为：给定市场价格与财富，消费者在其
\textbf{负担得起}的消费方案中选出\textbf{最令其满意}者。
本章将这一直观命题形式化为\textbf{效用最大化问题}
（utility maximization problem, UMP），并系统讨论其背后的假设条件。
在正式陈述问题之前，我们先逐一审视模型的构件：
商品空间、价格与财富、偏好与效用。

\subsubsection{商品空间与消费束}
\label{subsec:commodity-space}

\begin{definition}[商品空间与消费束]\label{def:consumption-bundle}
设经济中存在 $L$ 种商品，商品空间为非负实数向量空间
\[
  X=\R^L_+=\bigl\{x=(x_1,\dots,x_L): x_i\geq 0,\ i=1,\dots,L\bigr\}.
\]
向量 $x\in X$ 称为一个\textbf{消费束}（consumption bundle），
$x_i$ 为第 $i$ 种商品的消费量。
\end{definition}

这一定义看似简单，实则隐含了若干重要的结构性假设。
我们逐条讨论其含义与作用。

\paragraph{假设一（商品种数有限）。}
$L<\infty$。商品空间是有限维欧氏空间，因而
$\R^L_+$ 中\emph{闭且有界}的集合即为\emph{紧集}（Heine--Borel 定理）。
紧性是后续证明最优解存在的关键：
连续效用函数在紧的预算集上必然达到最大值。
若商品空间为无穷维（如连续时间上的消费路径空间），
紧性论证失效，最优解可能不存在。有限维假设因此并非无关紧要的技术便利，
而是整个消费者理论存在性结论的基石。

\paragraph{假设二（商品无限可分）。}
$x_i$ 可取任意非负实数值：大米可以买 0.37 千克，汽油可以加 12.6 升。
可分性使商品空间成为连续统而非离散网格，
从而使无差异曲线成为真正的\emph{曲线}（而非离散点集），
微积分工具得以进入消费者理论。
对于天然离散的商品（汽车、住房），连续化处理是一种近似：
当商品的基本单位相对于总消费量足够微小时，
把消费量视为连续变量所引入的误差可以忽略。

\paragraph{假设三（商品的时间与状态维度）。}
同一件物品，若其交付\emph{时间}或所处的\emph{自然状态}不同，
应视为\emph{不同的}商品（Debreu, 1959）。
"今天的面包"与"明天的面包"是两种商品；
"晴天时的冰淇淋"与"雨天时的冰淇淋"也是两种商品。
只要将 $L$ 扩展为（品种数 $\times$ 时期数）或（品种数 $\times$ 状态数），
跨期选择与不确定性选择便可纳入同一静态框架，
无需另起炉灶。这一视角是后续章节讨论跨期消费与不确定性的概念基础。

\subsubsection{价格、财富与预算集}
\label{subsec:budget-set}

\paragraph{假设四（价格的线性性）。}
每种商品有唯一的价格 $p_i>0$，且价格是线性的：
购买 $x_i$ 单位商品 $i$ 的支出恰为 $p_i x_i$，不存在数量折扣、
捆绑定价或其他非线性定价结构。
线性价格保证了预算集的凸性——两个可负担的消费束，
其凸组合仍然可负担——这是后续所有切点分析与凸优化方法的几何基础。

\paragraph{假设五（消费者是价格接受者）。}
消费者将价格向量 $p$ 视为外生给定的参数，其购买量不足以影响市场价格。
这是完全竞争市场结构在个体层面的表述。
若消费者具有市场力量（双边垄断、大宗买家），
则其选择会反过来影响价格，$p$ 不再是参数，
本节的静态优化框架需要根本性的修改。

\paragraph{假设六（完全信息）。}
消费者知晓全部价格、了解自身偏好，交易中不存在信息不对称、
搜寻成本或认知偏差。在这一假设下，
预算集被完整地感知，不存在"以为买得起实际买不起"的情形。

在假设四至假设六下，消费者的货币约束由价格向量 $p$ 与财富 $w$ 刻画。

\begin{definition}[预算集]\label{def:budget-set}
给定价格向量 $p\in\R^L$ 与财富 $w\geq 0$，\textbf{预算集}为
\[
  B(p,w)=\bigl\{x\in\R^L_+: p\cdot x\leq w\bigr\},
  \qquad
  p\cdot x=\sum_{i=1}^{L}p_i x_i .
\]
当 $p\gg 0$ 时，预算集是紧的：它显然是闭集；
又由于每个分量满足 $x_i\leq w/p_i$，它还是有界的。
\end{definition}

\subsubsection{效用最大化问题}
\label{subsec:ump}

消费者在其预算集内选择最令其满意的消费束。
满意程度由\textbf{偏好}刻画，偏好是消费者理论中先于效用的原始概念。
然而，偏好作为二元关系不便于计算，
我们希望将其转化为数值目标——这正是 \S\ref{sec:debreu}
中 Debreu 表示定理的主题。
当偏好可由连续效用函数 $u:\R^L_+\to\R$ 表示时，
消费者的选择行为归结为如下最优化问题：

\begin{definition}[效用最大化问题，UMP]\label{def:ump}
设 $u:\R^L_+\to\R$ 为表示偏好的连续效用函数。
给定 $p\gg 0$ 与 $w\geq 0$，\textbf{效用最大化问题}为
\begin{equation}\label{eq:ump}
  \max_{x\in B(p,w)} u(x).
\end{equation}
其解集称为 \textbf{Marshallian 需求}：
\[
  X(p,w)=\arg\max_{x\in B(p,w)} u(x).
\]
由于预算集 $B(p,w)$ 是紧集且 $u$ 连续，由 Weierstrass 定理，
解集 $X(p,w)$ 非空；若偏好还是凸的，则 $X(p,w)$ 是凸集
（从而可能包含多个点）；若偏好严格凸，则 $X(p,w)$ 退化为单点。
这些结论的完整证明将在 \S\ref{sec:demand-concepts} 系统给出，
届时还区分\emph{需求对应}与\emph{需求函数}两个概念。
\end{definition}

图~\ref{fig:ump-tangency} 展示了两种商品情形下 UMP 的几何结构：
预算线（预算集的上边界）与无差异曲线相切处即为最优消费束。
在切点上，无差异曲线的斜率（边际替代率）恰好等于预算线的斜率（价格比），
这一条件将在下一讲推导 Kuhn--Tucker 条件时严格化。

\begin{figure}[H]
\centering
\begin{tikzpicture}[scale=0.85, line cap=round, line join=round]
  % 坐标轴
  \draw[axisline] (0,0) -- (8.4,0) node[below, axislabel] {$x_1$};
  \draw[axisline] (0,0) -- (0,6.8) node[left, axislabel] {$x_2$};
  % 预算线
  \draw[thick] (0,4.8) -- (6.4,0);
  \node[figlabel] at (3.6,3.25) {预算线};
  % 预算集填充
  \fill[gray!10] (0,0) -- (6.4,0) -- (0,4.8) -- cycle;
  % 无差异曲线族（自低到高）
  \draw[curveA, thick] plot [smooth, tension=0.8]
    coordinates {(1.2,5.8) (2.0,4.2) (3.2,3.1) (4.6,2.5) (6.4,2.0)};
  \draw[curveA, thick] plot [smooth, tension=0.8]
    coordinates {(0.6,4.0) (1.4,2.6) (2.8,1.7) (4.4,1.2) (6.4,0.8)};
  \draw[curveA, thick] plot [smooth, tension=0.8]
    coordinates {(0.3,2.2) (1.0,1.4) (2.2,0.8) (3.6,0.5) (5.4,0.3)};
  % 最优点
  \fill (2.85,2.66) circle (2.5pt);
  \draw (2.85,2.66) node[above right] {$x^*$};
  % 切点虚线引导
  \draw[dashed, thin] (2.85,2.66) -- (2.85,0) node[below, figlabel] {$x_1^*$};
  \draw[dashed, thin] (2.85,2.66) -- (0,2.66) node[left, figlabel] {$x_2^*$};
\end{tikzpicture}
\caption{效用最大化：预算线与最高可达无差异曲线相切于 $x^*$}
\label{fig:ump-tangency}
\end{figure}
\figexplan{预算线由 $p\cdot x=w$ 给出，其斜率为 $-p_1/p_2$。
无差异曲线族由低到高排列，消费者希望达到尽可能高的无差异曲线，
但受预算约束限制。最优消费束 $x^*$ 位于预算线与某条无差异曲线的切点处：
在切点上，无差异曲线的斜率（边际替代率）恰好等于价格比 $p_1/p_2$。
低于 $x^*$ 的无差异曲线上的消费束虽然可负担，但未充分利用预算；
高于 $x^*$ 的无差异曲线上的消费束则超出预算集。
切点的存在性与唯一性依赖于偏好连续性与凸性假设，
将在后续各节严格讨论。}
\label{page:ump-figure}

\subsection{Debreu 表示定理：从偏好到连续效用}
\label{sec:debreu}

上一节提到，为了将偏好转化为可计算的数值目标，
我们需要一个\textbf{效用函数}来表示偏好。
但一个自然的问题是：什么样的偏好才能被效用函数表示？
本节讨论的 Debreu 表示定理给出了答案：
偏好必须\textbf{完备}、\textbf{传递}且\textbf{连续}。
更重要的是，这三个条件缺一不可——
每个条件单独去掉后，都存在无法被任何效用函数表示的反例。

\subsubsection{效用表示的定义}
\label{subsec:representation-def}

\begin{definition}[效用表示]\label{def:representation}
设 $\succeq$ 为商品空间 $X$ 上的偏好关系。
函数 $u:X\to\R$ \textbf{表示}偏好 $\succeq$，若
\[
  x\succeq y \iff u(x)\geq u(y), \qquad \forall\, x,y\in X.
\]
\end{definition}

这一定义的要点是"$\iff$"：$u$ 不仅需将更受偏好的消费束映到更大的数值，
还必须将无差异的消费束映到相等的数值。
定义中"$\Rightarrow$"保证保序性，"$\Leftarrow$"保证忠实性——
效用函数不多不少地刻画了偏好。

\subsubsection{Debreu 定理的陈述}
\label{subsec:debreu-statement}

\begin{theorem}[Debreu 表示定理，1954]\label{thm:debreu-representation}
设 $\succeq$ 为 $X=\R^L_+$ 上的偏好关系。若 $\succeq$ 满足：
\begin{enumerate}
  \item \textbf{完备性}：$\forall x,y\in X$，$x\succeq y$ 或 $y\succeq x$；
  \item \textbf{传递性}：$x\succeq y$ 且 $y\succeq z$ $\Rightarrow$ $x\succeq z$；
  \item \textbf{连续性}：$\forall y\in X$，
    上等值集 $\{x\in X: x\succeq y\}$ 与
    下等值集 $\{x\in X: y\succeq x\}$ 均为闭集，
\end{enumerate}
则存在连续函数 $u:\R^L_+\to\R$ 表示 $\succeq$。
\end{theorem}

\begin{proof}[证明思路（构造骨架）]
完整的严格证明技术性较强（Debreu, 1954, 1964），此处给出构造的核心思路。

\textbf{第一步（可数稠密子集）。}
取 $D=\{d_1,d_2,\dots\}$ 为 $\R^L_+$ 中坐标全为有理数的消费束全体。
$D$ 是可数集且在 $\R^L_+$ 中稠密——任何消费束 $x$ 附近都存在 $D$ 中的点。

\textbf{第二步（二进制编码）。}
对每个消费束 $x\in\R^L_+$，定义
\[
  u(x)=\sum_{n=1}^{\infty}\frac{1}{2^n}\,\mathbf{1}[x\succeq d_n],
\]
其中 $\mathbf{1}[\cdot]$ 为示性函数。直觉是：将 $D$ 中的消费束
按照偏好顺序"编码"为二进制小数，
若 $x$ 弱偏好于 $d_n$，则在第 $n$ 位计入 $2^{-n}$。

\textbf{第三步（保序性验证）。}
若 $x\succeq y$，则对每个 $n$，$y\succeq d_n$ 蕴含 $x\succeq d_n$
（由传递性），故 $u(x)\geq u(y)$。
反之，若 $x\succ y$（严格偏好），则由连续性，
集合 $\{z: z\succ y\}$ 是包含 $x$ 的开集；
由稠密性，存在 $d_n$ 落在该开集中且 $x\succeq d_n$
（对单调偏好可取 $d_n$ 为 $x$ 的有理近似"下方"点；
一般情形需更细致的论证，见 Debreu, 1954）。
于是 $u(x)\geq u(y)+2^{-n}>u(y)$，保序性成立。

\textbf{第四步（连续性验证）。}
由偏好的连续性，$x^k\to x$ 时，
集合 $\{n: x^k\succeq d_n\}$ 的"边界变化"是受控的——
若 $x\succ d_n$，则在 $x$ 的某个邻域内恒有 $z\succ d_n$；
若 $d_n\succ x$，同理。仅当 $x\sim d_n$（恰好无差异）时，
$x^k$ 与 $d_n$ 的比较结果才可能跳变，
而由于 $D$ 可数、每个跳变对 $u$ 的影响至多为 $2^{-n}$，
总影响可通过控制序列的逼近速度加以消除。
故 $u$ 连续。
\end{proof}

\subsubsection{三个条件为何缺一不可}
\label{subsec:debreu-necessity}

Debreu 定理的三个条件各自承担不可替代的角色。
我们逐条讨论去掉某个条件后会发生什么。

\paragraph{去掉完备性。}
若偏好不完备，存在消费束 $x,y$ 无法比较：
既非 $x\succeq y$，也非 $y\succeq x$。
但任何实值函数 $u$ 都必然给出 $u(x)\geq u(y)$ 或 $u(y)\geq u(x)$ 之一
（实数系的完全有序性），
因此 $u$ 必然为 $x,y$ 强行指定一个"虚假的"排序。
不完备的偏好无法被任何实值函数忠实地表示。
例如，消费者可能面对两种截然不同的生活方式
（在城市工作与在农村务农）时确实无法排序，
强行赋予数值便是武断。

\paragraph{去掉传递性。}
若偏好不传递，可能出现\textbf{偏好循环}：
$x\succ y$、$y\succ z$ 但 $z\succ x$。
此时若存在表示函数 $u$，则 $u(x)>u(y)>u(z)>u(x)$，
矛盾。更直观地说，持有 $x$ 的消费者会
以正价格卖出 $x$ 买入 $y$（因为他偏好 $y$），
再卖出 $y$ 买入 $z$，再卖出 $z$ 买入 $x$，
循环一轮后回到初始持有，却每一步都支付了正的价格——
这就是著名的\textbf{抽水泵悖论}（money pump argument）。
具有循环偏好的消费者在市场中会被无限剥削。

\paragraph{去掉连续性。}
若偏好不连续，即使在偏好完备且传递的"理性"情况下，
也可能不存在任何（哪怕不连续的）实值效用表示。
这类反例中最经典的是\textbf{字典序偏好}，
它具有完备性与传递性、满足局部非饱和，
却因不连续而无法被任何函数表示。
下一节将完整地展示这一反例的构造、性质及其不可表示性的证明。

\subsection{字典序偏好：完备、传递但不可表示}
\label{sec:lexicographic}

\subsubsection{定义与基本性质}
\label{subsec:lex-def}

\begin{definition}[字典序偏好]\label{def:lexicographic}
在 $\R^2_+$ 上定义偏好 $\succeq_L$ 如下：
$(x_1,x_2)\succeq_L (y_1,y_2)$ 当且仅当
\[
  x_2>y_2,\qquad\text{或者}\qquad x_2=y_2 \ \text{且}\ x_1\geq y_1.
\]
换言之，消费者\textbf{绝对优先}考虑第二坐标 $x_2$：
只要 $x_2$ 更大，无论 $x_1$ 是多少都严格偏好；
仅当 $x_2$ 相同时，才退而比较 $x_1$。
如同字典排序先看首字母、再看次字母，故名\textbf{字典序}
（lexicographic order）。
\end{definition}

\begin{proposition}[字典序偏好的完备性与传递性]\label{prop:lex-rational}
$\succeq_L$ 是完备且传递的（因而是一个理性偏好关系）。
\end{proposition}

\begin{proof}
\textbf{完备性。}任取 $x,y\in\R^2_+$。
若 $x_2\neq y_2$，则 $x_2>y_2$ 或 $y_2>x_2$，二者必居其一，
故 $x\succeq_L y$ 或 $y\succeq_L x$。
若 $x_2=y_2$，则比较 $x_1$ 与 $y_1$：由实数的完全有序性，
$x_1\geq y_1$ 或 $y_1\geq x_1$ 必居其一，结论同样成立。

\textbf{传递性。}设 $x\succeq_L y$ 且 $y\succeq_L z$。
由定义有 $x_2\geq y_2$ 且 $y_2\geq z_2$，故 $x_2\geq z_2$。
分两种情形：
若 $x_2>z_2$，则 $x\succ_L z$，结论成立。
若 $x_2=z_2$，则必须有 $x_2=y_2=z_2$
（中间不等式全部取等号）。
此时 $x\succeq_L y$ 蕴含 $x_1\geq y_1$，
$y\succeq_L z$ 蕴含 $y_1\geq z_1$，
故 $x_1\geq z_1$，从而 $x\succeq_L z$。
\end{proof}

\begin{proposition}[字典序偏好的局部非饱和性]\label{prop:lex-nonsatiation}
$\succeq_L$ 满足局部非饱和：对任意 $x\in\R^2_+$ 与任意 $\varepsilon>0$，
存在 $y\in B(x,\varepsilon)\cap\R^2_+$ 使 $y\succ_L x$。
\end{proposition}

\begin{proof}
取 $y=(x_1, x_2+\varepsilon/2)$。则 $\|y-x\|=\varepsilon/2<\varepsilon$
且 $y_2=x_2+\varepsilon/2>x_2$，故 $y\succ_L x$。
\end{proof}

由此可见，字典序偏好满足 Debreu 定理条件中的完备性与传递性，
甚至满足局部非饱和这一附加性质。
下面证明，它恰恰在\textbf{连续性}上失败，
并因此无法被任何实值效用函数表示。

\subsubsection{不连续性}
\label{subsec:lex-discontinuity}

\begin{proposition}[字典序偏好不连续]\label{prop:lex-discontinuous}
$\succeq_L$ 的上等值集不是闭集，因而 $\succeq_L$ 不连续。
具体地：取 $x^*=(0,1)$，则集合
$\{(x_1,x_2): (x_1,x_2)\succeq_L (0,1)\}$
包含点列 $(1/n, 0)$（$n=1,2,\dots$），却不包含其极限 $(0,0)$。
\end{proposition}

\begin{proof}
对每个 $n\geq 1$，点 $(1/n,0)$ 满足 $1/n>0=x_1^*$，
故 $(1/n,0)\succ_L (0,1)$，即 $(1/n,0)$ 属于 $(0,1)$ 的上等值集。
然而当 $n\to\infty$ 时 $(1/n,0)\to(0,0)$，
而 $(0,0)$ 的第二坐标 $0<1$，故 $(0,0)\prec_L(0,1)$，
即 $(0,0)$ 不在上等值集中。
上等值集包含了点列却不含其极限，故不是闭集。
由连续性定义（上等值集必须为闭集），$\succeq_L$ 不连续。
\end{proof}

直观地说，字典序偏好的"层级"沿 $x_2$ 方向排列，
而层内的排序由 $x_1$ 决定。当 $x_2$ 连续变化时，
偏好在 $x_2$ 轴的每个"截面"上都发生剧烈的跳变：
相邻截面之间的偏好关系突然翻转。
这种无穷多层级之间的跳变，正是连续性失效的根源。
图~\ref{fig:lex-structure} 展示了这一结构。

\begin{figure}[H]
\centering
\begin{tikzpicture}[scale=0.85, line cap=round, line join=round]
  % 坐标轴
  \draw[axisline] (0,0) -- (8.6,0) node[below, axislabel] {$x_1$};
  \draw[axisline] (0,0) -- (0,6.8) node[left, axislabel] {$x_2$};
  % 两条水平层级线
  \draw[dashed, thick, gray] (0.4,1.8) -- (8.2,1.8);
  \node[figlabel] at (7.5,1.52) {$x_2=b$};
  \draw[dashed, thick, gray] (0.4,4.0) -- (8.2,4.0);
  \node[figlabel] at (7.5,4.3) {$x_2=b'$};
  % 上方层级更优的标注
  \draw[arrblue, thick] (1.8,2.1) -- (1.8,3.7);
  \node[figlabel] at (1.8,2.9) {$\succ$};
  % 下方层内排序
  \draw[arrblue, thick] (4.8,1.72) -- (6.4,1.72);
  \node[figlabel] at (5.6,1.38) {$x_1\uparrow$ 层内更优};
  % 层内取值示例
  \fill (3.6,1.8) circle (2pt);
  \draw (3.6,1.8) node[above left, inner sep=2pt] {\small $y$};
  \fill (6.0,1.8) circle (2pt);
  \draw (6.0,1.8) node[above left, inner sep=2pt] {\small $z$};
  \node[figlabel] at (4.6,2.15) {$y\prec z$};
  % 整体说明
  \node[align=center, font=\small] at (4.3,5.8)
    {每个水平层级 $x_2=b$ 内部，按 $x_1$ 严格排序；\\
     层级之间按 $x_2$ 严格排序};
  \node[align=center, font=\small, text=gray] at (4.3,6.5)
    {不存在无差异曲线：每个无差异类退化为单点};
\end{tikzpicture}
\caption{字典序偏好的层级结构：竖直方向（$x_2$）优先，水平方向（$x_1$）次级}
\label{fig:lex-structure}
\end{figure}
\figexplan{字典序偏好按第二坐标 $x_2$ 分为无穷多个"层级"（图中虚线所示），
层级之间按 $x_2$ 严格排序（箭头 $\succ$ 指向更高层级）。
在每个层级内部（即固定 $x_2=b$），消费束再按第一坐标 $x_1$ 严格排序
（如图中 $y\prec z$）。
由于每个无差异类退化为单个点——不存在两个不同的消费束是无差异的——
传统的无差异曲线完全消失。
层内与层间均为\emph{严格}排序，没有"平台"，也没有等高线。
这种结构导致偏好在 $x_2$ 方向上无法连续变化：
当 $x_2$ 越过一个层级的边界时，
偏好关系瞬间发生不连续的跳变（命题~\ref{prop:lex-discontinuous}），
最终导致任何实值函数都无法忠实编码这一偏好结构
（定理~\ref{thm:lex-norepresentation}）。}

\subsubsection{不可表示性}
\label{subsec:lex-nonrep}

我们已经看到字典序偏好不完备于 Debreu 定理的连续性条件
（命题~\ref{prop:lex-discontinuous}），
因此定理不适用。
一个更深刻的事实是：不存在\emph{任何}实值函数——
无论连续与否——能够表示字典序偏好。

\begin{theorem}[字典序偏好不可表示]\label{thm:lex-norepresentation}
不存在任何函数 $u:\R^2_+\to\R$ 表示字典序偏好 $\succeq_L$。
\end{theorem}

\begin{proof}
反设存在表示函数 $u:\R^2_+\to\R$。
对每个 $t\in\R_+$，考虑两个消费束 $(0,t)$ 与 $(1,t)$。
由定义（$x_2$ 坐标相同、$x_1$ 更大），$(1,t)\succ_L(0,t)$，
故
\[
  u(1,t)>u(0,t).
\]
令 $I_t$ 为开区间
\[
  I_t=\bigl(u(0,t),\, u(1,t)\bigr),\qquad t\in\R_+ .
\]
每个 $I_t$ 非空。关键观察是：
\textbf{不同 $t$ 对应的区间两两不相交}。
设 $s<t$。由字典序定义，$(1,s)\prec_L(0,t)$
（因为 $t>s$，第二坐标更大者胜出），
故 $u(1,s)<u(0,t)$。
于是 $I_s$ 的右端点 $u(1,s)$ 小于 $I_t$ 的左端点 $u(0,t)$，
即 $I_s$ 完全位于 $I_t$ 的左侧，二者不相交。

现在，每个非空开区间 $I_t$ 都至少包含一个有理数
（有理数在 $\R$ 中稠密）。
为每个 $t\in\R_+$ 选定一个 $q(t)\in I_t\cap\mathbb{Q}$。
由两两不相交性，$s\neq t$ 蕴含 $q(s)\neq q(t)$，
故 $t\mapsto q(t)$ 是从 $\R_+$ 到 $\mathbb{Q}$ 的单射。
但 $\R_+$ 不可数而 $\mathbb{Q}$ 可数，矛盾。
因此表示函数不存在。
\end{proof}

\begin{proof}[另一种证明视角：层级太多，实数装不下]
定理~\ref{thm:lex-norepresentation} 的证明还可以从
"基数冲突"的角度直观理解。$\R_+$ 中有不可数无穷多个层级 $t$，
每个层级内部的 $x_1$ 维度又各自构成一个不可数无穷的严格排序链条。
若要将所有这些信息压缩到一条实数轴上，
就必须为每个层级 $t$ 在 $\R$ 中预留一段非退化的"空间"（区间 $I_t$），
而这些空间两两不重叠。
但实数轴上两两不相交的非空开区间至多只有可数多个
（每个区间可分配一个互不相同的有理数）。
不可数多个层级无法容纳于可数多个区间中，
故表示不可能存在。
\end{proof}

\subsubsection{对 Debreu 定理条件的回顾}
\label{subsec:lex-lesson}

字典序反例的价值在于，它揭示了 Debreu 定理中
\textbf{连续性}条件的实质经济内涵。
完备性要求消费者对任何两个消费束都能排序——
这是关于消费者"判断能力"的假设；
传递性要求偏好不出现循环——
这是关于偏好"内在一致性"的假设，
并与市场中的套利论证（抽水泵悖论）相联系；
而连续性要求偏好在消费束的微小变动下不发生"跳变"——
它排除了一种极端的心理结构：
消费者将某个商品维度视为具有压倒一切的绝对优先级
（无论代价多小，多一点 $x_2$ 总比任何数量的 $x_1$ 更好）。
这种"词典式"的优先结构虽然在形式上是理性的
（完备、传递、甚至局部非饱和），
但它过于极端，以至于实数系统——
一维的、连续的、可分的——装不下这种二维的严格层级结构。

最后值得指出的是，连续性条件在经济学中通常被视为无害的技术性假设：
"消费束的微小变化不应导致偏好剧烈反转"与直觉高度一致。
字典序反例恰恰说明，这一看似温和的条件
排除了什么样的病态结构，
从而为后续章节中以效用函数为出发点的所有分析扫清了障碍。
一旦 $u$ 存在且连续，我们便可放心地运用微积分、凸分析等工具
来研究消费者的最优化行为。

% ============================================================
% 第五讲 · Part 02：需求概念与最优值的存在性
% ============================================================

\section{需求概念与上确界：Marshallian 需求、概念家族与存在性}
\label{sec:demand-concepts}

效用最大化问题的解——\textbf{需求}——是消费者理论的核心输出。
本节首先厘清\textbf{需求对应}与\textbf{需求函数}的区别，
然后系统梳理围绕 UMP 构建的整个\textbf{需求概念家族}。

\subsection{Marshallian 需求：对应与函数}
\label{subsec:correspondence-vs-function}

\begin{definition}[Marshallian（Walrasian）需求对应]\label{def:marshallian-demand}
设 $u:\R^L_+\to\R$ 为效用函数，$p\gg 0$ 为价格向量，$w\ge 0$ 为财富。
\textbf{Marshallian 需求对应}定义为 UMP 的最优解集：
\[
  X(p,w)=\arg\max\bigl\{\,u(x)\;:\;x\in\R^L_+,\;p\cdot x\le w\,\bigr\}.
\]
当 $X(p,w)$ 对给定的 $(p,w)$ 恰好只含一个元素时，
记 $x(p,w)$ 为其中唯一的消费束，称为\textbf{Marshallian 需求函数}。
\end{definition}

$X(p,w)$ 之所以称为\textbf{对应}（correspondence）而非\textbf{函数}（function），
是因为从 $(p,w)$ 到消费束的映射可能是\emph{多值}的：
同一个价格—财富组合下，可能存在多个互不相同的最优消费束。
何时出现多值？何时退化为单值？
答案是：取决于偏好的\textbf{凸性}。

\begin{itemize}
\item \textbf{偏好凸但非严格凸}（如完全替代偏好）：
  无差异曲线可能含有\emph{线性段}。
  当预算线的斜率恰好等于线性段的斜率时，
  线性段上\emph{每一点}都与预算线"相切"，
  整个线段上的消费束同时达到最大效用，
  此时 $X(p,w)$ 是一段线段（多值）。
\item \textbf{偏好严格凸}（如 Cobb--Douglas 偏好）：
  无差异曲线处处严格弯曲，不含线性段。
  每条无差异曲线与预算线至多只有一个接触点，
  因而 $X(p,w)$ 退化为单点集 $\{x(p,w)\}$，
  需求成为\emph{函数}。
\end{itemize}

图~\ref{fig:convex-vs-strict} 直观展示两种情形的差异。

\begin{figure}[H]
\centering
\begin{tikzpicture}[scale=0.78, line cap=round, line join=round]
  % ===== 左图：凸偏好（非严格凸），需求对应为一段线段 =====
  \begin{scope}
    % 坐标轴
    \draw[axisline] (0,0) -- (6.2,0) node[below, axislabel] {$x_1$};
    \draw[axisline] (0,0) -- (0,5.8) node[left, axislabel] {$x_2$};
    % 预算线：从 (0,4.5) 到 (5.5,0)，斜率 -4.5/5.5 ≈ -0.818
    \draw[thick] (0,4.5) -- (5.5,0);
    \node[figlabel, anchor=south west] at (0.1,4.55) {$w/p_2$};
    \node[figlabel, anchor=north east] at (5.45,-0.05) {$w/p_1$};
    % 无差异曲线：上弯部分 + 线性段（与预算线重合）+ 下弯部分
    \draw[curveA] (0.4, 5.3) .. controls (0.6, 4.2) and (1.0, 3.1) .. (1.6, 2.55);
    \draw[curveA, line width=1.8pt] (1.6, 2.55) -- (3.8, 0.75);
    \draw[curveA] (3.8, 0.75) .. controls (4.3, 0.44) and (4.8, 0.32) .. (5.6, 0.22);
    % 最优线段端点标记
    \fill (1.6, 2.55) circle (2.2pt);
    \fill (3.8, 0.75) circle (2.2pt);
    \node[figlabel, anchor=south east] at (1.5, 2.7) {$x^a$};
    \node[figlabel, anchor=north west] at (3.9, 0.6) {$x^b$};
    % 标注：整段最优
    \node[align=center, font=\footnotesize] at (2.7, 3.6)
      {凸（非严格凸）偏好：\\ $X(p,w)=[x^a,x^b]$ 为线段};
    % 一条更低的无差异曲线（虚线）
    \draw[curveB] (0.5, 3.6) .. controls (1.0, 2.2) and (2.0, 1.3) .. (3.2, 0.85);
    \draw[curveB] (3.2, 0.85) -- (4.6, 0.30);
    \draw[curveB] (4.6, 0.30) .. controls (4.9, 0.20) and (5.2, 0.15) .. (5.7, 0.10);
    \node[curveB] at (5.2, 0.75) {};
  \end{scope}
  % ===== 右图：严格凸偏好，唯一最优点 =====
  \begin{scope}[xshift=7.4cm]
    % 坐标轴
    \draw[axisline] (0,0) -- (6.2,0) node[below, axislabel] {$x_1$};
    \draw[axisline] (0,0) -- (0,5.8) node[left, axislabel] {$x_2$};
    % 预算线：从 (0,4.5) 到 (5.5,0)
    \draw[thick] (0,4.5) -- (5.5,0);
    \node[figlabel, anchor=south west] at (0.1,4.55) {$w/p_2$};
    \node[figlabel, anchor=north east] at (5.45,-0.05) {$w/p_1$};
    % 严格凸无差异曲线：x₂ = k/x₁ 型，与预算线相切
    % 切点在 (2.5, 2.45)（预算线上），k = 2.5×2.45 = 6.125
    \draw[curveA] plot [smooth, tension=0.8]
      coordinates {(1.2, 5.1) (1.5, 4.08) (2.0, 3.06) (2.5, 2.45) (3.0, 2.04) (3.5, 1.75) (4.2, 1.46) (5.2, 1.18)};
    % 最优点标记
    \fill (2.5, 2.45) circle (2.2pt);
    \node[figlabel, anchor=north west] at (2.65, 2.30) {$x^* = x(p,w)$};
    % 标注
    \node[align=center, font=\footnotesize] at (2.7, 3.9)
      {严格凸偏好：\\ $X(p,w)=\{x^*\}$ 单点};
    % 一条更低的无差异曲线（虚线）：k = 3.5
    \draw[curveB] plot [smooth, tension=0.8]
      coordinates {(1.0, 3.5) (1.4, 2.5) (2.0, 1.75) (2.8, 1.25) (3.5, 1.0) (4.5, 0.78) (5.5, 0.64)};
  \end{scope}
\end{tikzpicture}
\caption{偏好凸性与需求对应的单值性：左图为凸但非严格凸偏好，
  需求对应为一段线段；右图为严格凸偏好，需求退化为唯一的切点。}
\label{fig:convex-vs-strict}
\end{figure}

\figexplan{
左图：无差异曲线（实线）在中间段为直线（线性段），且线性段恰好落在预算线上。
此时线性段两端点 $x^a$ 与 $x^b$ 之间（含端点）的每个消费束，
其所在的无差异曲线都"同时"与预算线重合，效用值相同且均达到最优，
因此 Marshallian 需求对应 $X(p,w)$ 是整个线段 $[x^a,x^b]$。
虚线为一条更低的无差异曲线，其上的消费束同样可负担但未达最优。
右图：无差异曲线处处严格弯曲（不含线性段），
与预算线只在唯一切点 $x^*$ 处接触。
该无差异曲线是预算集内所能达到的最高无差异曲线，
切点 $x^*$ 是唯一的最优消费束，故需求对应退化为需求函数 $x(p,w)$。
虚线为一条更低的无差异曲线，不与预算线相切。}

\subsection{需求概念家族一览}
\label{subsec:demand-family}

围绕 UMP，微观经济学构建了一族紧密关联的概念。
表~\ref{tab:demand-family} 给出完整的概览。

\begin{definition}[间接效用函数、Hicksian 需求与支出函数]\label{def:indirect-hicksian}
设 $u:\R^L_+\to\R$ 连续，$p\gg 0$，$w\ge 0$，目标效用 $\bar u$ 满足可达性。
\begin{enumerate}
\item \textbf{间接效用函数}（indirect utility function）：
  Marshallian 需求所对应的效用水平
  \[
    v(p,w)=\max\bigl\{\,u(x)\;:\;x\in\R^L_+,\;p\cdot x\le w\,\bigr\}
      = u\bigl(x(p,w)\bigr).
  \]
\item \textbf{Hicksian（补偿）需求}（compensated demand）：
  在\emph{达到目标效用}的前提下，最小化支出
  \[
    h(p,\bar u)=\arg\min\bigl\{\,p\cdot x\;:\;x\in\R^L_+,\;u(x)\ge\bar u\,\bigr\}.
  \]
\item \textbf{支出函数}（expenditure function）：
  上述最小化问题的最优值
  \[
    e(p,\bar u)=\min\bigl\{\,p\cdot x\;:\;x\in\R^L_+,\;u(x)\ge\bar u\,\bigr\}
      = p\cdot h(p,\bar u).
  \]
\end{enumerate}
\end{definition}

\begin{table}[H]
\centering
\small
\begin{tabular}{@{}llll@{}}
\toprule
\textbf{概念} & \textbf{数学定义} & \textbf{固定什么} & \textbf{回答什么问题} \\
\midrule
Marshallian 需求 & $\max u(x)$ s.t.\ $p\cdot x\le w$ & 价格 $p$、财富 $w$ & 给定预算，买什么？ \\
$x(p,w)$ & 的最优解 & & \\
\addlinespace[3pt]
间接效用 $v(p,w)$ & 同上问题的最优值 & 价格 $p$、财富 $w$ & 给定预算，最高效用是多少？ \\
\addlinespace[3pt]
Hicksian 需求 $h(p,\bar u)$ & \multicolumn{1}{l}{$\min p\cdot x$ s.t.\ $u(x)\ge\bar u$} & 价格 $p$、目标效用 $\bar u$ & 达到目标效用，买什么？ \\
 & 的最优解 & & \\
\addlinespace[3pt]
支出函数 $e(p,\bar u)$ & 同上问题的最优值 & 价格 $p$、目标效用 $\bar u$ & 达到目标效用，最少花多少钱？ \\
\bottomrule
\end{tabular}
\caption{需求概念家族一览}
\label{tab:demand-family}
\end{table}

这一家族的核心联系由如下\textbf{对偶关系}给出：
\begin{equation}\label{eq:duality-relations}
  h(p,\bar u) = x\bigl(p,\,e(p,\bar u)\bigr),
  \qquad
  e\bigl(p,\,v(p,w)\bigr) = w.
\end{equation}
第一式含义：最小化支出得到的最优消费束，
恰好等于"给消费者恰好够用的财富 $e(p,\bar u)$ 让其自己最大化效用"的结果；
第二式含义：最大化效用之后再问"达到 $v(p,w)$ 至少要花多少钱"，
答案就是 $w$。
这两条对偶关系将在后续章节（支出最小化问题）中系统展开，
并构成 Slutsky 分解的理论基础：
\textbf{价格变化的总效应}可以拆分为
\emph{替代效应}（Hicksian 需求的变化，效用不变）与
\emph{财富效应}（由实际购买力变化引起）。
没有 Hicksian 需求与支出函数这对工具，
Slutsky 分解便无从谈起。

\begin{remark}[为什么需要两套概念？]
Marshallian 需求回答"给定预算怎么花"；
Hicksian 需求回答"给定效用怎么最省"。
前者是消费者真实面对的问题，与数据直接对应；
后者是理论分析的理想化问题，结构更简洁。
两者通过 \eqref{eq:duality-relations} 的对偶关系紧密联结：
$v$ 是 $e$ 的"反函数"（在适当的意义下），
$h$ 可由 $e$ 经 Shephard 引理直接恢复，
正如 $x$ 可由 $v$ 经 Roy 恒等式恢复。
对偶结构的完整刻画留待后续章节。
\end{remark}
\label{sec:sup-max}

\subsection{sup 与 max 的严格辨析}
\label{subsec:sup-vs-max}

设 $S\subseteq\R$ 为一实数集，$f$ 为定义在某集合 $D$ 上的实值函数。

\begin{definition}[上确界与最大值]\label{def:sup-max}
\begin{enumerate}
\item $M\in\R$ 是 $S$ 的\textbf{最大值}，若 (i) $M\in S$ 且 (ii) $s\le M$ 对一切 $s\in S$ 成立。
\item $\bar s\in\R$ 是 $S$ 的\textbf{上确界}（supremum），若 (i) $s\le\bar s$ 对一切 $s\in S$（上界），
  且 (ii) 对任意 $\bar s'<\bar s$，存在 $s\in S$ 使 $s>\bar s'$（最小上界）。
\end{enumerate}
最大值必定是上确界；但上确界可能不属于 $S$，此时最大值不存在。
\end{definition}

\begin{example}[sup 存在而 max 不存在]\label{ex:sup-not-max}
\leavevmode
\begin{enumerate}
\item 开区间 $S=(0,1)$：$\sup S=1\notin S$，最大值不存在。
\item 定义在开集 $D=(0,1)$ 上的函数 $f(x)=x$：
  $\sup_{x\in D} f(x)=1$，但 $f$ 在 $D$ 上没有最大值
  （因为 $x=1\notin D$，而 $f(x)$ 可以无限接近但永远不等于 $1$）。
\end{enumerate}
在 UMP 中，预算集 $B(p,w)$ 是\emph{闭}的
（它包含边界 $\{x:p\cdot x=w\}$），
这保证了"最优解不会被遗漏在可行集之外"。
\end{example}

上确界与最大值何时一致？
这正是 \textbf{Weierstrass 定理}回答的问题：
在紧集上，连续函数的最大值必然达到。

\subsection{Weierstrass 定理：完整证明}
\label{subsec:weierstrass}

\begin{theorem}[Weierstrass 最大值定理]\label{thm:weierstrass}
设 $K\subseteq\R^L$ 为非空紧集，$f:K\to\R$ 为连续函数。
则 $f$ 在 $K$ 上达到最大值与最小值：
存在 $x^*,x_*\in K$ 使 $f(x^*)\ge f(x)\ge f(x_*)$ 对一切 $x\in K$ 成立。
\end{theorem}

\begin{proof}
只证明最大值的情形（最小值通过对 $-f$ 取最大值得到）。

\textbf{第一步：$f$ 在 $K$ 上有界。}
反设 $f$ 在 $K$ 上无上界。
则对每个 $n\in\N$，存在 $x^n\in K$ 使 $f(x^n)>n$。
由于 $K$ 是紧集（$\R^L$ 中紧 $\Leftrightarrow$ 闭且有界 $\Leftrightarrow$ 序列紧），
序列 $\{x^n\}$ 存在收敛子列 $\{x^{n_k}\}$，其极限 $x^*=\lim_k x^{n_k}\in K$。
由 $f$ 连续，$f(x^{n_k})\to f(x^*)$。
但 $f(x^{n_k})>n_k\to\infty$，与 $f(x^*)<\infty$ 矛盾。
故 $f$ 在 $K$ 上有上界，从而 $M=\sup_{x\in K}f(x)<\infty$ 存在。

\textbf{第二步：上确界可以达到。}
由上确界的定义（最小上界性质），
对每个 $n\in\N$ 存在 $x^n\in K$ 使 $f(x^n)>M-\frac{1}{n}$。
再次由 $K$ 的序列紧性，
存在收敛子列 $\{x^{n_k}\}$，其极限 $x^*\in K$。
由 $f$ 连续，
\[
  f(x^*)=\lim_{k\to\infty}f(x^{n_k}).
\]
另一方面，$M-\frac{1}{n_k}<f(x^{n_k})\le M$ 对一切 $k$ 成立，
令 $k\to\infty$，由夹逼定理得
\[
  f(x^*)=M=\sup_{x\in K}f(x).
\]
又 $x^*\in K$，故 $M\in f(K)$，
即上确界作为函数值被 $x^*$ 达到，$x^*$ 即为最大点。
\end{proof}

\begin{remark}[证明结构解读]
证明的核心是\textbf{序列紧性}（sequential compactness）：
紧集中的任意序列必有收敛于该集合内的子列。
第一步排除了目标值发散到无穷的可能（保证 sup 有限），
第二步将"无限逼近"转化为"确实达到"。
两步共同的关键前提是：$K$ 是\textbf{闭}的（子列极限不会跑出可行集）
且\textbf{有界}（序列不会跑向无穷远）。
这解释了为何在 UMP 中，
预算集 $B(p,w)$ 的闭性与有界性缺一不可：
闭性要求预算约束包含等号（$\le w$ 而非 $<w$），
有界性要求所有价格为正（$p\gg 0$）。
\end{remark}

\subsection{需求对应的三条基本定理}
\label{subsec:demand-three-theorems}

将 Weierstrass 定理应用于预算集，
我们立即得到需求理论的三条基本定理。

\begin{theorem}[需求对应非空]\label{thm:demand-nonempty}
设 $u:\R^L_+\to\R$ 连续，$p\gg 0$，$w\ge 0$。
则 Marshallian 需求对应 $X(p,w)$ 非空（且为紧集）。
\end{theorem}

\begin{proof}
预算集
\[
  B(p,w)=\bigl\{x\in\R^L_+:p\cdot x\le w\bigr\}
\]
是紧集：
\emph{闭}，因为它是不等式 $p\cdot x-w\le 0$ 的零点下方集
（连续函数的闭集的原像）与非负象限之交；
\emph{有界}，因为 $p_i>0$ 且 $x_i\ge 0$ 对每个 $i$ 成立，
故 $x_i\le w/p_i$。
又 $0\in B(p,w)$，故 $B(p,w)$ 非空。
由 Weierstrass 定理（定理~\ref{thm:weierstrass}），
连续函数 $u$ 在非空紧集 $B(p,w)$ 上达到最大值，
即 $X(p,w)=\arg\max_{B(p,w)}u\neq\varnothing$。

$X(p,w)$ 的紧性：$X(p,w)$ 是闭集
（它是 $\{x\in B(p,w):u(x)=v(p,w)\}$，
即连续函数 $u$ 的水平集与闭集 $B(p,w)$ 之交，故为闭集）
且含于有界集 $B(p,w)$，故 $X(p,w)$ 为紧集。
\end{proof}

\begin{theorem}[凸偏好下需求对应为凸集]\label{thm:demand-convex}
设 $u:\R^L_+\to\R$ 连续且拟凹（代表凸偏好），
$p\gg 0$，$w\ge 0$。
则 $X(p,w)$ 是凸集。
\end{theorem}

\begin{proof}
设 $x^a, x^b\in X(p,w)$，$t\in[0,1]$。
记 $x^t=tx^a+(1-t)x^b$。

\textbf{$x^t$ 可行}：由 $p\cdot x^a\le w$、$p\cdot x^b\le w$ 与 $p$ 的线性性，
\[
  p\cdot x^t = t\,p\cdot x^a + (1-t)\,p\cdot x^b \le t\,w+(1-t)\,w = w,
\]
且 $x^t\in\R^L_+$（非负象限是凸集），故 $x^t\in B(p,w)$。

\textbf{$x^t$ 最优}：由 $u$ 的拟凹性，
\[
  u(x^t) \ge \min\bigl\{u(x^a),\,u(x^b)\bigr\}.
\]
而 $x^a, x^b$ 均为最优解，故 $u(x^a)=u(x^b)=v(p,w)$，其中
\[
  v(p,w)=\max_{x\in B(p,w)}u(x)
\]
为间接效用函数。
于是 $u(x^t)\ge v(p,w)$。
但 $v(p,w)$ 是 $u$ 在 $B(p,w)$ 上的最大值，
故 $u(x^t)\le v(p,w)$。
综合得 $u(x^t)=v(p,w)$，即 $x^t\in X(p,w)$。
\end{proof}

\begin{theorem}[严格凸偏好下需求唯一]\label{thm:demand-unique}
设 $u:\R^L_+\to\R$ 连续且严格拟凹（代表严格凸偏好），
$p\gg 0$，$w\ge 0$。
则 $X(p,w)$ 至多含有一个元素，
即 Marshallian 需求退化为单值的需求函数 $x(p,w)$。
\end{theorem}

\begin{proof}
反设 $X(p,w)$ 含有两个不同元素 $x^a\neq x^b$。
由定理~\ref{thm:demand-nonempty}，$X(p,w)\neq\varnothing$，
故 $v(p,w)$ 有限且 $u(x^a)=u(x^b)=v(p,w)$。

取 $t\in(0,1)$，记 $x^t=tx^a+(1-t)x^b$。
由定理~\ref{thm:demand-convex} 的证明（第一步），
$x^t\in B(p,w)$。
由 $u$ 的\emph{严格}拟凹性，
\[
  u(x^t) > \min\bigl\{u(x^a),\,u(x^b)\bigr\} = v(p,w).
\]
但 $v(p,w)$ 是 $u$ 在 $B(p,w)$ 上的最大值，
$u(x^t)>v(p,w)$ 与此矛盾。
故 $X(p,w)$ 不可能含有两个不同元素。
\end{proof}

\subsection{Walras 定律：预算约束取等}
\label{subsec:walras-law}

在局部非饱和偏好下，消费者永远不会"留钱不花"：
最优消费束必定恰好花完财富。
这一性质称为 Walras 定律，
它是需求理论中最基础的可检验含义之一。

\begin{theorem}[Walras 定律]\label{thm:walras}
设偏好局部非饱和（对任意 $x$ 及任意 $\varepsilon>0$，
存在 $y\in\R^L_+$ 使 $\|y-x\|<\varepsilon$ 且 $y\succ x$），
$u$ 为其连续表示，$p\gg 0$，$w>0$。
则对任意 $x^*\in X(p,w)$，有
\[
  p\cdot x^* = w.
\]
\end{theorem}

\begin{proof}
由 $x^*\in X(p,w)$，显然 $p\cdot x^*\le w$。
反设 $p\cdot x^*<w$，即 $x^*$ 位于预算集的\emph{内部}
（严格满足不等式）。
令
\[
  \varepsilon = \frac{w - p\cdot x^*}{2\|p\|} > 0.
\]
由局部非饱和，存在 $y\in\R^L_+$ 使 $\|y-x^*\|<\varepsilon$ 且 $u(y)>u(x^*)$。

考察 $y$ 的可负担性。由 Cauchy--Schwarz 不等式，
\[
  p\cdot y = p\cdot x^* + p\cdot(y-x^*)
  \le p\cdot x^* + \|p\|\,\|y-x^*\|
  < p\cdot x^* + \|p\|\,\varepsilon
  = p\cdot x^* + \frac{w-p\cdot x^*}{2}
  < w.
\]
故 $y\in B(p,w)$ 且 $u(y)>u(x^*)$。
这与 $x^*$ 是 $B(p,w)$ 上的最大点矛盾。
故必有 $p\cdot x^*=w$。
\end{proof}

图~\ref{fig:walras-interior} 给出几何直观：
若消费束位于预算集内部（$p\cdot x<w$），
则存在包含它的一个小邻域仍然在预算集内；
由局部非饱和，邻域内总存在严格更好的消费束；
因此内部点不可能是最优的——最优点必须落在预算集的边界上，
即恰好花完财富。

\begin{figure}[H]
\centering
\begin{tikzpicture}[scale=0.82, line cap=round, line join=round]
  % 坐标轴
  \draw[axisline] (0,0) -- (7.2,0) node[below, axislabel] {$x_1$};
  \draw[axisline] (0,0) -- (0,6.2) node[left, axislabel] {$x_2$};
  % 预算线
  \draw[thick] (0,5) -- (6,0);
  \node[figlabel, anchor=south west] at (0.1,5.05) {$w/p_2$};
  \node[figlabel, anchor=north east] at (5.95,-0.05) {$w/p_1$};
  % 内点 x（非最优候选）
  \fill (2.5, 2.0) circle (2.2pt);
  \node[figlabel, anchor=south east] at (2.3, 2.2) {$x$};
  % 邻域（小圆）
  \draw[dashed] (2.5, 2.0) circle (1.2);
  \node[figlabel, anchor=north west] at (3.7, 3.1) {\small $\varepsilon$-邻域仍在 $B(p,w)$ 内};
  % 邻域内的更优点 y
  \fill (3.0, 2.5) circle (2.2pt);
  \node[figlabel, anchor=south west] at (3.2, 2.7) {$y\succ x$};
  % 箭头指向更优点
  \draw[arrblue, thick] (2.65, 2.15) -- (2.85, 2.35);
  % 预算边界上的最优解
  \fill (4.0, 1.67) circle (2.2pt);
  \node[figlabel, anchor=west] at (4.2, 1.67) {$x^*\in X(p,w)$：必须在边界上};
  % 过 x 的无差异曲线（示意，虚线）
  \draw[curveB] (1.2, 4.2) .. controls (2.0, 3.0) and (3.0, 2.6) .. (4.2, 2.4);
  \draw[curveB] (4.2, 2.4) .. controls (5.0, 2.2) and (5.8, 2.1) .. (6.8, 2.0);
\end{tikzpicture}
\caption{Walras 定律的几何直观：内部点不可能是最优解。
  若 $p\cdot x<w$，则 $x$ 的一个小邻域整体落在预算集内，
  局部非饱和保证邻域内存在严格更好的消费束 $y$，
  因而 $x$ 不可能最优。最优解必须落在预算线上，
  即 $p\cdot x^*=w$。}
\label{fig:walras-interior}
\end{figure}

\figexplan{
预算线（实线）将可行区域（左下方）与不可行区域分开。
候选消费束 $x$（左点）严格位于预算集内部，即 $p\cdot x<w$。
以 $x$ 为中心、适当小的半径画一个小圆（虚线），
由 $p\cdot x<w$ 且不等式是严格的，可将半径选得足够小使整个圆仍在预算集内。
由局部非饱和假设，圆内必存在严格优于 $x$ 的消费束 $y$（右点），
且 $y$ 同样可负担。
因此 $x$ 被支配，不可能是最优解。
反观位于预算线上的点 $x^*$（右侧圆点），
预算线之外不存在任何可负担的消费束，
故它可能成为最优解。
这从几何上解释了 Walras 定律：
在局部非饱和偏好下，消费者将财富恰好用尽。
}

\begin{remark}[Walras 定律的含义与应用]
Walras 定律 $p\cdot x^*=w$ 有两个重要推论。
其一，\textbf{需求的零次齐次性}：
当价格与财富同比例变化 $(p,w)\to(\lambda p,\lambda w)$（$\lambda>0$）时，
预算集 $B(p,w)$ 不变，因此 $X(p,w)=X(\lambda p,\lambda w)$。
消费者不受"货币幻觉"影响，只关心相对价格与实际购买力。
其二，\textbf{加总与市场层面}：
将所有消费者的 Walras 定律加总，
可得市场总超额需求的价值恒为零。
这一性质在一般均衡理论中用于证明：
若 $L-1$ 个市场出清，则第 $L$ 个市场自动出清
（Walras 定律在一般均衡中的版本）。
\end{remark}

\begin{remark}[对偶预告]
本节建立的三条基本定理（非空性、凸值性、单值性）与 Walras 定律
构成了需求理论的公理化骨架。
后续章节将在此骨架上叠加微积分结构：
对偶关系 \eqref{eq:duality-relations} 将给出
Hicksian 与 Marshallian 需求之间的精确换算（Slutsky 方程），
间接效用函数 $v$ 的可微性将产生 Roy 恒等式，
支出函数 $e$ 的可微性将产生 Shephard 引理。
这些结果共同构成现代需求理论的完整体系。
\end{remark}

\section{间接效用函数：定义、结构性质与例题}
% ============================================================
% 第五讲 5.3：间接效用函数 v(p,w)
% 内容：定义与 sup/max 之辨、零次齐次性、对 w 严格递增、
%       对 p 非增、Berge 连续性、对 (p,w) 拟凸、
%       偏好性质对照表、Cobb--Douglas 例题逐条验证
% ============================================================

\subsection{间接效用函数：定义、上确界与最大值}

前两节回答了消费者\emph{选什么}（需求对应 $X(p,w)$）与\emph{何时能选}（存在性与唯一性）；\emph{如何求解}（Kuhn--Tucker 条件）将在 \S 5.4 展开。本节转向一个更紧凑的对象：把效用最大化问题的\emph{最优值本身}当作价格与财富的函数，研究它的结构性质。这个“值函数”就是间接效用函数 $v(p,w)$。它是消费者理论的核心中间产品：Slutsky 分解、Roy 恒等式、补偿变差与等价变差、价格指数理论，都建立在 $v$ 的性质之上。

记号约定：本节统一以 $w\ge 0$ 表示财富，价格 $p\in\R^L_{++}$，消费束 $x\in X=\R^L_+$，预算集
\[
B(p,w)=\{x\in X:\ p\cdot x\le w\}.
\]

\begin{definition}[间接效用函数]
\label{def:v}
给定效用函数 $u:X\to\R$，间接效用函数定义为
\[
v(p,w)\;=\;\sup_{x\in B(p,w)}u(x),
\]
即预算可行集上效用的上确界；Walrasian 需求对应为 $X(p,w)=\arg\max_{x\in B(p,w)}u(x)$。
\end{definition}

\subsubsection{为什么用 \texorpdfstring{$\sup$}{sup} 而不是 \texorpdfstring{$\max$}{max}}

\begin{remark}[上确界的三重严格性]
用 $\sup$ 而非 $\max$ 并非记号洁癖，而是有三层理由。
\textbf{其一（良定性）}：只要 $B(p,w)$ 非空，$\sup$ 在广义实数中恒有定义，\emph{完全不依赖最优解是否存在}；而 $\max$ 要求 $\arg\max$ 非空。换言之，$v$ 的定义把“值”与“选择”分离开：即使消费者暂时选不出最优束，“他最好能好到什么程度”仍有明确数值。若 $B(p,w)=\varnothing$（例如消费集取开集 $\R^L_{++}$ 且 $w$ 过小），我们约定 $v(p,w)=-\infty$。
\textbf{其二（对极限封闭）}：上确界与单调极限运算交换：若 $B^n\uparrow B$，则 $\sup_{B^n}u\uparrow\sup_Bu$。这使 $v$ 在逼近论证（如本节连续性证明中“近似最优束”的构造）里可以自由取极限，而 $\max$ 无此便利。
\textbf{其三（诚实的陈述）}：把存在性问题与值函数问题分开陈述，可以准确指出每条性质究竟依赖什么——本节将看到，$v$ 的齐次性、单调性与拟凸性\emph{不需要}解的存在性，而“$\sup$ 化为 $\max$”与连续性才需要 $u$ 连续。
\end{remark}

\begin{proposition}[$\sup$ 化为 $\max$ 的充分条件]
\label{prop:max}
若 $u$ 在 $X$ 上连续，$p\gg0$ 且 $0\le w<\infty$，则 $B(p,w)$ 非空紧，于是 $v(p,w)=\max_{x\in B(p,w)}u(x)$，且 $X(p,w)$ 非空、紧。
\end{proposition}
\begin{proof}
非空性：$0\in B(p,w)$。有界性：对 $x\in B(p,w)$，有 $0\le x_i\le w/p_i$（$p_i>0$），故 $B(p,w)\subseteq\prod_i[0,w/p_i]$。闭性：$B(p,w)$ 是闭半空间 $\{x:p\cdot x\le w\}$ 与闭锥 $\R^L_+$ 之交。由 Heine--Borel 定理，$B(p,w)$ 紧。连续函数在紧集上达到最大值（Weierstrass），故最大值存在；$X(p,w)$ 是紧集的水平集切片，是闭子集，从而紧。
\end{proof}

\begin{example}[两种失效方式]
\label{ex:fail}
$\sup$ 未必是 $\max$，且闭性与连续性缺一不可。
（i）\emph{可行集不闭}：取 $B=(0,1)$、$u(x)=x$，则 $\sup_Bu=1$，但无最大值点。
（ii）\emph{效用不连续}：取 $B=[0,1]$，令 $u(x)=x$（$x<1$）而 $u(1)=0$。可行集紧，但 $\sup_Bu=1$ 仍不可达，$\arg\max=\varnothing$。
在标准环境（$u$ 连续、$p\gg0$、闭消费锥）下，命题~\ref{prop:max} 的条件自动满足；但在应用设定中——整数商品、不可分选择（住址、职业）、带最低生存约束 $\underline x$ 的消费集——必须逐条检查闭性与连续性，$\sup$ 记号正是为这些情形保留的诚实性。
\end{example}

\subsection{间接效用的结构性质}

本节系统建立 $v$ 的五条代数与拓扑性质：零次齐次、对 $w$ 严格递增、对 $p_i$ 非增、连续、对 $(p,w)$ 拟凸。请特别注意各条性质的\emph{假设账单}：齐次性、对 $p$ 的非增性与拟凸性是预算约束线性结构的直接后果，对\emph{任意}偏好成立；对 $w$ 的严格递增需要局部非饱和；连续性需要 $u$ 连续（通过 Berge 最大值定理）。

\subsubsection{零次齐次性：无货币幻觉}

\begin{theorem}[零次齐次性]
\label{thm:homog}
对任意 $t>0$，有 $B(tp,tw)=B(p,w)$，从而
\[
v(tp,tw)=v(p,w).
\]
\end{theorem}
\begin{proof}
设 $x\in X$。因 $t>0$，
\[
tp\cdot x\le tw\iff p\cdot x\le w,
\]
故 $x\in B(tp,tw)$ 当且仅当 $x\in B(p,w)$，两集合作为 $X$ 的子集完全相同。同一集合上取 $\sup$，值相等。
\end{proof}

\notehead{经济解释：无货币幻觉}
零次齐次性即“无货币幻觉”（absence of money illusion）：物价与财富同步放大 $t$ 倍时，\emph{真实选择集} $B$ 一动不动，最优效用自然不变。直觉有三层。
（i）\emph{计价单位无关}：把货币单位从元改为分，所有数字乘 $100$，可行的消费计划集合不变；币制改革（如欧元转换、删去两个零）不改变任何人的福利。
（ii）\emph{名义与实际}：纯通胀若伴随收入的完全指数化，则 $(p,w)$ 同比例变动，$v$ 不变——福利只取决于\emph{相对}价格与\emph{实际}财富。这是需求系统最早的可检验含义之一：名义变量的同比例变化不应出现在需求方程中。
（iii）\emph{理论便利}：齐次性允许规范化（normalize），例如固定 $w=1$ 或把某一价格定为 $1$（计价物），把 $(p,w)$ 的 $L+1$ 个自由度减少一个。后文的对偶图示与比较静态反复使用这一自由度。

\subsubsection{对财富的严格递增性}

\begin{theorem}[对 $w$ 严格递增]
\label{thm:incw}
设 $u$ 连续且偏好局部非饱和，$p\gg0$。则对 $0\le w<w'$，有
\[
v(p,w)<v(p,w').
\]
\end{theorem}
\begin{proof}
第一步（非严格单调）。$w\le w'$ 时 $B(p,w)\subseteq B(p,w')$，上确界随集合扩大而不减，故 $v(p,w)\le v(p,w')$。

第二步（严格化）。设 $w<w'$。由命题~\ref{prop:max}，$v(p,w)=u(x^*)$ 在某个 $x^*\in B(p,w)$ 处达到。记松弛量 $\delta=w'-p\cdot x^*\ge w'-w>0$。局部非饱和给出：对任意 $\varepsilon>0$，存在 $y\in X$，$\|y-x^*\|<\varepsilon$ 且 $y\succ x^*$。取 $\varepsilon$ 充分小，使 $\|y-x^*\|<\varepsilon$ 蕴含 $p\cdot y<p\cdot x^*+\delta\le w'$（内积 $y\mapsto p\cdot y$ 连续）。于是该 $y$ 满足 $y\in B(p,w')$ 且
\[
v(p,w')\ \ge\ u(y)\ >\ u(x^*)\ =\ v(p,w).\qedhere
\]
\end{proof}

\begin{remark}[局部非饱和不可省]
若无局部非饱和，更多财富可能买不到更好的东西。取 $u(x)=-(x_1-1)^2-(x_2-1)^2$：点 $\bar x=(1,1)$ 是严格优于一切的饱和点（satiation point），在其任意邻域内不存在更优的点，局部非饱和在 $\bar x$ 处失效。当 $w\ge p_1+p_2$ 时 $\bar x$ 可被负担，$v(p,w)=u(\bar x)$ 对一切更大的 $w$ 保持平坦。局部非饱和恰好排除“钱多到花不出去”的角落。另外注意定理证明中“$v(p,w)$ 被达到”来自 $u$ 的连续性；若 $u$ 不连续（例~\ref{ex:fail}（ii）型），$\sup$ 不可达时严格性论证失效——这也是 $\sup$ 与 $\max$ 差异的又一体现。
\end{remark}

\notehead{经济解释：财富的边际效用}
$\partial v/\partial w$（存在时）是财富的边际效用。定理~\ref{thm:incw} 说明在局部非饱和下它严格为正：任何财富增量都能换到严格更优的消费计划。这是福利度量政策的基石：补偿变差 $CV=e(p',v(p,w))-w$、等价变差 $EV=w-e(p',v(p,w))$ 之所以有意义，正因为 $v$ 对 $w$ 严格单调，从而 $e(p,\cdot)$ 与 $v(p,\cdot)$ 互为（广义）反函数（见本章对偶一节）。

\begin{figure}[H]
\centering
\begin{tikzpicture}[scale=1.0]
  \draw[axisline] (0,0) -- (6.6,0) node[below,axislabel] {$w$};
  \draw[axisline] (0,0) -- (0,4.4) node[left,axislabel] {$v(p,w)$};
  % v at price p (solid, strictly increasing, diminishing slope)
  \draw[curveA] plot[smooth] coordinates {(0,0) (1.2,1.85) (2.4,2.6) (3.6,3.05) (4.8,3.3) (5.8,3.45)};
  % v at higher price p' >= p (dashed, weakly lower)
  \draw[curveB] plot[smooth] coordinates {(0,0) (1.2,1.1) (2.4,1.6) (3.6,1.95) (4.8,2.2) (5.8,2.35)};
  % marked points on curveA
  \node[nodewhite] at (1.2,1.85) {};
  \node[nodewhite] at (3.6,3.05) {};
  \draw[thin,dashed,black!60] (1.2,1.85) -- (1.2,0);
  \draw[thin,dashed,black!60] (3.6,3.05) -- (3.6,0);
  \draw[thin,dashed,black!60] (1.2,1.85) -- (0,1.85);
  \draw[thin,dashed,black!60] (3.6,3.05) -- (0,3.05);
  \node[figlabel] at (1.2,-0.22) {$w_1$};
  \node[figlabel] at (3.6,-0.22) {$w_2$};
  \node[figlabel] at (-0.15,1.85) {$v_1$};
  \node[figlabel] at (-0.15,3.05) {$v_2$};
  \node[figlabel] at (6.0,3.45) {$v(p,\cdot)$};
  \node[figlabel] at (6.0,2.35) {$v(p',\cdot)$};
\end{tikzpicture}
\caption{固定价格下间接效用随财富严格递增}
\figexplan{实线为固定 $p$ 时 $w\mapsto v(p,w)$ 的图像：曲线自 $(0,u(0))$ 出发（图中规范化 $u(0)=0$）并严格上升（定理~\ref{thm:incw}），示意性地画出边际财富效用递减。虚线为更高价格 $p'\ge p$ 下的间接效用：整条曲线被压低且不再过原点上方，体现对价格的非增性（命题~\ref{prop:decp}）。点 $(w_1,v_1)$、$(w_2,v_2)$ 标注财富增加带来的效用严格增益。}
\end{figure}

\subsubsection{对价格的非增性}

\begin{proposition}[对 $p_i$ 非增]
\label{prop:decp}
设 $p'\ge p$（分量式），则 $B(p',w)\subseteq B(p,w)$，从而
\[
v(p',w)\le v(p,w).
\]
\end{proposition}
\begin{proof}
对 $x\in B(p',w)$：$p\cdot x\le p'\cdot x\le w$（因 $p\le p'$ 且 $x\ge0$），故 $x\in B(p,w)$。集合缩小，上确界不增。
\end{proof}

非增性不需要任何偏好假设：涨价只会砍掉消费机会。什么时候\emph{严格}下降？答案取决于该商品是否“被需要”。

\begin{proposition}[严格递减的充分条件]
\label{prop:decpstrict}
设 $u$ 连续且严格递增（$x\ge y,\ x\ne y\Rightarrow u(x)>u(y)$）。若每个最优束 $x\in X(p,w)$ 都满足 $x_i>0$，则对任意 $p'\ge p$、$p'_i>p_i$，有 $v(p',w)<v(p,w)$。
\end{proposition}
\begin{proof}
由命题~\ref{prop:max} 取 $x'\in X(p',w)$，则 $v(p',w)=u(x')\le v(p,w)$（命题~\ref{prop:decp}）。反设 $u(x')=v(p,w)$。则 $x'\in B(p',w)\subseteq B(p,w)$ 且取得 $B(p,w)$ 上的最大值，即 $x'\in X(p,w)$，由假设 $x'_i>0$。又严格递增偏好下预算必绑定：若 $p\cdot x'<w$，取 $\varepsilon>0$ 充分小使 $x'+\varepsilon e_i\in B(p,w)$，则 $u(x'+\varepsilon e_i)>u(x')$，与最优性矛盾；故 $p\cdot x'=w$。于是
\[
p'\cdot x'\;=\;p\cdot x'+(p'_i-p_i)\,x'_i\;=\;w+(p'_i-p_i)x'_i\;>\;w,
\]
与 $x'\in B(p',w)$ 矛盾。故 $v(p',w)<v(p,w)$。
\end{proof}

\notehead{经济解释：通胀的福利损失}
对 $p$ 非增、对 $w$ 严格递增合起来刻画了“价格上涨即福利受损”。严格性的条件同样有内容：若商品 $i$ 在最优束中消费量为零（消费者视其为无关品甚至厌恶品），其涨价不损失任何效用。命题~\ref{prop:decpstrict} 的逆否形式在应用中常用：观测到 $v$ 对 $p_i$ 严格下降，等价于该商品在最优点处被正消费。此外，“涨价压低 $v$”正是补偿变差度量通胀损失的出发点：把消费者补回原效用水平所需的财富 $e(p',v(p,w))-w$，依赖的全部结构就是 $v$ 对 $p$ 的单调性与对 $w$ 的严格单调性。

\subsubsection{连续性：Berge 最大值定理}

前三条性质是纯代数的。连续性则需要拓扑工具——Berge 最大值定理，它是整个比较静态分析的合法性来源：若参数微扰能让最优值跳变，一切边际分析（导数、包络、Roy 恒等式）都无从谈起。

\begin{lemma}[预算对应的连续性]
\label{lem:budget}
在 $\Theta=\R^L_{++}\times\R_+$ 上，预算对应 $(p,w)\mapsto B(p,w)$ 非空、紧值、上半连续且下半连续。
\end{lemma}
\begin{proof}
非空（$0\in B$）与紧值（命题~\ref{prop:max} 的证明）已述。

上半连续。设 $(p^n,w^n)\to(p,w)$，$x^n\in B(p^n,w^n)$。因 $p^n\to p\gg0$，存在 $c>0$ 与 $N$ 使 $n\ge N$ 时 $\min_i p^n_i\ge c$；于是 $0\le x^n_i\le w^n/c$ 有界。取任一收敛子列 $x^{n_k}\to x$；在 $p^{n_k}\cdot x^{n_k}\le w^{n_k}$ 中令 $k\to\infty$（内积连续）得 $p\cdot x\le w$，且 $x\ge0$，故 $x\in B(p,w)$。即：任意序列的任意收敛子列的极限都落回极限预算集，这是紧值情形下上半连续的序列刻画。

下半连续。设 $(p^n,w^n)\to(p,w)$，任取 $x\in B(p,w)$。若 $x=0$，取 $x^n=0$ 即可。设 $x\ne0$。
（i）若 $p\cdot x<w$（约束松弛），则 $n$ 充分大时 $p^n\cdot x<w^n$，取 $x^n\equiv x$。
（ii）若 $p\cdot x=w$（约束绑定，此时 $w>0$），令
\[
\lambda^n=\min\Bigl\{1,\ \frac{w^n}{p^n\cdot x}\Bigr\},\qquad x^n=\lambda^n x .
\]
则 $x^n\ge0$ 且 $p^n\cdot x^n=\lambda^n(p^n\cdot x)\le w^n$，故 $x^n\in B(p^n,w^n)$；又 $w^n/(p^n\cdot x)\to w/(p\cdot x)=1$，故 $\lambda^n\to1$，$x^n\to x$。
\end{proof}

\begin{theorem}[Berge 最大值定理（度量情形，序列形式）]
\label{thm:berge}
设 $\Theta\subseteq\R^K$，$u:X\to\R$ 连续，$F:\Theta\to\to X$ 非空、紧值，且满足
（i）\emph{上半连续}：对任意 $\theta^n\to\theta$ 与 $x^n\in F(\theta^n)$，$\{x^n\}$ 有界，且其每个收敛子列的极限都属于 $F(\theta)$；
（ii）\emph{下半连续}：对任意 $\theta^n\to\theta$ 与 $x\in F(\theta)$，存在 $x^n\in F(\theta^n)$ 使 $x^n\to x$。
则值函数 $\varphi(\theta)=\max_{x\in F(\theta)}u(x)$ 在 $\Theta$ 上连续；最优值对应 $\Psi(\theta)=\arg\max_{x\in F(\theta)}u$ 非空、紧值、上半连续。
\end{theorem}
\begin{proof}
非空与紧值：紧集上连续函数达到最大（Weierstrass），$\Psi(\theta)$ 是紧集的闭子集。

$\varphi$ 的上半连续。设 $\theta^n\to\theta$，取 $x^n\in\Psi(\theta^n)$，$\varphi(\theta^n)=u(x^n)$。取子列实现 $\limsup$：设 $u(x^{n_k})\to\limsup_n\varphi(\theta^n)$；由（i），$\{x^{n_k}\}$ 有再收敛子列 $x^{n_{k_j}}\to x\in F(\theta)$；$u$ 连续给出
\[
\limsup_n\varphi(\theta^n)=\lim_j u(x^{n_{k_j}})=u(x)\le\varphi(\theta).
\]

$\varphi$ 的下半连续。任给 $\varepsilon>0$，取 $x\in F(\theta)$ 使 $u(x)>\varphi(\theta)-\varepsilon$（最大值达到）。由（ii）取 $x^n\in F(\theta^n)$，$x^n\to x$，则
\[
\liminf_n\ \varphi(\theta^n)\ \ge\ \lim_n\ u(x^n)\ =\ u(x)\ >\ \varphi(\theta)-\varepsilon .
\]
令 $\varepsilon\downarrow0$。两个半连续性合并即连续性。

$\Psi$ 的上半连续。设 $\theta^n\to\theta$，$x^n\in\Psi(\theta^n)$。由（i）取收敛子列 $x^{n_k}\to x\in F(\theta)$。对任意 $z\in F(\theta)$，由（ii）取 $z^n\in F(\theta^n)$，$z^n\to z$；最优性给出 $u(x^{n_k})\ge u(z^{n_k})$，令 $k\to\infty$ 得 $u(x)\ge u(z)$。故 $x\in\Psi(\theta)$。
\end{proof}

\begin{corollary}[间接效用的连续性]
\label{cor:cont}
设 $u$ 连续。则 $v$ 在 $\R^L_{++}\times\R_+$ 上连续，$X(p,w)$ 非空、紧值、上半连续；若偏好还严格凸，则 $X$ 单值，从而 $X$ 是连续的需求函数 $x(p,w)$。
\end{corollary}
\begin{proof}
引理~\ref{lem:budget} 验证定理~\ref{thm:berge} 的条件（i）（ii），结论直接代入。单值情形：上半连续 + 单值 = 连续。
\end{proof}

\notehead{上、下半连续的直觉}
\emph{上半连续}（upper hemicontinuity）说的是“极限可行仍可行”：参数序列下的可行点，其极限点不会跌出极限预算集之外——预算约束的图 $\{(p,w,x):p\cdot x\le w\}$ 是闭集，取极限保号。直觉上，约束条件在极限中不会“凭白放宽”。\emph{下半连续}（lower hemicontinuity）说的是“可行可被逼近”：极限预算集里的任何计划，都能被邻近参数下的可行计划逼近——约束不会“突然塌缩”。证明中的构造 $x^n=\lambda^nx$ 正是经济学内容本身：预算线上的消费计划按\emph{实际购买力} $w^n/(p^n\cdot x)$ 等比例缩放，恰好适配新价格与新财富。两个方向的失效都有经济对应物：消费集带生存下限 $\underline x$ 时，$w$ 逼近生存成本处下半连续失效（预算集塌缩），$v$ 可能跳变；这就是我们坚持闭消费锥与 $p\gg0$ 的原因。

\subsubsection{对 \texorpdfstring{$(p,w)$}{(p,w)} 的拟凸性}

\begin{theorem}[拟凸性]
\label{thm:qvx}
对任意 $(p,w),(p',w')\in\R^L_{++}\times\R_+$ 与 $t\in[0,1]$，记 $\tilde p=tp+(1-t)p'$，$\tilde w=tw+(1-t)w'$，则
\[
v(\tilde p,\tilde w)\ \le\ \max\{v(p,w),\ v(p',w')\}.
\]
等价地，每个下水平集 $\{(p,w)\in\R^L_{++}\times\R_+:\ v(p,w)\le \bar u\}$ 都是凸集。
\end{theorem}
\begin{proof}
第一步（预算集的并集控制）。设 $x\in X$ 满足组合预算约束
\[
\tilde p\cdot x\ =\ t\,p\cdot x+(1-t)\,p'\cdot x\ \le\ tw+(1-t)w'\ =\ \tilde w .
\]
反设 $x\notin B(p,w)$ 且 $x\notin B(p',w')$，即 $p\cdot x>w$ 且 $p'\cdot x>w'$。以 $t$ 与 $1-t$ 加权相加：
\[
t\,p\cdot x+(1-t)\,p'\cdot x\ >\ t\,w+(1-t)w' ,
\]
与组合预算约束矛盾。故
\[
x\in B(p,w)\ \cup\ B(p',w').
\]

第二步（取上确界）。由第一步，$B(\tilde p,\tilde w)\subseteq B(p,w)\cup B(p',w')$，于是
\[
v(\tilde p,\tilde w)
=\sup_{x\in B(\tilde p,\tilde w)}u(x)
\le\sup_{x\in B(p,w)\cup B(p',w')}u(x)
=\max\{v(p,w),\ v(p',w')\}.
\]

第三步（与下水平集凸性的等价）。设 $v$ 在两点处均 $\le\bar u$，则由已证不等式，$v(\tilde p,\tilde w)\le\max\{\bar u,\bar u\}=\bar u$，故下水平集在凸组合下封闭，是凸集。反之设下水平集皆凸，取 $\bar u=\max\{v(p,w),v(p',w')\}$，两点都在水平集内，凸组合亦然，即得定理中的不等式。
\end{proof}

\begin{example}[拟凹性一般不成立]
\label{ex:notqc}
不少讲义把定理~\ref{thm:qvx} 误称为“$v$ 拟凹”。拟凹要求 $v(\tilde p,\tilde w)\ge\min\{v(p,w),v(p',w')\}$，这在一般情形\emph{不成立}。取完美替代偏好 $u(x)=x_1+x_2$，此时 $v(p,w)=w/\min\{p_1,p_2\}$。令 $\theta=((1,10),1)$、$\theta'=((10,1),1)$，则 $v(\theta)=v(\theta')=1$；而中点 $\tilde\theta=((5.5,5.5),1)$ 给出
\[
v(\tilde\theta)=\frac{1}{5.5}\approx0.18\ <\ 1=\min\{v(\theta),v(\theta')\},
\]
拟凹性被违反。注意该偏好是凸的（甚至线性），可见凸偏好也救不回拟凹性；反过来，定理~\ref{thm:qvx} 的证明完全没用凸偏好——拟凸性是\emph{预算约束的线性结构}的直接后果。
\end{example}

\notehead{经济解释：混合环境不创造机会}
定理~\ref{thm:qvx} 的第一步有独立的经济学内容：\emph{在两个价格—财富环境的凸组合下可行的消费计划，必然在其中一个原始环境下已经可行}。换言之，“混合/平均”价格—财富环境只会销毁消费机会，不会创造新机会——组合环境的可行集被控制在其两个来源可行集的\emph{并集}之内。因此平均环境的间接效用不超过两个环境中较好者。用下水平集的语言：使消费者福利不超过给定水平 $\bar u$ 的所有“坏环境”构成凸集——两个坏环境的任何凸组合仍是坏环境。政策含义：若两组（税后）价格—财富组合都使某家庭福利低于贫困线，则它们的任何凸组合也不高于贫困线；“福利上低于某阈值的环境”这一集合没有凹陷的漏洞。把该结论限制在固定 $w$ 上，就得到常见陈述“$v$ 对 $p$ 拟凸”。

\begin{figure}[H]
\centering
\begin{tikzpicture}[scale=1.0]
  \draw[axisline] (0,0) -- (6.4,0) node[below,axislabel] {$p_1$（$p_2$ 固定）};
  \draw[axisline] (0,0) -- (0,4.6) node[left,axislabel] {$w$};
  % shaded lower level set below the middle contour
  \fill[LightRed,fill opacity=0.6] plot[smooth] coordinates {(0,0) (0.25,0.7) (1,1.4) (2.25,2.1) (3.6,2.66) (5,3.13)} -- (5,0) -- cycle;
  % three contours w = c sqrt(p1), increasing level upward
  \draw[curveA] plot[smooth] coordinates {(0,0) (1,0.9) (2.25,1.35) (4,1.8) (5.5,2.11)};
  \draw[curveB] plot[smooth] coordinates {(0,0) (0.25,0.7) (1,1.4) (2.25,2.1) (3.6,2.66) (5,3.13)};
  \draw[curveC] plot[smooth] coordinates {(0,0) (0.25,1.0) (1,2.0) (2.25,3.0) (3.2,3.57)};
  % labels
  \node[figlabel] at (5.65,2.11) {$\bar v_1$};
  \node[figlabel] at (5.15,3.13) {$\bar v_2=\bar u$};
  \node[figlabel] at (3.35,3.57) {$\bar v_3$};
  \node[figlabel] at (2.5,0.85) {$\{v\le\bar u\}$（凸集）};
  \draw[arrblue] (4.6,3.9) -- (3.7,4.35);
  \node[figlabel] at (4.75,4.25) {$v$ 递增方向};
\end{tikzpicture}
\caption{价格—财富平面上的间接效用等值线与凸的下水平集}
\label{fig:vcontour}
\figexplan{固定 $p_2$，在 $(p_1,w)$ 平面上画三条间接效用等值线 $\bar v_1<\bar v_2<\bar v_3$。效用沿左上方增大：财富更高、价格更低则间接效用更高。阴影为下水平集 $\{v\le\bar u\}$（$\bar u=\bar v_2$），它是凸集——这正是拟凸性（定理~\ref{thm:qvx}）的几何内容。等值线向左上方移动对应福利改善；任两点的“坏区域”无凹陷漏洞。}
\end{figure}

\subsection{偏好性质与需求、间接效用性质的对照}

把前几节的结论放在一起，可以列一张“假设账单”：偏好的每一条性质，分别给需求对应 $x(p,w)$ 与间接效用 $v(p,w)$ 买到什么。表中“任意偏好”一行特别值得注意：$v$ 以 $\sup$ 良定、零次齐次、对 $p$ 非增、对 $(p,w)$ 拟凸，这四件事不花任何偏好假设——它们纯粹是预算结构（加上 $\sup$ 记号）的赠品。

\begin{table}[H]
\centering
\caption{偏好性质与 Walrasian 需求、间接效用的性质对照}
\begin{tabular}{p{3.1cm}p{5.4cm}p{5.6cm}}
\toprule
偏好/效用假设 & Walrasian 需求 $x(p,w)$ & 间接效用 $v(p,w)$ \\
\midrule
任意偏好 & 可能空 & 以 $\sup$ 良定；零次齐次、对 $p$ 非增、对 $(p,w)$ 拟凸 \\
$u$ 连续 & 非空、紧值、上半连续（Berge） & 连续；$\sup=\max$ \\
局部非饱和 & 预算绑定 $p\cdot x=w$ & 对 $w$ 严格递增 \\
严格单调 & 预算绑定；最优束处无松弛 & 对 $p_i$ 严格递减（若每个最优束 $x_i>0$） \\
凸偏好 & 凸值：$x(p,w)$ 是凸集 & 无新增（拟凸性本就不需要凸偏好） \\
严格凸偏好 & 单值且连续，成为需求函数 & 无新增，但为可微性与 Roy 恒等式铺路 \\
\bottomrule
\end{tabular}
\end{table}

对表中两行作简短证明与说明。\emph{凸偏好 $\Rightarrow$ 需求凸值}：设 $x,y\in X(p,w)$，则 $u(x)=u(y)=v(p,w)$；偏好凸意味着 $u(tx+(1-t)y)\ge\min\{u(x),u(y)\}=v(p,w)$，而 $p\cdot(tx+(1-t)y)\le tw+(1-t)w=w$，故凸组合可行且最优，$tx+(1-t)y\in X(p,w)$。注意这与 $v$ 的拟凸性是\emph{两件不同的事}：前者靠偏好凸性，后者靠预算线性结构。\emph{严格凸偏好 $\Rightarrow$ 单值}：若 $x\ne y$ 均最优，则 $tx+(1-t)y\succ x$ 且仍可行，与最优性矛盾；单值加半连续（推论~\ref{cor:cont}）得连续。最后，\emph{局部非饱和 $\Rightarrow$ 预算绑定}（Walras 法则）：若 $p\cdot x^*<w$，则 $x^*$ 的邻域内有更优且仍可负担的束（证明同定理~\ref{thm:incw} 第二步），矛盾。

\subsection{例题：Cobb--Douglas 间接效用与六条性质验证}

\begin{example}[Cobb--Douglas 的 $v(p,w)$]
\label{ex:cd}
设 $u(x)=x_1^{\alpha}x_2^{1-\alpha}$，$0<\alpha<1$，$p\gg0$，$w>0$。内点解满足
\[
\mathrm{MRS}=\frac{\alpha x_2}{(1-\alpha)x_1}=\frac{p_1}{p_2}
\quad\Longrightarrow\quad
p_1x_1=\frac{\alpha}{1-\alpha}\,p_2x_2 ,
\]
与预算绑定 $p_1x_1+p_2x_2=w$ 联立，得支出份额 $p_1x_1=\alpha w$、$p_2x_2=(1-\alpha)w$，即
\[
x_1(p,w)=\frac{\alpha w}{p_1},\qquad x_2(p,w)=\frac{(1-\alpha)w}{p_2}.
\]
代回效用函数：
\[
v(p,w)=\Bigl(\frac{\alpha w}{p_1}\Bigr)^{\alpha}\Bigl(\frac{(1-\alpha)w}{p_2}\Bigr)^{1-\alpha}
= A\,\frac{w}{p_1^{\alpha}p_2^{1-\alpha}},
\qquad A=\alpha^{\alpha}(1-\alpha)^{1-\alpha}.
\]
记 $f(p)=p_1^{\alpha}p_2^{1-\alpha}$（Cobb--Douglas 价格指数），则 $v(p,w)=A\cdot w/f(p)$：间接效用正比于\emph{实际财富} $w/f(p)$，且财富的边际效用 $A/f(p)$ 只依赖价格、不依赖财富——这是位似偏好的特征。
\end{example}

\textbf{逐条验证六条性质。}

（i）\emph{$\sup=\max$。}$u$ 连续、$B(p,w)$ 紧（命题~\ref{prop:max}），且显式解 $x(p,w)\gg0$ 达到 $v$：$v=u(x(p,w))$。

（ii）\emph{零次齐次。}对 $t>0$，
\[
v(tp,tw)=A\,\frac{tw}{(tp_1)^{\alpha}(tp_2)^{1-\alpha}}
=A\,\frac{tw}{t\,p_1^{\alpha}p_2^{1-\alpha}}=v(p,w),
\]
其中用到 $(tp_1)^{\alpha}(tp_2)^{1-\alpha}=t^{\alpha+1-\alpha}p_1^{\alpha}p_2^{1-\alpha}=t\,f(p)$。

（iii）\emph{对 $w$ 严格递增。}$\partial v/\partial w=A/f(p)=v/w>0$，与定理~\ref{thm:incw} 一致（Cobb--Douglas 偏好单调，自然局部非饱和）。

（iv）\emph{对 $p_i$ 非增（此处严格递减）。}
\[
\frac{\partial v}{\partial p_1}=-\frac{\alpha v}{p_1}<0,\qquad
\frac{\partial v}{\partial p_2}=-\frac{(1-\alpha)v}{p_2}<0 .
\]
严格性正与命题~\ref{prop:decpstrict} 印证：内点最优束两种商品均正消费。

（v）\emph{连续性。}$v$ 是初等函数，处处连续；这当然也可由推论~\ref{cor:cont}（Berge）直接得到。

（vi）\emph{拟凸性。}设 $v(p,w)=Ac$，$v(p',w')\le Ac$，即 $w\le c\,f(p)$、$w'\le c\,f(p')$，其中 $c=\max\{w/f(p),\,w'/f(p')\}$。由加权 AM--GM 不等式，加权几何平均 $f$ 是凹函数：
\[
f\bigl(tp+(1-t)p'\bigr)\ \ge\ t\,f(p)+(1-t)\,f(p') .
\]
于是
\[
\tilde w=tw+(1-t)w'\ \le\ c\bigl[t\,f(p)+(1-t)f(p')\bigr]\ \le\ c\,f(\tilde p),
\]
即 $v(\tilde p,\tilde w)=A\tilde w/f(\tilde p)\le Ac=\max\{v(p,w),v(p',w')\}$，定理~\ref{thm:qvx} 的不等式成立。几何上，下水平集 $\{v\le Ac\}=\{(p,w):w\le c\,f(p)\}$ 是凹函数 $f$ 的下图形（hypograph），当然是凸集——这正是图~\ref{fig:vcontour} 阴影区域的来源。

\section{Marshallian 需求求解与 Kuhn--Tucker 条件}
% ============================================================
% 第五讲第四分节：Marshallian 求解、存在性与 Kuhn--Tucker 条件
% ============================================================

\subsection{Marshallian 需求的求解与可微性准备}
\label{sec:kt-setup}

设消费集 $X=\R^n_+$，价格 $p\in\R^n_{++}$，财富 $w\ge 0$。Marshallian（非补偿）效用最大化问题为
\begin{equation}
 v(p,w)=\max_{x\in\R^n_+}\ u(x)\quad\st\quad p\cdot x\le w,
 \qquad
 x(p,w)\in\argmax_{x\in\R^n_+}\{u(x):p\cdot x\le w\}.
 \label{eq:ump-kt}
\end{equation}
存在性不需要任何微分工具：若 $u$ 连续且 $p\gg 0$，则 $B(p,w)$ 非空紧，最大值达到；偏好严格凸时最优束唯一，需求成为单值函数（存在性与唯一性的完整证明见 \S 5.2）。本节处理更精细的问题：\emph{给定解存在，如何把它解出来}。路线有两条：或猜测内点解，用等式约束的拉格朗日法求解后再回头验证猜测；或直接使用允许不等式与角点的 Kuhn--Tucker 条件。后者的适用范围覆盖前者，失效模式更少，是本节的重点。

在动用一阶条件之前，先澄清一个概念性的负债：可微性从何而来。

\begin{remark}[可微性是表示的性质，不是偏好的性质]
\label{rem:diff}
偏好关系 $\succeq$ 本身没有导数的概念。说“偏好可微”，严格含义是\emph{存在}一个可微的效用表示。三点澄清如下。
（一）可微性不是序数不变的。$u(x)=x_1x_2$ 可微；取严格递增但带扭结的 $\phi(t)=t+\tfrac12|t-1|$，则 $\phi\circ u$ 是同一偏好的合法表示，却在 $x_1x_2=1$ 处不可微——同一偏好，一个表示光滑，一个带扭结。反过来，某些偏好\emph{任何}表示都不可微：完全互补 $u=\min\{x_1,x_2\}$ 的无差异曲线有直角，MRS 从 $+\infty$ 跳到 $0$，这种跳跃无法被任何单调变换抹平。因此“存在可微表示”是对偏好的真实限制，而“手头这个表示恰好可微”只是运气。
（二）可微性不能被有限数据证伪。对任何有限组满足显示偏好一致性的价格—选择观测，既可以构造穿过观测点的光滑无差异曲线，得到可微的理性化偏好；也可以构造折线无差异曲线，得到带扭结的理性化偏好。于是可微性不是一条可检验的假设，而是为了书写一阶条件而支付的建模代价。
（三）这笔代价换来的内容是序数的。MRS 在严格递增的可微变换下不变（见下一分节的证明）；内点最优条件 $\mathrm{MRS}_{jk}=p_j/p_k$ 中，基数乘子 $\lambda$ 被完全消去。换言之，“哪个函数可微”依赖表示的选择，但“最优点的梯度平行于价格向量”这一几何事实不依赖。
\end{remark}

\subsection{拉格朗日函数的构造逻辑}
\label{sec:lagrange-build}

Kuhn--Tucker 方法的第一步，是把所有约束统一写成 $g(x)\ge 0$ 的形式。预算约束与非负约束改写为
\[
 g_0(x)=w-p\cdot x\ge 0,
 \qquad
 g_i(x)=x_i\ge 0\quad(i=1,\dots,n).
\]
为什么坚持这个方向？这纯粹是符号约定，但统一方向能让乘子的符号规则唯一化：对最大化问题，约束一律写 $g\ge 0$、拉格朗日函数中乘子一律取正号并入，则最优乘子一律非负。若有人坚持把预算写成 $p\cdot x-w\le 0$ 仍以正号并入，乘子的最优符号会翻转为非正；两种约定自洽时都正确，混乱只出现在混用之时。

按此约定，拉格朗日函数为
\begin{equation}
 \mathcal L(x,\lambda,\mu)
 =u(x)+\lambda\,(w-p\cdot x)+\sum_{i=1}^n\mu_i\,x_i,
 \qquad \lambda\in\R,\ \mu\in\R^n.
 \label{eq:lagrangian}
\end{equation}

\textbf{为什么预算约束只配一个乘子，而 $n$ 个非负约束各配一个乘子？}乘子度量“放宽某条约束的价值”，其个数由约束的条数决定，更准确地说，由\emph{可能分别绑定}的约束决定。预算只有一条，其影子价值是\emph{财富的边际效用}：把 $w$ 提高 $\mathrm dw$，可达效用大约增加 $\lambda\,\mathrm dw$。非负约束有 $n$ 条，且哪些商品落在角点因价格与偏好而异：Cobb--Douglas 偏好下全部商品内点消费；完全替代偏好下永远有商品落在角点；更一般的情形是部分商品零消费、部分商品内点，必须逐商品记录各自的影子价值。若强行共用一个乘子，就无法表达“商品 1 在内点、商品 2 被压在角点”的差异化结构。

$\mu_i$ 的经济含义是商品 $i$ 非负约束的影子补贴：假如允许 $x_i<0$（消费者可以按市价“倒卖”商品 $i$），最优值还能提高多少。由于我们不允许，$\mu_i$ 度量这条禁令的机会成本；它只在该约束真正起作用（$x_i=0$）时才可能为正——这正是后面互补松弛条件的内容。

还有一个观察值得记录：对乘子求偏导恰好恢复对应的约束，
\[
 \frac{\partial\mathcal L}{\partial\lambda}=w-p\cdot x\ge 0,
 \qquad
 \frac{\partial\mathcal L}{\partial\mu_i}=x_i\ge 0,
\]
对决策变量求偏导并令其为零则给出驻点条件。Kuhn--Tucker 系统的“四组条件”并非拼凑，而是“$\mathcal L$ 对全体变量的最优性条件 + 符号约束”的自然拆解，下节逐条展开。

\subsection{Kuhn--Tucker 四组条件逐条详解}
\label{sec:kt-conditions}

\begin{theorem}[Kuhn--Tucker 必要条件]
\label{thm:kt}
设 $u\in C^1$，$x^*$ 是问题\eqref{eq:ump-kt}的局部解。由于预算约束与非负约束都是仿射函数，约束规格自动成立。故存在乘子 $\lambda\in\R$、$\mu\in\R^n$，使得：
\begin{enumerate}
\item \textbf{原始可行性}：$p\cdot x^*\le w$，$x^*\ge 0$；
\item \textbf{驻点条件}：$u_i(x^*)=\lambda p_i-\mu_i$，$i=1,\dots,n$；
\item \textbf{对偶可行性}：$\lambda\ge 0$，$\mu_i\ge 0$，$i=1,\dots,n$；
\item \textbf{互补松弛}：$\lambda\,(w-p\cdot x^*)=0$，$\mu_i x_i^*=0$，$i=1,\dots,n$。
\end{enumerate}
\end{theorem}

四组条件缺一不可。下面逐条说明“为什么需要”，并为每条配一个删去它就会失效的反例。

\subsubsection{原始可行性：候选点必须先合法}

第一条看起来最平凡：候选点必须落在预算集内。但它的地位是逻辑上最先的：四组条件里只有这一条约束“点本身”，其余三条构造的是关于乘子的系统。删去它，一个不可行点可以顺利通过其余三条件。

\begin{example}[删去可行性的反例：买不起的极乐点]
\label{ex:no-feas}
取 $u(x)=8-(x_1-1)^2-(x_2-1)^2$（极乐点偏好，唯一最大点为 $(1,1)$，$u=8$），$p=(1,1)$，$w=1$。考察 $x=(1,1)$：花费 $2>w$，不可行。但在该点 $\nabla u=(0,0)$，取 $\lambda=\mu_1=\mu_2=0$：驻点条件 $0=0-0$ 成立；乘子非负成立；互补松弛 $\lambda(w-p\cdot x)=0$、$\mu_ix_i=0$ 也成立。也就是说，\emph{除可行性外全部条件都满足的点是全局无约束最大点}——消费者最想要的点，却买不起。该问题真正的解是 $(\tfrac12,\tfrac12)$，效用 $7.5$。没有第①条，KT 系统会把“最好但买不起”认证为解。
\end{example}

\subsubsection{驻点条件：沿可行方向不能再改进}

驻点条件 $u_i(x^*)=\lambda p_i-\mu_i$ 可以读作：边际效用 = 预算的影子价值 $\times$ 价格 $-$ 非负约束的影子补偿。它的几何内容是：在最优点，效用沿任何\emph{可行方向}的一阶变化不能为正。设 $x^*$ 处预算绑定，可行方向 $d$ 满足 $p\cdot d\le 0$，且若 $x_i^*=0$ 则 $d_i\ge 0$。用驻点条件与符号约束计算方向导数：
\[
 \nabla u(x^*)\cdot d=(\lambda p-\mu)\cdot d
 =\underbrace{\lambda\,p\cdot d}_{\le 0}
  \underbrace{-\sum_{i:\,x_i^*=0}\mu_i d_i}_{\le 0}\ \le 0,
\]
其中用到 $x_i^*>0$ 时 $\mu_i=0$（互补松弛）、$d_i\ge 0$ 时 $\mu_i d_i\ge 0$。反过来，若存在可行方向使方向导数为正，沿它走一小步即可改进，与最优性矛盾——这就是“为什么需要②”的核心。乘子 $(\lambda,-\mu)$ 恰好是把 $\nabla u(x^*)$ 分解到“起作用约束的法锥”上的那组权重。

\begin{example}[删去驻点条件的反例：可行但可改进的内点]
\label{ex:no-stat}
取 $u=x_1x_2$，$p=(1,1)$，$w=2$，候选点 $x=(1.5,\,0.5)$：可行且预算绑定，$u=0.75$。检验驻点条件：$u_1=x_2=0.5$，$u_2=x_1=1.5$，要求 $0.5=\lambda=1.5$，无 $\lambda$ 可满足——条件②正确地拒绝了它。事实上把一元钱从商品 2 移到商品 1（移向 $(1,1)$，$u=1>0.75$）即可严格改进。注意：删去②后，\emph{任何}可行点都能平凡地通过其余三组条件——取 $\lambda=\mu=0$，互补松弛全部成立。四组条件中只有②携带一阶信息。
\end{example}

\subsubsection{乘子非负：影子价格的方向必须正确}

乘子是约束的影子价格，非负性说的是\emph{方向}：放宽一条约束，最优值只能不降。对 $\lambda$：财富增加扩大预算集（新旧预算集嵌套），最大效用单调不减，故在可微点处 $\lambda=\partial v/\partial w\ge 0$。对 $\mu_i$：把 $x_i\ge 0$ 放宽为 $x_i\ge-\varepsilon$ 同样只会扩大可行集，故 $\mu_i\ge 0$。若解出负乘子，等于宣称“约束越紧、处境越好”，这在单调嵌套的可行集结构下是荒谬的。

\begin{example}[删去乘子非负的反例：伪角点]
\label{ex:no-sign}
取 $u=x_1x_2$，$p=(1,1)$，$w=2$，候选点 $x=(2,0)$：可行角点，$u=0$，而真正的最优是 $(1,1)$，$u=1$。检验：$u_1=x_2=0$，$u_2=x_1=2$。驻点条件要求 $0=\lambda-\mu_1$、$2=\lambda-\mu_2$；互补松弛中 $\mu_1x_1=2\mu_1=0$ 迫使 $\mu_1=0$，于是 $\lambda=0$，进而 $\mu_2=-2<0$——条件③被违反。若允许负乘子，$(\lambda,\mu_1,\mu_2)=(0,0,-2)$ 通过①②④，效用为零的伪角点被错误认证。对称的角点 $(0,2)$ 同理。正是③排除了这类“看似角点、实则劣于内点”的候选。
\end{example}

\subsubsection{互补松弛：约束不紧则乘子为零，乘子为正则约束必紧}

互补松弛把“约束是否起作用”与“乘子是否为零”绑在一起，两个方向都要：
（a）若 $w-p\cdot x^*>0$（预算没花完），则 $\lambda=0$——不稀缺的资源没有影子价格；
（b）若 $\lambda>0$，则 $p\cdot x^*=w$——财富有正的边际价值，预算必然用尽。
对每个商品同理：$\mu_i>0$ 意味着非负约束“真的在咬”（$x_i^*=0$）；$x_i^*>0$ 则该约束形同虚设，$\mu_i=0$。在单调偏好下 $\lambda>0$，于是预算恒绑定，这就是 Walras 定律的乘子版本。

\begin{example}[删去互补松弛的反例：不起作用的约束挂上了正乘子]
\label{ex:no-comp}
取 $u(x)=8-(x_1-1)^2-(x_2-1)^2$，$p=(1,1)$，$w=4$，候选点 $x=(3,1)$：可行、预算绑定、$u=4$；真正的最优是可负担的极乐点 $(1,1)$，$u=8$。在 $(3,1)$ 处 $\nabla u=(-4,0)$。驻点条件：$-4=\lambda-\mu_1$、$0=\lambda-\mu_2$。取 $(\lambda,\mu_1,\mu_2)=(0,4,0)$：①②③全部通过。但互补松弛要求 $\mu_1x_1=4\times 3=12=0$，不可能——商品 1 消费为 3，非负约束根本不起作用，$\mu_1$ 却为正，等于给一条虚设的约束标了正的影子价格。删去④，该点通过；保留④，它被正确拒绝。
\end{example}

\subsection{必要而非充分：反例、凸问题与拟凹效用的处理}
\label{sec:kt-suff}

\subsubsection{KT 点未必是最大点}

定理\ref{thm:kt}只保证“最优点满足 KT”，不保证“KT 点是最优点”。满足四组条件的候选可以是鞍点、甚至局部极小。

\begin{example}[KT 点是全局最小点]
\label{ex:saddle}
单商品问题：$\max\ x^3$ s.t. $0\le x\le w$（即 $p=1$）。考察 $x^*=0$：驻点条件 $3(x^*)^2=\lambda-\mu$ 即 $0=\lambda-\mu$；互补松弛 $\lambda(w-0)=0$ 迫使 $\lambda=0$（$w>0$），于是 $\mu=0$；乘子非负成立；可行性成立。四组条件在 $x^*=0$ 全部满足，但 $x^3$ 严格递增，$x^*=0$ 是可行集上的全局\emph{最小}点。多商品情形同理：非凸偏好可以在预算线上制造多个局部极大，每一个都满足 KT；仅凭一阶条件无法在其中挑选。要把“满足条件”升级为“确为最优”，需要全局结构——凹性与凸性。
\end{example}

\subsubsection{凸问题：KT 条件充分必要}

\begin{theorem}[凸问题的充分性]
\label{thm:convex-suff}
设 $u\in C^1$ 且凹，$p\gg 0$。若 $(x^*,\lambda,\mu)$ 满足定理\ref{thm:kt}的四组条件，则 $x^*$ 是问题\eqref{eq:ump-kt}的全局最大点；若 $u$ 严格凹，则最优唯一。
\end{theorem}
\begin{proof}
预算集由仿射不等式之交构成，自动是凸集。对任意可行 $x$（$p\cdot x\le w$，$x\ge0$），凹效用的一阶特征给出
\[
 u(x)-u(x^*)\le\nabla u(x^*)\cdot(x-x^*)=(\lambda p-\mu)\cdot(x-x^*),
\]
其中驻点条件代入 $\nabla u(x^*)=\lambda p-\mu$。展开右边并逐项估计：
\[
 (\lambda p-\mu)\cdot(x-x^*)
 =\underbrace{\lambda\,(p\cdot x-w)}_{\le 0}
  +\underbrace{\lambda\,(w-p\cdot x^*)}_{=0}
  -\underbrace{\mu\cdot x}_{\le 0}
  +\underbrace{\mu\cdot x^*}_{=0}.
\]
第一项非正，因为 $\lambda\ge0$ 且 $p\cdot x\le w$；第二、四项由互补松弛为零；第三项非正，因为 $\mu\ge0$、$x\ge0$。故 $u(x)\le u(x^*)$对一切可行 $x$ 成立，$x^*$ 是全局最大点。严格凹时，一阶特征中的等号仅当 $x=x^*$ 成立，最优唯一。
\end{proof}

\begin{remark}[鞍点视角]
\label{rem:saddle}
上述证明其实证明了更多：在给定乘子下，$x^*$ 在整个 $\R^n_+$ 上最大化 $\mathcal L(\cdot,\lambda^*,\mu^*)$（不必再检查预算），而 $(\lambda^*,\mu^*)$ 在给定 $x^*$ 下最小化 $\mathcal L(x^*,\cdot)$。这就是拉格朗日鞍点：
\[
 \mathcal L(x,\lambda^*,\mu^*)\ \le\ \mathcal L(x^*,\lambda^*,\mu^*)\ \le\ \mathcal L(x^*,\lambda,\mu),
 \qquad x\ge0,\ \lambda\ge0,\ \mu\ge0.
\]
“乘子是惩罚参数”的直觉由此严格化：最小化一方把违反预算的惩罚调到无利可图，最大化一方则在惩罚后的全景中选中 $x^*$。乘子非负（条件③）正是惩罚参数的合法取值域。
\end{remark}

\subsubsection{拟凹而非凹：Cobb--Douglas 型效用的处理}

凹性是充分性的便利条件，不是必要条件，而且许多常用效用并不凹。

\begin{remark}[拟凹非凹效用的两种处理]
\label{rem:quasi}
$u=x_1x_2$ 拟凹（上水平集 $\{x:x_1x_2\ge c\}$ 是凸区域），但沿对角射线 $u(t,t)=t^2$ 是凸函数，故 $u$ 不凹——直接套用定理\ref{thm:convex-suff}不合法。两种标准处理如下。
其一，\textbf{换一个凹表示}。取 $\ln u=\ln x_1+\ln x_2$：两个凹函数之和，凹；且与 $u$ 表示同一偏好。Marshallian 需求是 $\arg\max$，不随表示的单调变换改变，故用 $\ln u$ 求出的需求与用 $u$ 求出的完全相同；改变的只是 $\lambda$——乘子是基数的，随表示缩放，而 MRS 条件不变（下一分节证明）。凡是经单调变换可凹化的拟凹效用（整个 CES/CD 族均可）都可用这条路走通凸问题定理。
其二，\textbf{直接用拟凹版本的充分性}（Kuhn--Tucker--Arrow--Enthoven 定理）：$u$ 拟凹可微、且在候选点边际效用非负（例如强单调偏好）时，KT 条件仍充分。证明思路：拟凹性使上水平集 $\{x:u(x)\ge u(x^*)\}$ 是凸集，$x^*$ 位于其边界；在梯度非零处，边界有以 $\nabla u(x^*)$ 为法向的支撑超平面，于是 $u(x)>u(x^*)$ 蕴含 $\nabla u(x^*)\cdot(x-x^*)\ge0$；再复用定理\ref{thm:convex-suff}证明中“$(\lambda p-\mu)\cdot(x-x^*)\le0$”的链条即得矛盾。梯度为零的退化边界情形需要附加论证，此处从略。
\end{remark}

\subsection{与等式约束拉格朗日方法的对比}
\label{sec:compare}

经济学初级课程里的“拉格朗日法”默认两件事：预算约束写成等式 $p\cdot x=w$，所有商品内点消费 $x_i>0$。这两个先验假设一旦不真，等式法或无解、或解出负需求。表\ref{tab:compare}给出系统对比。

\begin{table}[H]
\centering
\caption{等式约束拉格朗日法与 Kuhn--Tucker 方法}
\label{tab:compare}
\begin{tabular}{@{}lp{5.1cm}p{5.9cm}@{}}
\toprule
维度 & 等式拉格朗日 & Kuhn--Tucker \\
\midrule
约束形式 & $p\cdot x=w$（先验绑定）；$x_i>0$ 先验假设 & $p\cdot x\le w$、$x_i\ge0$ 全部保留为不等式 \\
乘子个数与符号 & 一个 $\lambda$，无符号限制 & $\lambda\ge0$ 加每商品 $\mu_i\ge0$ \\
驻点条件 & $u_i=\lambda p_i$，全部商品 & $u_i=\lambda p_i-\mu_i$，逐商品 \\
互补松弛 & 无 & $\lambda(w-p\cdot x)=0$，$\mu_ix_i=0$ \\
角点处理 & 无：解出 $x_i<0$ 时只能人工截断 & 内生：$x_i=0$ 由 $\mu_i>0$ 解释 \\
预算是否绑定 & 强制绑定 & 由互补松弛判定（单调偏好下绑定） \\
典型失效 & 完全替代、$p_1\ne p_2$ 时方程组无解 & 不失效，自动给出角点 \\
\bottomrule
\end{tabular}
\end{table}

以完全替代 $u=x_1+x_2$ 说明最后一行。等式法的驻点条件是 $1=\lambda p_1$ 且 $1=\lambda p_2$，当 $p_1\ne p_2$ 时直接矛盾：方程组无解。初学者容易误读为“该消费者没有最优选择”；正确结论是\emph{没有内点最优}，最优在角点上，而等式法的框架里根本没有角点的位置。Kuhn--Tucker 把同样的计算写成 $1=\lambda p_1-\mu_1$、$1=\lambda p_2-\mu_2$：方程组不再矛盾，多余的约束压力由 $\mu$ 吸收，解自动指向角点（见例\ref{ex:corner-kt}）。

公平地说，等式法可以作为两步流程安全使用：先假设内点求解，再验证解确实满足 $x_i>0$。验证通过时两法等价；验证失败时必须回到 KT 补做角点分析。KT 的优势是把“求解—验证—修正”一步化，并且把“为什么不买商品 $i$”翻译成一个有经济含义的量 $\mu_i$，而不是一句“方程无解”。

\subsection{内点解与角点解}
\label{sec:interior-corner}

\subsubsection{内点解：切点条件}

\begin{proposition}[内点一阶条件]
\label{prop:interior-foc}
设 $u\in C^1$ 强单调（各 $u_i>0$），$p\gg0$，$x^*\gg0$ 是问题\eqref{eq:ump-kt}的解。则预算绑定 $p\cdot x^*=w$，且存在 $\lambda>0$ 使
\[
 u_i(x^*)=\lambda p_i,\qquad i=1,\dots,n,
\]
从而对任意商品对 $j,k$：
\[
 \mathrm{MRS}_{jk}(x^*)=\frac{u_j(x^*)}{u_k(x^*)}=\frac{p_j}{p_k},
 \qquad\text{等价地}\qquad
 \frac{u_j(x^*)}{p_j}=\frac{u_k(x^*)}{p_k}=\lambda .
\]
\end{proposition}
\begin{proof}
$x_i^*>0$ 与互补松弛 $\mu_ix_i^*=0$ 给出 $\mu_i=0$，驻点条件化为 $u_i(x^*)=\lambda p_i$。局部非饱和使预算绑定。强单调与 $p\gg0$ 保证 $\lambda=u_i/p_i>0$。对 $j,k$ 两式相除，$\lambda$ 消去，得 MRS 等于价格比。
\end{proof}

条件的两个读法都重要。\textbf{边际读法}：每一元花在任何商品上的边际效用相同，$u_i/p_i=\lambda$；若不等，把一元从比值低的商品移到比值高的商品即可严格改进。\textbf{交换读法}：$p_j/p_k$ 是市场允许的“以 $k$ 换 $j$”的比率，$\mathrm{MRS}_{jk}$ 是消费者自己的意愿比率；内点最优即二者打平——无差异曲线与预算线在 $x^*$ 相切，见图\ref{fig:interior}。注意相切的说法只在 $x^*$ 处梯度与价格向量平行这一意义上使用；图形上即 $\nabla u(x^*)$ 与 $p$ 同向。

\begin{figure}[H]
\centering
\begin{tikzpicture}[scale=1.0]
  \draw[axisline] (0,0)--(6.3,0) node[below,axislabel]{$x_1$};
  \draw[axisline] (0,0)--(0,5.3) node[left,axislabel]{$x_2$};
  % 预算线：(0,4) -- (5,0)，斜率 -0.8
  \draw[thick] (0,4)--(5,0);
  \node[figlabel,anchor=west] at (0.18,4.18) {$p\cdot x=w$};
  % 无差异曲线 x1*x2 = 5（经过 (2.5,2)）
  \draw[curveA] plot[domain=1.06:5.9,samples=60] (\x,{5/\x});
  \node[figlabel] at (5.62,1.22) {$\bar u$};
  % 较低无差异曲线 x1*x2 = 2
  \draw[curveB] plot[domain=0.72:5.7,samples=60] (\x,{2/\x});
  \node[figlabel] at (5.42,0.55) {$\underline u$};
  % 切点
  \node[nodewhite] at (2.5,2) {};
  \node[figlabel] at (2.72,1.78) {$x^*$};
  % 梯度箭头：方向 (2,2.5)
  \draw[arrblue] (2.5,2)--(3.31,3.01);
  \node[figlabel,anchor=south west] at (3.36,3.04) {$\nabla u(x^*)$};
  % 价格向量箭头：方向 (4,5)，从预算线上一点出发
  \draw[arrblue] (0.9,3.28)--(1.44,3.95);
  \node[figlabel,anchor=south] at (1.5,4.02) {$p$};
\end{tikzpicture}
\caption{内点最优：无差异曲线与预算线相切。}
\label{fig:interior}
\end{figure}
\figexplan{粗实线为无差异曲线 $x_1x_2=\bar u$，虚线为较低水平 $\underline u$；直线为预算线，斜率 $-p_1/p_2=-0.8$，切点为 $x^*$。梯度 $\nabla u(x^*)=(x_2^*,x_1^*)$ 垂直于无差异曲线，价格向量 $p$ 垂直于预算线；相切等价于两箭头平行，即 $\nabla u(x^*)=\lambda p$（本图取 $\lambda=0.5$，$p=(0.8,1)$）。虚线曲线在预算集内可达但严格劣。}

\subsubsection{角点解：完全替代的分情形求解}

\begin{example}[完全替代的角点解与 KT 验证]
\label{ex:corner-kt}
设 $u=x_1+x_2$，$p\gg0$，$w>0$。$u$ 线性故凹，定理\ref{thm:convex-suff}适用：任何满足 KT 的点即全局最优。$u_1=u_2=1$ 恒成立，分三种情形。
\textbf{情形一，$p_1<p_2$。}猜测 $x^*=(w/p_1,\,0)$。验证：$x_1^*>0$ 迫使 $\mu_1=0$，驻点条件给 $1=\lambda p_1$，即 $\lambda=1/p_1>0$；对商品 2，$1=\lambda p_2-\mu_2$ 给 $\mu_2=p_2/p_1-1>0$；互补松弛：$\lambda(w-p\cdot x^*)=\lambda\cdot0=0$，$\mu_1x_1=0$，$\mu_2\cdot0=0$，全部成立。故 $(w/p_1,0)$ 为全局最优。经济解读：$\mathrm{MRS}_{12}=1>p_1/p_2$，消费者对商品 1 的内部评价高于市场要价，全部预算买商品 1。
\textbf{情形二，$p_1>p_2$。}对称地，$x^*=(0,\,w/p_2)$，$\lambda=1/p_2$，$\mu_1=p_1/p_2-1>0$。
\textbf{情形三，$p_1=p_2$。}预算线上任一点 $p\cdot x=w$ 都是最优（$u=w/p_1$ 为常数），需求对应是整段预算线（多值）。KT 检验：$\lambda=1/p_1$、$\mu=0$ 在线上任一点成立；预算集内部的点（$p\cdot x<w$）则被互补松弛排除——$\lambda>0$ 与 $w-p\cdot x>0$ 不能并存，对应“钱没花完必可改进”。
三情形合并：比较 $\mathrm{MRS}_{12}=1$ 与 $p_1/p_2$，大于则全买商品 1，小于则全买商品 2，相等则整线最优。见图\ref{fig:corner}。
\end{example}

\begin{figure}[H]
\centering
\begin{tikzpicture}[scale=1.0]
  \draw[axisline] (0,0)--(6.2,0) node[below,axislabel]{$x_1$};
  \draw[axisline] (0,0)--(0,4.6) node[left,axislabel]{$x_2$};
  % 预算线：(0,2) -- (4,0)，对应 p1=1, p2=2, w=4
  \draw[thick] (0,2)--(4,0);
  \node[figlabel,anchor=west] at (0.14,2.17) {$p\cdot x=w$};
  \node[figlabel,anchor=east] at (-0.06,2.0) {$w/p_2$};
  \node[figlabel] at (4.0,-0.3) {$w/p_1$};
  % 无差异直线 x1+x2=4（穿过角点）
  \draw[curveA] (0.7,3.3)--(4,0);
  \node[figlabel] at (1.02,3.48) {$u=4$};
  % 较低无差异直线 x1+x2=2
  \draw[curveB] (0.3,1.7)--(2,0);
  \node[figlabel] at (0.62,1.9) {$u=2$};
  % 不可达无差异直线 x1+x2=5
  \draw[curveC] (1.3,3.7)--(5,0);
  \node[figlabel] at (1.62,3.88) {$u=5$};
  % 角点
  \node[nodeblue] at (4,0) {};
  \node[figlabel] at (3.62,0.42) {$x^*$};
\end{tikzpicture}
\caption{完全替代的角点最优。}
\label{fig:corner}
\end{figure}
\figexplan{参数取 $p_1=1<p_2=2$、$w=4$：预算线斜率 $-p_1/p_2=-0.5$，比无差异直线的斜率 $-1$ 平缓。最高可达的无差异直线 $u=4$ 与预算集只接触于 $x_1$ 轴上的截距 $x^*=(w/p_1,0)$；虚线 $u=2$ 可达但劣；点线 $u=5$ 完全在预算集之外。此时 $\mathrm{MRS}_{12}=1>p_1/p_2=0.5$，边际上消费者对商品 1 的评价高于市场，最优反应是全部预算购买商品 1。}

一般地，角点的判定法则从驻点条件读出：若 $x_j^*=0$，则
\[
 u_j(x^*)=\lambda p_j-\mu_j\le\lambda p_j,
\]
即商品 $j$ 的“每元边际效用”不超过已购商品的相应值——不买它不是被排除在最优之外，而是在边际上不值得买。内点条件是等式，角点条件是不等式，KT 框架把两者统一在同一个方程 $u_i=\lambda p_i-\mu_i$ 里。

\subsection{Cobb--Douglas 效用最大化：完整求解}
\label{sec:cd-solve}

\begin{example}[Cobb--Douglas 的完整求解、验证与乘子]
\label{ex:cd-full}
设 $u(x)=x_1^{\alpha}x_2^{1-\alpha}$，$0<\alpha<1$，$p\gg0$，$w>0$。按 KT 流程完整走一遍。

\textbf{第一步（内点判定）。}边界上 $u=0$；而 $w>0$ 时任何严格正的消费束（如 $(\varepsilon,\varepsilon)$ 适当小）可行且 $u>0$，故最优束 $x^*\gg0$，于是 $\mu_1=\mu_2=0$，驻点条件简化为 $u_i=\lambda p_i$。偏导数为
\[
 u_1=\alpha x_1^{\alpha-1}x_2^{1-\alpha}=\frac{\alpha u}{x_1},
 \qquad
 u_2=(1-\alpha)x_1^{\alpha}x_2^{-\alpha}=\frac{(1-\alpha)u}{x_2}.
\]

\textbf{第二步（解方程组）。}驻点条件两边同乘 $x_i$：
\[
 \frac{\alpha u}{\lambda}=p_1x_1,
 \qquad
 \frac{(1-\alpha)u}{\lambda}=p_2x_2.
\]
相加并用预算绑定 $p_1x_1+p_2x_2=w$（局部非饱和保证）：
\[
 w=\frac{u}{\lambda}
 \quad\Longrightarrow\quad
 \lambda=\frac{u(x^*)}{w},
 \qquad
 x_1^*=\frac{\alpha w}{p_1},
 \qquad
 x_2^*=\frac{(1-\alpha)w}{p_2}.
\]
支出份额为常数：$p_1x_1^*/w=\alpha$，$p_2x_2^*/w=1-\alpha$，与 $(p,w)$ 无关——这是 CD 偏好的指纹。

\textbf{第三步（Walras 定律验证）。}$p_1x_1^*+p_2x_2^*=\alpha w+(1-\alpha)w=w$。预算恰好用尽，与局部非饱和的预测一致。

\textbf{第四步（零次齐次验证）。}对 $t>0$，
\[
 x_1(tp,tw)=\frac{\alpha tw}{tp_1}=\frac{\alpha w}{p_1}=x_1(p,w),
 \qquad x_2(tp,tw)=x_2(p,w),
\]
需求只依赖实际预算集，无货币幻觉。

\textbf{第五步（充分性与唯一性）。}$u$ 的对数变换 $\ln u=\alpha\ln x_1+(1-\alpha)\ln x_2$ 是严格凹函数（Hessian 负定），表示同一偏好；由定理\ref{thm:convex-suff}，所求解是唯一全局最优。（$u$ 本身也是凹的——指数和为 1 的 CD 函数凹——两种论证皆可。）

\textbf{第六步（乘子与包络验证）。}代回得间接效用
\[
 v(p,w)=u(x^*)=\Big(\frac{\alpha w}{p_1}\Big)^{\alpha}
 \Big(\frac{(1-\alpha)w}{p_2}\Big)^{1-\alpha}
 =\alpha^{\alpha}(1-\alpha)^{1-\alpha}\,
 \frac{w}{p_1^{\alpha}p_2^{1-\alpha}} .
\]
一方面 $\partial v/\partial w=v/w$；另一方面 $\lambda=u(x^*)/w=v/w$，二者相等——这正是包络定理 $\lambda=\partial v/\partial w$ 的直接验证。若改用 $\ln u$ 求解，需求与 $v$ 的形状（经变换后）不变，但乘子变为 $\tilde\lambda=1/w\ne\lambda$：乘子是基数的，随表示改变，而需求与 MRS 条件不随表示改变（对照注\ref{rem:diff}与下一分节）。
\end{example}

\section{边际替代率与可分性}
% ============================================================
% 第五讲第五分节：边际替代率与可分性
% ============================================================

\subsection{边际替代率：定义与序数性质}
\label{sec:mrs-def}

\begin{definition}[边际替代率]
\label{def:mrs}
设 $u\in C^1$ 表示偏好 $\succeq$，且在 $x$ 处 $u_k(x)\ne 0$。商品 $j$ 对商品 $k$ 的边际替代率定义为
\[
 \mathrm{MRS}_{jk}(x)=\frac{u_j(x)}{u_k(x)}
 =\frac{\partial u/\partial x_j}{\partial u/\partial x_k}.
\]
几何上，过 $x$ 的无差异曲线的切线斜率为 $-u_j/u_k$，MRS 是其绝对值；经济上，它是在保持效用不变的前提下，为多得到一单位 $j$ 而愿意放弃的 $k$ 的数量的边际值。
\end{definition}

效用函数是序数的：同一偏好有无穷多个表示。一个自然的问题是，MRS 依赖表示的选择吗？答案是否定的。

\begin{proposition}[MRS 是序数概念]
\label{prop:mrs-ordinal}
设 $u\in C^1$，$\phi:\R\to\R$ 可微且严格递增（$\phi'>0$）。则 $\phi\circ u$ 与 $u$ 表示同一偏好，且在两者 MRS 同时有定义之处，
\[
 \mathrm{MRS}^{\phi\circ u}_{jk}(x)=\mathrm{MRS}^{u}_{jk}(x).
\]
\end{proposition}
\begin{proof}
链式法则给出
\[
 (\phi\circ u)_j(x)=\phi'\big(u(x)\big)\,u_j(x),
 \qquad
 (\phi\circ u)_k(x)=\phi'\big(u(x)\big)\,u_k(x).
\]
由于 $\phi'(u(x))>0$，分子分母的公共正因子消去：
\[
 \mathrm{MRS}^{\phi\circ u}_{jk}(x)
 =\frac{\phi'(u(x))\,u_j(x)}{\phi'(u(x))\,u_k(x)}
 =\frac{u_j(x)}{u_k(x)}
 =\mathrm{MRS}^{u}_{jk}(x). \qedhere
\]
同时 $u_k(x)\ne0$ 当且仅当 $(\phi\circ u)_k(x)\ne0$，故两者的定义域也一致。
\end{proof}

\begin{remark}[MRS 不含乘子，但一般含整个消费束]
\label{rem:mrs-lambda}
两个常被混淆的事实需要分开陈述。
其一，MRS 与乘子 $\lambda$ 无关，而 $\lambda$ 不是序数的。$\lambda=u_i/p_i$ 是\emph{基数}量：把 $u$ 换成 $\phi\circ u$，乘子变为 $\phi'(u)\,\lambda$，随表示改变；同一偏好下“财富的边际效用”没有表示无关的数值。而 MRS 在严格递增变换下不变（命题\ref{prop:mrs-ordinal}）。正因如此，最优性条件一旦写成 $\mathrm{MRS}_{jk}=p_j/p_k$，其中已不含任何依赖表示的对象；写成 $u_i=\lambda p_i$ 则还拖着随表示缩放的 $\lambda$。
其二，MRS 一般依赖消费束的\emph{全部}坐标。$u=x_1x_2$ 时 $\mathrm{MRS}_{12}=x_2/x_1$ 同时依赖两者；商品更多时，$\mathrm{MRS}_{jk}$ 还可能依赖第三种商品，例如 $u=x_1(x_2+x_3)$ 给出 $\mathrm{MRS}_{12}=(x_2+x_3)/x_1$，随 $x_3$ 变化。“组内 MRS 只依赖组内数量”是一种特殊的结构性质——可分性，第\ref{sec:separability}节专门讨论；在此之前，应当默认 MRS 是 $x_1,\dots,x_n$ 的函数。
\end{remark}

\subsection{内点最优：MRS 等于价格比}
\label{sec:mrs-price}

\begin{proposition}[从 Kuhn--Tucker 到 MRS 条件]
\label{prop:mrs-price}
设 $u\in C^1$ 强单调（各 $u_i>0$），$p\gg0$，$x^*\gg0$ 是效用最大化问题 $\max\{u(x):p\cdot x\le w,\ x\ge0\}$ 的解。则存在 $\lambda>0$ 使 $u_i(x^*)=\lambda p_i$ 对一切 $i$ 成立，从而对任意商品对 $j,k$：
\[
 \mathrm{MRS}_{jk}(x^*)=\frac{u_j(x^*)}{u_k(x^*)}=\frac{p_j}{p_k},
 \qquad\text{等价地}\qquad
 \frac{u_j(x^*)}{p_j}=\frac{u_k(x^*)}{p_k}=\lambda .
\]
\end{proposition}
\begin{proof}
由 $x_i^*>0$ 与互补松弛 $\mu_ix_i^*=0$ 得 $\mu_i=0$；驻点条件 $u_i=\lambda p_i-\mu_i$ 化为 $u_i=\lambda p_i$。强单调与 $p\gg0$ 保证 $\lambda=u_i/p_i>0$。对商品 $j,k$ 的两个方程相除，$\lambda$ 消去：
\[
 \frac{u_j(x^*)}{u_k(x^*)}=\frac{\lambda p_j}{\lambda p_k}=\frac{p_j}{p_k}. \qedhere
\]
\end{proof}

这一推导值得细看：乘子在相除中消去并非巧合，而是命题\ref{prop:mrs-ordinal}的体现——最优条件中唯一可能依赖表示的对象 $\lambda$ 必须也从任何表示无关的表述中消失，而 $\mathrm{MRS}_{jk}=p_j/p_k$ 恰是消去后的形式。条件本身有两个等价读法：\textbf{交换读法}，$p_j/p_k$ 是市场按 $k$ 计的 $j$ 的相对价格，$\mathrm{MRS}_{jk}(x^*)$ 是消费者在 $x^*$ 处的内部意愿比率，内点最优即两者打平，无差异曲线与预算线相切；\textbf{每元读法}，$u_j/p_j=\lambda$ 对所有 $j$ 成立，每一元花在任何商品上的边际效用相同——若商品 $j$ 的比值更高，把一元从别的商品挪到 $j$ 即可严格改进，与最优矛盾。

\begin{remark}[角点版本：等式弱化为不等式]
\label{rem:corner-mrs}
当 $x_j^*=0$ 而 $x_k^*>0$ 时，驻点条件是 $u_j(x^*)=\lambda p_j-\mu_j\le\lambda p_j$，于是
\[
 \frac{u_j(x^*)}{p_j}\ \le\ \lambda=\frac{u_k(x^*)}{p_k},
 \qquad\text{即}\qquad
 \mathrm{MRS}_{jk}(x^*)\ \le\ \frac{p_j}{p_k}
\]
（当 $u_j>0$ 时）。不买商品 $j$ 的原因由此有了精确表述：每一元花在 $j$ 上的边际效用不超过花在已购商品上的边际效用。完全替代 $u=x_1+x_2$、$p_1<p_2$ 即是 $1/p_1>1/p_2$（等价地 $\mathrm{MRS}_{12}=1>p_1/p_2$）的角点情形；等式与不等式分别对应内点与角点，这正是 Kuhn--Tucker 框架统一的两种状态。
\end{remark}

\subsection{可分偏好与组内边际替代率}
\label{sec:separability}

\begin{definition}[弱可分性]
\label{def:weak-sep}
设商品向量分块为 $(x,y)$，其中 $x\in\R^m_+$ 为关注组（$m\ge2$），$y$ 为其余商品。若存在函数 $v:\R^m_+\to\R$ 与 $U:\R\times Y\to\R$，使得偏好可表示为
\[
 u(x,y)=U\big(v(x),\,y\big),
\]
其中 $v$ 对组内向量严格递增、$U$ 对第一自变量严格递增，则称偏好对分组 $X|Y$ \textbf{弱可分}，$v$ 称为组子效用函数。
\end{definition}

含义是：组内商品只通过其聚合指标 $v(x)$ 进入整体效用——组内“怎么搭配”与组外商品无关，组内偏好自成体系。弱可分性的全部边际含义浓缩在下述命题中。

\begin{proposition}[组内 MRS 不依赖组外商品]
\label{prop:sep-mrs}
设 $u(x,y)=U(v(x),y)$，$u,U,v$ 均可微，$U_1>0$（$U$ 对第一自变量严格递增），$v_k(x)\ne0$。则对组内任意商品 $j,k$：
\[
 \mathrm{MRS}_{jk}(x,y)=\frac{u_j(x,y)}{u_k(x,y)}
 =\frac{v_j(x)}{v_k(x)},
\]
右端不含任何组外商品 $y$。
\end{proposition}
\begin{proof}
链式法则：
\[
 \frac{\partial u}{\partial x_j}
 =U_1\big(v(x),y\big)\,v_j(x),
 \qquad
 \frac{\partial u}{\partial x_k}
 =U_1\big(v(x),y\big)\,v_k(x),
\]
其中 $U_1$ 是 $U$ 对第一自变量的偏导数。两式相除，公共正因子 $U_1$ 消去，得 $\mathrm{MRS}_{jk}=v_j/v_k$，只依赖组内数量 $x$，与 $y$ 无关。
\end{proof}

证明的结构与命题\ref{prop:mrs-ordinal}完全同构：外层函数的边际权重（$\phi'$ 或 $U_1$）在比值中消去。这提示两个命题的逆命题也都有内容：在正则条件下（MRS 光滑且处处非零），“组内 MRS 与组外商品无关”反过来蕴含存在弱可分表示，证明需要沿无差异曲线积分的构造，此处从略。要点是：弱可分性是\emph{偏好}的性质而非某个表示的性质——命题\ref{prop:sep-mrs} 的结论本身用 MRS 表述，而 MRS 是序数的。

\begin{remark}[可加可分：特例与更强的含义]
\label{rem:additive}
\textbf{可加可分} $u(x)=\sum_{i}u_i(x_i)$ 是弱可分的极端情形：每组只有一件商品，组子效用就是单变量函数 $u_i$。此时
\[
 \mathrm{MRS}_{jk}=\frac{u_j'(x_j)}{u_k'(x_k)},
\]
只依赖 $x_j$ 与 $x_k$ 各自的数量，不依赖任何其他商品。相比一般弱可分（组内 MRS 不依赖组外商品），可加性额外要求组内 MRS 可按商品\emph{逐个}分离。
两点提醒。其一，“某个表示可加”不等于“偏好可加可分”：可加性在单调变换下一般不保持，但\emph{存在}可加表示是偏好的性质（有完整的刻画定理，条件较繁，从略）。经典例子：$u=x_1x_2$ 本身不可加，但 $\ln u=\ln x_1+\ln x_2$ 可加，故该偏好可加可分；一般地，二商品偏好可加可分当且仅当 MRS 可写成 $f(x_1)g(x_2)$ 的乘积形式。其二，可加性给出可检验的交叉效应限制：它要求第 $i$ 种商品的边际效用不随其他商品的数量变化（$u_i'(x_i)$ 只看 $x_i$），这在需求估计中表现为具体的系数约束。
\end{remark}

\begin{remark}[实证意义：两阶段预算]
\label{rem:two-stage}
弱可分性是需求系统估计中“分而治之”的许可证。理论逻辑：组内偏好只通过子效用 $v(x)$ 与组外商品互动，于是在最优处，组内需求只依赖组内相对价格与分配到该组的（间接）支出，组外价格只通过组支出总额发生作用。计量上这允许\textbf{两阶段预算}：第一阶段估计组间支出份额方程（食品、住房、交通等大类如何瓜分总支出），第二阶段在各组内部估计组内需求。AIDS、QUAIDS 以及嵌套 CES/CD 结构的实证需求系统大量依赖这一性质，其收益是参数维度的大幅下降。
可分性由此从“记号便利”变成“可证伪假设”：弱可分性蕴含（在适当的组价格指数化下）组内需求不直接依赖其他组的价格；若估计出的组内需求对组外价格仍有显著敏感反应，弱可分结构即被数据拒绝。理论结构、参数节约与假设检验在此合为一体。
\end{remark}

\subsection{例题：Cobb--Douglas、CES 与可分性检验}
\label{sec:mrs-examples}

\begin{example}[Cobb--Douglas 的 MRS 与可分性]
\label{ex:cd-mrs}
设 $u=x_1^{\alpha}x_2^{1-\alpha}$，$0<\alpha<1$。偏导数 $u_1=\alpha u/x_1$、$u_2=(1-\alpha)u/x_2$，故
\[
 \mathrm{MRS}_{12}=\frac{\alpha}{1-\alpha}\cdot\frac{x_2}{x_1}.
\]
\textbf{序数性检验：}换表示 $\ln u=\alpha\ln x_1+(1-\alpha)\ln x_2$，则
\[
 \frac{(\ln u)_1}{(\ln u)_2}=\frac{\alpha/x_1}{(1-\alpha)/x_2}
 =\frac{\alpha}{1-\alpha}\cdot\frac{x_2}{x_1},
\]
与原表示一致，验证命题\ref{prop:mrs-ordinal}。
\textbf{三商品与弱可分：}设 $u=x_1^{a}x_2^{b}x_3^{c}$（$a,b,c>0$），则
\[
 \mathrm{MRS}_{12}=\frac{a}{b}\cdot\frac{x_2}{x_1},
\]
不含 $x_3$，故对分组 $\{1,2\}|\{3\}$ 弱可分：取 $v(x_1,x_2)=x_1^{a/(a+b)}x_2^{b/(a+b)}$ 与 $U(v,x_3)=v^{a+b}x_3^{c}$ 即可，且 $v_1/v_2=(a/b)(x_2/x_1)$，与命题\ref{prop:sep-mrs}一致。又因 $\ln u=a\ln x_1+b\ln x_2+c\ln x_3$ 可加，该偏好还是可加可分的：CD 族在两种可分性下都表现良好。
\end{example}

\begin{example}[CES 的 MRS、替代弹性与可分性]
\label{ex:ces-mrs}
设 $u=(x_1^{\rho}+x_2^{\rho})^{1/\rho}$，$\rho\le1$、$\rho\ne0$（$\rho<1$ 保证拟凹）。偏导数
\[
 u_1=(x_1^{\rho}+x_2^{\rho})^{1/\rho-1}x_1^{\rho-1},
 \qquad
 u_2=(x_1^{\rho}+x_2^{\rho})^{1/\rho-1}x_2^{\rho-1},
\]
公共因子 $(x_1^{\rho}+x_2^{\rho})^{1/\rho-1}$ 在比值中消去：
\[
 \mathrm{MRS}_{12}=\Big(\frac{x_1}{x_2}\Big)^{\rho-1}.
\]
\textbf{序数性检验：}换表示 $\ln u=\tfrac1\rho\ln(x_1^{\rho}+x_2^{\rho})$，偏导数为 $(\ln u)_i=x_i^{\rho-1}/(x_1^{\rho}+x_2^{\rho})$，比值仍为 $(x_1/x_2)^{\rho-1}$，不变。
\textbf{替代弹性的读法：}MRS 对数量比 $(x_1/x_2)$ 的弹性为 $\rho-1$，替代弹性 $\sigma=1/(1-\rho)$：$\sigma$ 越大（$\rho\to1$），MRS 对数量比越敏感，无差异曲线越平直（极限是完全替代的直线）；$\rho\to0$ 时 MRS $\to x_2/x_1$，退回 CD（$\sigma=1$）；$\rho\to-\infty$ 时趋于完全互补的直角。
\textbf{最优条件与可分性：}内点最优要求 $(x_1/x_2)^{\rho-1}=p_1/p_2$，即
\[
 \frac{x_1}{x_2}=\Big(\frac{p_2}{p_1}\Big)^{\sigma},
\]
需求比只依赖价格比。又因单调变换 $x_1^{\rho}+x_2^{\rho}$ 可加，CES 偏好可加可分；CD 是其 $\rho\to0$ 极限，与例\ref{ex:cd-mrs}的结论相洽。
\end{example}

\begin{example}[一个不可分偏好]
\label{ex:nonsep}
设 $u=(x_1+x_2)(x_2+x_3)$。检验分组 $\{1,3\}|\{2\}$ 是否弱可分：
\[
 \frac{\partial u}{\partial x_1}=x_2+x_3,
 \qquad
 \frac{\partial u}{\partial x_3}=x_1+x_2,
 \qquad
 \mathrm{MRS}_{13}=\frac{x_2+x_3}{x_1+x_2},
\]
随组外商品 $x_2$ 变化，故该分组不弱可分。这个例子同时说明一个容易被忽视的平凡性：把\emph{单件}商品算作一组时，任何效用都能写成 $U(v(x_j),\text{其余})$（取 $v(x_j)=x_j$、$U=u$ 即可），弱可分性因此只在组内至少两件商品时才有内容；命题\ref{prop:sep-mrs}的价值正在于把非平凡的情形（组内 MRS 摆脱组外商品）与平凡情形区分开。
\end{example}

\section{包络定理与 Roy identity}
% ============================================================
% 第五讲 5.6：包络定理与 Roy identity
% ============================================================
\subsection{包络定理}

前几节把 Marshallian 需求 $x(p,w)$ 与间接效用 $v(p,w)$ 定义为最优化问题的解与最优值。本节回答一个更一般的问题：\textbf{当问题的参数变动时，最优值如何变动？}答案是包络定理（envelope theorem）：最优值对参数的导数，等于目标函数（或拉格朗日函数）直接对参数求偏导，而\emph{不需要}考虑最优解本身随参数移动带来的间接效应。这一定理是整个对偶理论（Roy identity、Shephard 引理、Slutsky 方程）的技术引擎。

\begin{theorem}[一般包络定理（无约束形式）]\label{thm:env-general}
设 $f:\R^{n}\times\Theta\to\R$ 连续可微，$\Theta\subseteq\R$ 为开区间，且对每个 $\theta\in\Theta$，最大化问题
\[
 v(\theta)=\max_{x\in D}\,f(x,\theta)
\]
在 $D$ 的内部有唯一解 $x^*(\theta)$，且 $x^*(\cdot)$ 可微。若 $v$ 在 $\theta$ 处可微，则
\[
 \boxed{\;v'(\theta)=\frac{\partial f}{\partial \theta}\bigl(x^*(\theta),\theta\bigr).\;}
\]
\end{theorem}

\begin{proof}
由定义，$v(\theta)=f\bigl(x^*(\theta),\theta\bigr)$ 对所有 $\theta$ 成立。两边对 $\theta$ 用链式法则求全导数：
\[
 v'(\theta)
 =\underbrace{\sum_{i=1}^{n}\frac{\partial f}{\partial x_i}\bigl(x^*(\theta),\theta\bigr)\,\frac{\mathrm dx_i^*}{\mathrm d\theta}(\theta)}_{\text{间接效应：解随参数移动}}
 \;+\;\underbrace{\frac{\partial f}{\partial \theta}\bigl(x^*(\theta),\theta\bigr)}_{\text{直接效应}}.
\]
关键在第一项。由于 $x^*(\theta)$ 是\emph{内部}最优解且 $f$ 可微，它必须满足一阶条件
\[
 \frac{\partial f}{\partial x_i}\bigl(x^*(\theta),\theta\bigr)=0,\qquad i=1,\ldots,n,
\]
否则沿某个坐标方向微调 $x_i$ 即可增大 $f$，与最优性矛盾。把一阶条件代入上式，间接效应整项为零，只剩下直接效应。
\end{proof}

\begin{remark}[为什么间接项会消失]
直觉如下：在最优解处，决策者已经把“通过调整 $x$ 获利”的空间用尽了——\textbf{最优性即一阶意义上无利可图}。参数 $\theta$ 变动一点点时，最优解 $x^*(\theta)$ 确实会移动，但这个移动是二阶小事：移动 $x$ 带来的收益在 $x^*$ 附近的边际值为零（一阶条件），所以调整量乘以零边际值，一阶贡献为零。用泰勒展开的语言说，$f\bigl(x^*(\theta+\mathrm d\theta),\theta+\mathrm d\theta\bigr)$ 与 $f\bigl(x^*(\theta),\theta+\mathrm d\theta\bigr)$ 之差是 $O(\mathrm d\theta^2)$，故对导数（一阶量）无贡献。几何上，值函数曲线在每一点都与某条“固定 $x$ 的”曲线相切，故称“包络”：$v$ 是一族曲线 $\{f(x,\cdot)\}_{x\in D}$ 的上包络线。
\end{remark}

\begin{remark}[带约束的形式]
经济学的最优化问题通常带约束，而参数往往恰好出现在约束里（价格、财富都进入预算约束而非效用函数）。设
\[
 v(\theta)=\max_{x}\;\{f(x,\theta):g(x,\theta)=0\},
\]
构造拉格朗日函数 $\mathcal L(x,\lambda,\theta)=f(x,\theta)+\lambda\,g(x,\theta)$，内点最优解 $(x^*,\lambda^*)$ 满足 $\mathcal L_x=0$。由于约束在最优解处绑定，$g(x^*(\theta),\theta)=0$，故
\[
 v(\theta)=f\bigl(x^*(\theta),\theta\bigr)=\mathcal L\bigl(x^*(\theta),\lambda^*(\theta),\theta\bigr).
\]
对 $\theta$ 全微分并利用 $\mathcal L_x=0$ 消去 $x^{*\prime}(\theta)$ 项，再利用 $g=0$ 使含 $\lambda^{*\prime}(\theta)$ 的项也消失，得
\[
 v'(\theta)=\frac{\partial \mathcal L}{\partial\theta}\bigl(x^*(\theta),\lambda^*(\theta),\theta\bigr).
\]
结论：\textbf{求值函数对参数的导数，只需对拉格朗日函数直接求偏导，然后在最优解处取值。}约束随参数移动的部分已被乘子 $\lambda$ 计价。
\end{remark}

\subsection{包络定理应用于效用最大化}

回到 UMP。设 $u$ 连续可微、偏好严格凸、局部非饱和，$p\gg0$、$w>0$，最优解 $x(p,w)$ 唯一且为内点，则 $v$ 可微。UMP 的拉格朗日函数为
\[
 \mathcal L(x,\lambda;p,w)=u(x)+\lambda\,(w-p\cdot x),
\]
其中预算约束写成 $w-p\cdot x=0$（局部非饱和保证绑定）。

\begin{proposition}[间接效用的包络性质]\label{prop:env-ump}
在上述假设下，
\[
 \frac{\partial v(p,w)}{\partial w}=\lambda(p,w),\qquad
 \frac{\partial v(p,w)}{\partial p_i}=-\lambda(p,w)\,x_i(p,w),\quad i=1,\ldots,n.
\]
\end{proposition}

\begin{proof}
对 $\mathcal L=u(x)+\lambda(w-p\cdot x)$ 分别关于 $w$ 与 $p_i$ 求偏导：
\[
 \frac{\partial\mathcal L}{\partial w}=\lambda,\qquad
 \frac{\partial\mathcal L}{\partial p_i}=-\lambda x_i.
\]
由带约束的包络定理，$v_w=\partial\mathcal L/\partial w=\lambda$，$v_{p_i}=\partial\mathcal L/\partial p_i$，并在最优解 $(x(p,w),\lambda(p,w))$ 处取值，即得两式。
\end{proof}

两条等式各有清晰的经济含义。

\begin{remark}[财富的影子价值：$\partial v/\partial w=\lambda$]
$\lambda$ 是预算约束的乘子，度量\textbf{约束放松一单位的边际价值}：财富增加一元，最优效用增加 $\lambda$ 个效用单位。因此 $\lambda$ 被称为\emph{财富的边际效用}或\emph{收入的影子价格}。它把“效用”这个不可观测的序数刻度与“元”这个基数刻度连起来：$\lambda$ 是效用对元的汇率。注意 $\lambda$ 本身依赖 $(p,w)$——越富的人一元钱的边际效用（按此效用表示）通常越低，这正是 $\lambda(p,w)$ 随 $w$ 递减的形态（在效用函数为凹的货币变换下）。
\end{remark}

\begin{remark}[涨价等于变穷：$\partial v/\partial p_i=-\lambda x_i$]
价格 $p_i$ 上升一元时，若消费者维持原消费束 $x(p,w)$ 不变，第 $i$ 种商品每单位多花一元，总支出增加 $x_i$ 元。也就是说，\textbf{$p_i$ 上涨一元在预算上的效果，等价于财富损失 $x_i$ 元}；而财富损失 $x_i$ 元的效用代价是 $\lambda x_i$，故 $v_{p_i}=-\lambda x_i$。注意没有消费的商品（$x_i=0$）涨价不直接损害消费者——这解释了为什么间接效用只通过\emph{被购买的}商品的价格起作用。把两式相除，$x_i$ 恰好是被“价格损失/财富边际价值”约分还原出来，这正是 Roy identity 的来源。
\end{remark}

\subsection{Roy identity}

\begin{theorem}[Roy identity]\label{thm:roy}
设 $u$ 连续可微、偏好严格凸且局部非饱和，$x(p,w)$ 唯一、内点，$v$ 在 $(p,w)$ 处可微且 $v_w(p,w)\neq0$。则对每种商品 $i$，
\[
 \boxed{\;x_i(p,w)=-\frac{\partial v(p,w)/\partial p_i}{\partial v(p,w)/\partial w}.\;}
\]
\end{theorem}

\begin{proof}
由命题~\ref{prop:env-ump}，
\[
 -\frac{v_{p_i}(p,w)}{v_w(p,w)}
 =-\frac{-\lambda(p,w)\,x_i(p,w)}{\lambda(p,w)}=x_i(p,w).\qedhere
\]
\end{proof}

Roy identity 的价值在于\textbf{方向的可逆性}：给定任何一个行为良好的间接效用函数 $\tilde v(p,w)$，不必重新解 UMP，直接求偏导即可恢复 Marshallian 需求。这把“解最优化问题”变成“求两个偏导数之比”，也是宏观与计量中由间接效用设定（如 LOGEX、Almost Ideal 体系的对偶写法）直接导出需求系统的标准手法。

\begin{remark}[用包络定理重证预算耗尽]
包络性质还能反过来推出 Walras 法则。$v$ 对 $(p,w)$ 零次齐次：$v(tp,tw)=v(p,w)$ 对一切 $t>0$ 成立。在 $t=1$ 处对 $t$ 求导（Euler 定理）：
\[
 \sum_{i=1}^{n}v_{p_i}(p,w)\,p_i+v_w(p,w)\,w=0 .
\]
代入命题~\ref{prop:env-ump} 的 $v_{p_i}=-\lambda x_i$、$v_w=\lambda$，约去 $\lambda>0$：
\[
 \sum_{i=1}^{n}p_i\,x_i(p,w)=w .
\]
局部非饱和下的预算耗尽——第五章前半用“邻域内存在更好束”证明的结论——在此作为齐次性与包络定理的代数推论重新出现。同一条思路也给出 $\lambda$ 的齐次性质：对 $v(tp,tw)=v(p,w)$ 分别关于 $tp_i$ 与 $tw$ 求偏导可得 $\lambda(tp,tw)=t^{-1}\lambda(p,w)$，即财富的边际效用随货币单位的放大而同比缩小——货币只是标签。
\end{remark}

\begin{remark}[包络定理的普适性：一个定理，三处引理]
包络定理不专属消费者理论。把“目标函数对参数求偏导”套到不同的问题上，就得到不同领域的经典引理：
\begin{itemize}
\item UMP（参数 $p,w$ 进入约束）$\Rightarrow$ $v_w=\lambda$、$v_{p_i}=-\lambda x_i$，进而 Roy identity；
\item EMP（参数 $p$ 进入目标）$\Rightarrow$ Shephard 引理 $h_i=\partial e/\partial p_i$（下节）；
\item 厂商利润最大化 $\pi(p)=\max_x\{p\cdot y(x)-c(x)\}$（参数 $p$ 进入目标）$\Rightarrow$ Hotelling 引理：供给 $y_i(p)=\partial\pi/\partial p_i$；对成本最小化值函数 $c(w,y)$ 关于产出 $y$ 求导 $\Rightarrow$ 边际成本 $\partial c/\partial y$。
\end{itemize}
识别一个引理是不是包络定理的化身，只需看一句话：\textbf{参数出现在哪里，导数就从哪里来；最优解的移动不贡献一阶项。}
\end{remark}

\begin{example}[Cobb--Douglas 的数值验证]\label{ex:roy-cd}
取 $u(x_1,x_2)=x_1^{1/2}x_2^{1/2}$，价格 $p=(4,1)$，财富 $w=8$。第五章前半已得
\[
 v(p,w)=\alpha^{\alpha}(1-\alpha)^{1-\alpha}\frac{w}{p_1^{\alpha}p_2^{1-\alpha}}
 =\frac{1}{2}\cdot\frac{8}{4^{1/2}\cdot 1^{1/2}}=2 .
\]
\textbf{第一步：两个包络导数。}$v_w=v/w=2/8=1/4$，即财富的边际效用为 $1/4$；$v_{p_1}=-\alpha v/p_1=-\tfrac12\cdot 2/4=-1/4$，即第一种商品涨价一元使效用损失 $1/4$。
\textbf{第二步：Roy identity。}
\[
 x_1=-\frac{v_{p_1}}{v_w}=-\frac{-1/4}{1/4}=1,
\]
与直接求解 UMP 得到的 $x_1=\alpha w/p_1=\tfrac12\cdot 8/4=1$ 完全一致。同理 $x_2=-v_{p_2}/v_w=-( -\tfrac12\cdot 2/1)/(1/4)=4=\tfrac12\cdot 8/1$。验证了 $\lambda=v_w=1/4$，且 $v_{p_1}=-\lambda x_1=-\tfrac14\cdot 1$、$v_{p_2}=-\lambda x_2=-\tfrac14\cdot 4$，包络关系逐项吻合。
\end{example}

\begin{remark}[财富边际效用的递减形态]
接例~\ref{ex:roy-cd}：对 Cobb--Douglas，$\lambda=v_w=K/(p_1^{\alpha}p_2^{1-\alpha})$（$K=\alpha^{\alpha}(1-\alpha)^{1-\alpha}$）\emph{不依赖} $w$——同位偏好下财富边际效用恒定。这不是常态：换成效用 $u(x)=\ln(x_1^{\alpha}x_2^{1-\alpha})$（同一偏好、另一表示），则 $v=\alpha\ln\alpha+(1-\alpha)\ln(1-\alpha)+\alpha\ln w-\alpha\ln p_1-\cdots$，$\lambda=\alpha/w$ 随 $w$ 严格递减。两件事同时说明：\textbf{$\lambda$ 的具体形态依赖效用表示的选取}（序数性），但“给定表示后 $\lambda$ 度量约束的边际价值”这一定性含义（以及 Roy identity 的比值）在表示替换下不变。
\end{remark}

\begin{remark}[一个常见的误用]
Roy identity 的分母 $v_w=\lambda$ 不能为零：局部非饱和保证 $\lambda>0$（多一元财富总能严格改善），故在标准假设下除法合法。但若把不可微的 $v$（如完全替代偏好的角点解附近）硬套进公式，就会得到错误结果——包络定理要求内点、唯一与可微，这些条件在角点解处失效，此时应回到 Kuhn--Tucker 条件与对应的互补松弛不等式。
\end{remark}

\section{支出最小化与对偶：EMP、Shephard 引理与对偶恒等式}
% ============================================================
% 第五讲 5.7：支出最小化、Shephard 引理与对偶恒等式
% ============================================================
\subsection{为什么需要第二个最优化问题}

UMP 把 $(p,w)$ 作为给定，最大化 $u(x)$。细看会发现一个结构上的不便：\textbf{参数 $p$ 与 $w$ 全都进入约束集 $B(p,w)$，而不进入目标函数}。参数一动，整个可行域就变了形状，值函数 $v$ 与需求 $x$ 的性质（凹凸、单调、齐次）因此难以直接刻画：例如 $v$ 对 $(p,w)$ 是拟凹而非凹，对 $p$ 的二阶性质更无从谈起。

经济学家的对策是把问题“翻转”一次：\textbf{固定效用水平，最小化支出}。这样一来，价格 $p$ 进入目标函数，而约束（达到目标效用）不再依赖 $p$——目标函数对参数的全部二阶性质立刻可用。这就是支出最小化问题（expenditure minimization problem, EMP）。

先看它为什么必要，再定义它。

\begin{example}[苹果涨价：两种效应方向可能相反]\label{ex:apple}
苹果价格上升。一方面，\emph{相对价格}变了：苹果相对变贵，消费者会\textbf{用其他商品替代苹果}——替代效应使苹果消费减少。另一方面，\emph{实际财富}变了：同样的货币财富买不到原来的消费束，消费者实际变穷——若苹果是正常品（越富吃得越多），财富效应同样使苹果消费减少；但若苹果是劣等品（越富吃得越少，如某种廉价主食），财富效应反而使苹果消费\textbf{增加}。两种效应方向相反时，总效应的符号成为经验问题。UMP 把两者搅在一起；要把它们干净地拆开，需要一个“固定效用、让支出变动”的问题——替代效应正是在那里定义的。本节先把 EMP 的代数结构建好，Slutsky 分解留到下节。
\end{example}

\subsection{支出最小化问题及其与成本最小化的同构}

\begin{definition}[支出最小化问题]\label{def:emp}
给定价格 $p\gg0$ 与目标效用水平 $\bar u$，
\[
 \min_{x\in\R^n_+}\;p\cdot x\qquad\st\;u(x)\ge\bar u.
\]
其最优值称为\textbf{支出函数}
\[
 e(p,\bar u)=\min\{p\cdot x:x\ge0,\;u(x)\ge\bar u\},
\]
其最优解（严格凸偏好下唯一）称为\textbf{Hicksian（补偿）需求}
\[
 h(p,\bar u)=\arg\min\{p\cdot x:x\ge0,\;u(x)\ge\bar u\}.
\]
\end{definition}

\begin{remark}[解的存在性]
设 $u$ 连续、$\bar u$ 可达（存在 $x$ 使 $u(x)\ge\bar u$）、$p\gg0$。可行集 $\{x\ge0:u(x)\ge\bar u\}$ 非空且为闭集（$u$ 连续）。任取可行点 $\bar x$，则最优解若存在必落在子集 $\{x\ge0:u(x)\ge\bar u,\;p\cdot x\le p\cdot\bar x\}$ 中；该子集有界（$p\gg0$ 与 $x\ge0$ 给出 $0\le x_i\le(p\cdot\bar x)/p_i$）且闭，故为紧集，连续函数 $p\cdot x$ 在其上达到最小值。因此 EMP 的解存在；偏好严格凸时唯一（与 UMP 中同款的凸组合论证）。
\end{remark}

\begin{remark}[与厂商成本最小化的同构]\label{rem:isomorphism}
EMP 与厂商理论中的成本最小化问题在数学上\textbf{完全同构}，只需把消费者理论的对象按下表“翻译”成厂商理论的对象：
\begin{center}
\begin{tabular}{ll}
\toprule
消费者理论（EMP） & 厂商理论（成本最小化）\\
\midrule
消费束 $x\in\R^n_+$ & 投入向量 $x\in\R^n_+$\\
价格 $p\gg0$ & 投入价格 $w\gg0$\\
效用 $u(x)$（视作“享乐生产函数”） & 生产函数 $f(x)$\\
目标效用 $\bar u$ & 目标产量 $y$\\
Hicksian 需求 $h(p,\bar u)$ & 条件要素需求 $x(w,y)$\\
支出函数 $e(p,\bar u)$ & 成本函数 $c(w,y)$\\
\bottomrule
\end{tabular}
\end{center}
把 $u$ 读作“用消费束 $x$ 生产出享乐水平 $u(x)$ 的生产函数”，消费者就是一个把效用当产出的厂商：他以市场价格 $p$ 购买“投入品”（商品），以最省钱的方式“生产”出 $\bar u$ 单位的效用。因此，厂商理论中关于成本函数的全部定理（齐次、凹性、Shephard 引理）可以原封不动地搬到 $e$ 上；反过来，本节对 $e$ 的每条证明也都是成本函数相应性质的证明。这种一鱼两吃正是对偶方法的价值。
\end{remark}

\begin{remark}[EMP 的一阶条件：与 UMP 相同的切点]
内点最优解 $h(p,\bar u)$ 满足（拉格朗日函数 $\mathcal L=p\cdot x+\mu(\bar u-u(x))$ 对 $x$ 的一阶条件）
\[
 p_i=\mu\,u_i\bigl(h(p,\bar u)\bigr),\qquad i=1,\ldots,n,
\]
即 $\mathrm{MRS}_{ij}=u_i/u_j=p_i/p_j$：\textbf{等支出线与无差异曲线相切}。这与 UMP 的一阶条件（$u_i=\lambda p_i$，同样 $\mathrm{MRS}_{ij}=p_i/p_j$）形式完全相同——两个问题在各自的对偶点 $(p,w)\leftrightarrow(p,\bar u=v(p,w))$ 上是\emph{同一个}切点。这正是对偶恒等式 $h(p,v(p,w))=x(p,w)$（定理~\ref{thm:duality}）的几何内容：UMP 的切点在预算线上，EMP 的切点在同一条无差异曲线上，而两个约束在同一束上同时绑定。区别只在外围：UMP 沿一族预算线找最高切点，EMP 沿一条无差异曲线找最平的等支出线。
\end{remark}

\subsection{支出函数与 Hicksian 需求的性质}

\begin{proposition}[支出函数的基本性质]\label{prop:e-props}
设 $u$ 连续、局部非饱和，$p\gg0$，$\bar u$ 可达。则：
\begin{enumerate}
\item $e(p,\bar u)$ 对 $p$ \textbf{一次齐次}：$e(tp,\bar u)=t\,e(p,\bar u)$，$t>0$；
\item $e(p,\bar u)$ 对 $p$ \textbf{凹}；
\item $e(p,\bar u)$ 对 $\bar u$ \textbf{非减}；若偏好严格单调且 $\bar u$ 在可达范围内变动，则严格递增；
\item $h(p,\bar u)$ 对 $p$ \textbf{零次齐次}：$h(tp,\bar u)=h(p,\bar u)$，$t>0$。
\end{enumerate}
\end{proposition}

\begin{proof}
(1) 价格同比例放大不改变可行集 $\{x\ge0:u(x)\ge\bar u\}$，只把目标函数乘以 $t$：
\[
 e(tp,\bar u)=\min_{u(x)\ge\bar u}(tp)\cdot x=t\min_{u(x)\ge\bar u}p\cdot x=t\,e(p,\bar u).
\]
(2) 任取 $p^1,p^2\gg0$，$t\in(0,1)$，记 $p=tp^1+(1-t)p^2$，$x^*=h(p,\bar u)$。关键观察：$x^*$ 对两个子问题都可行（可行集只看 $u(x)\ge\bar u$，与价格无关），故
\[
 e(p,\bar u)=p\cdot x^*=t\,(p^1\cdot x^*)+(1-t)(p^2\cdot x^*)
 \;\ge\;t\,e(p^1,\bar u)+(1-t)\,e(p^2,\bar u),
\]
其中不等号用了“每个子问题的最小值不超过任一可行点的成本”。这正是凹性。
(3) 若 $\bar u'\ge\bar u$，则 $\{x:u(x)\ge\bar u'\}\subseteq\{x:u(x)\ge\bar u\}$；在更小的集合上求最小只会更大：$e(p,\bar u')\ge e(p,\bar u)$。严格递增的论证：若 $u$ 严格递增且 $e(p,\bar u')=e(p,\bar u)$，则公共最小点 $x$ 满足 $u(x)\ge\bar u'$，再沿任一坐标加 $\varepsilon e_i$ 可使效用更高而成本 $e+\varepsilon p_i$，与 $e(p,\bar u')$ 已是达到 $\bar u'$ 的最小成本矛盾（当 $\varepsilon$ 足够小时 $e+\varepsilon p_i<e(p,\bar u'')$ 对介于其间的 $\bar u''$ 成立）。
(4) 由 (1) 的证明，$\min_{u\ge\bar u}(tp)\cdot x=t\min_{u\ge\bar u}p\cdot x$，两边的最优解集合相同：把成本乘以正常数不改变哪个可行点成本最小。故 $h(tp,\bar u)=h(p,\bar u)$。
\end{proof}

\begin{remark}[对照 UMP 一侧]
这些性质与 $v$ 一侧形成镜像：$v$ 对 $(p,w)$ 零次齐次（规模不变），$e$ 对 $p$ 一次齐次（货币单位）；$v$ 对 $(p,w)$ 拟凹，$e$ 对 $p$ 凹。凹性是 $e$ 最有用的性质：凹函数在区间中点的函数值不低于端点值的平均——直觉是“混合价格下的最小支出，不超过分别面对两种价格的最小支出的平均”，因为混合价格下的最优束 $h(p,\bar u)$ 对两种原始价格都可行，可作为两者的备选方案。这个直觉将在下节直接生成 Slutsky 矩阵的负半定性。
\end{remark}

图~\ref{fig:emp-tangency} 给出 EMP 的几何：一族斜率为 $-p_1/p_2$ 的等支出线（iso-expenditure lines）从原点向外平移，最先触到无差异曲线 $u=\bar u$ 的那条决定了最小支出，切点即 Hicksian 需求。

\begin{figure}[H]
\centering
\begin{tikzpicture}[scale=.92]
  \draw[axisline] (0,0)--(6.1,0) node[below]{$x_1$};
  \draw[axisline] (0,0)--(0,4.7) node[left]{$x_2$};
  % 无差异曲线 xy=4
  \draw[curveA,domain=0.95:5.55,smooth] plot (\x,{4/\x});
  \node[figlabel] at (4.85,1.15) {$u=\bar u$};
  % 等支出线族（斜率 -1）：低、切、高
  \draw[curveC] (0,3.2)--(3.2,0);
  \draw[curveB] (0,4)--(4,0);
  \draw[curveC] (0,4.65)--(4.65,0);
  \node[figlabel] at (3.05,3.55) {$E_3$};
  \node[figlabel] at (3.95,4.25) {$E_2$};
  \node[figlabel] at (2.25,2.45) {$E_1$};
  \node[nodewhite] at (2,2) {};
  \node[figlabel,above right] at (2.05,2.05) {$x^*=h(p,\bar u)$};
\end{tikzpicture}
\caption{EMP：等支出线与无差异曲线的切点}
\figexplan{实线为无差异曲线 $u=\bar u$（其右上方向为可行区域 $u\ge\bar u$）。三条平行线段为等支出线 $p\cdot x=E_k$，斜率 $-p_1/p_2$，$E_1<E_2<E_3$。$E_1$（点线）与可行区域不相交，支出太低达不到 $\bar u$；$E_3$（点线）整段落入可行区域，支出浪费；$E_2$（虚线）恰好与无差异曲线相切于 $x^*$，是最小支出 $e(p,\bar u)=E_2$，切点即 Hicksian 需求 $h(p,\bar u)$。}
\label{fig:emp-tangency}
\end{figure}

\subsection{Shephard 引理}

\begin{theorem}[Shephard 引理]\label{thm:shephard}
设 $u$ 连续可微、偏好严格凸且局部非饱和，$\bar u>u(0)$ 可达，$p\gg0$，$h(p,\bar u)$ 为内点。若 $e$ 在 $(p,\bar u)$ 处对 $p$ 可微，则
\[
 \boxed{\;\frac{\partial e(p,\bar u)}{\partial p_i}=h_i(p,\bar u),\qquad i=1,\ldots,n.\;}
\]
\end{theorem}

\begin{proof}
EMP 的拉格朗日函数（约束写成 $\bar u-u(x)\le0$）为
\[
 \mathcal L(x,\mu;p,\bar u)=p\cdot x+\mu\bigl(\bar u-u(x)\bigr).
\]
\textbf{第一步：效用约束绑定。}若 $u\bigl(h(p,\bar u)\bigr)>\bar u$，由 $u$ 连续与 $\bar u>u(0)$，存在 $t\in(0,1)$ 使 $u\bigl(t\,h(p,\bar u)\bigr)=\bar u$，而其成本 $t\,p\cdot h<p\cdot h$，与 $h$ 的最小性矛盾。故 $u(h)=\bar u$，从而在最优解处
\[
 e(p,\bar u)=p\cdot h(p,\bar u)=\mathcal L\bigl(h(p,\bar u),\mu(p,\bar u);p,\bar u\bigr).
\]
\textbf{第二步：用包络定理。}把 $e$ 看作 EMP 的值函数，参数为 $p_i$。约束 $\bar u-u(x)\ge0$ 不含 $p$，故带约束的包络定理给出
\[
 \frac{\partial e}{\partial p_i}
 =\frac{\partial\mathcal L}{\partial p_i}\Bigl|_{(h,\mu)}
 =h_i(p,\bar u),
\]
其中 $\partial\mathcal L/\partial p_i=x_i$ 在 $x=h$ 处取值。与 Roy identity 同理：价格只出现在目标函数里，直接求偏导即可；解随价格移动的间接效应被一阶条件 $\mathcal L_x=0$ 消去。
\end{proof}

\begin{remark}[乘子的对照与“效用价格”]
同一个包络论证还给出 $\partial e/\partial\bar u=\mu$：$\mu$ 是\textbf{目标效用的影子价格}——多要求一个单位的效用要多花多少元。对偶地，UMP 的乘子 $\lambda=v_w$ 是\textbf{一元财富的边际效用}。比较两个问题的一阶条件（内点情形）
\[
 \text{UMP：}\;u_i=\lambda p_i;\qquad
 \text{EMP：}\;\mu u_i=p_i\;\text{即}\;u_i=\tfrac{1}{\mu}p_i,
\]
可见在对偶点上有 $\lambda=1/\mu$：一个把效用换成钱，一个把钱换成效用，两个乘子互为倒数。这也解释了 $e$ 与 $v$ 的反函数关系（见下节对偶恒等式）。
\end{remark}

\begin{example}[直接解 CD 的 EMP 并验证 Shephard]
取 $u=x_1^{\alpha}x_2^{1-\alpha}$。EMP 的一阶条件给出 $x_2=\frac{1-\alpha}{\alpha}\frac{p_1}{p_2}x_1$；代入约束 $x_1^{\alpha}x_2^{1-\alpha}=\bar u$ 解出
\[
 h_1=\bar u\Bigl(\frac{\alpha p_2}{(1-\alpha)p_1}\Bigr)^{1-\alpha},\qquad
 h_2=\bar u\Bigl(\frac{(1-\alpha)p_1}{\alpha p_2}\Bigr)^{\alpha},
\]
于是
\[
 e=p_1h_1+p_2h_2=\frac{\bar u\,p_1^{\alpha}p_2^{1-\alpha}}{\alpha^{\alpha}(1-\alpha)^{1-\alpha}} .
\]
Shephard 验证：$\partial e/\partial p_1=\alpha e/p_1$。另一方面
\[
 \alpha\frac{e}{p_1}
 =\frac{\alpha\,\bar u\,p_1^{\alpha-1}p_2^{1-\alpha}}{K}
 =\bar u\Bigl(\frac{\alpha p_2}{(1-\alpha)p_1}\Bigr)^{1-\alpha}=h_1,\qquad K=\alpha^{\alpha}(1-\alpha)^{1-\alpha},
\]
其中第二个等号是初等代数整理（把 $\alpha^{\alpha}(1-\alpha)^{1-\alpha}$ 拆进幂次）。同理 $\partial e/\partial p_2=(1-\alpha)e/p_2=h_2$。数值例：$\alpha=1/2$、$p=(4,1)$、$\bar u=2$，则 $e=2\cdot 2\cdot 1/0.5=8$，$h_1=\tfrac12\cdot 8/4=1$、$h_2=\tfrac12\cdot 8/1=4$恰为 \S 5.6 中 Roy identity 数值例（例~\ref{ex:roy-cd}）中 $(p,w)=(4,1)$、$w=8$ 的 Marshallian 需求，印证 $h(p,v(p,w))=x(p,w)$。
\end{example}

图~\ref{fig:hicksian-slide} 展示 Hicksian 需求如何随价格变化：$p_1$ 下降时，等支出线变平，切点沿\textbf{同一条}无差异曲线 $u=\bar u$ 向右下滑动——效用被钉死，只有相对价格在起作用。这正是“补偿”的含义，也是下节替代效应的定义。

\begin{figure}[H]
\centering
\begin{tikzpicture}[scale=.92]
  \draw[axisline] (0,0)--(6.1,0) node[below]{$x_1$};
  \draw[axisline] (0,0)--(0,4.7) node[left]{$x_2$};
  % 同一条无差异曲线 xy=4
  \draw[curveA,domain=0.95:5.55,smooth] plot (\x,{4/\x});
  \node[figlabel] at (4.85,1.15) {$u=\bar u$};
  % 两条等支出线
  \draw[curveB] (0,4)--(4,0);
  \node[figlabel] at (3.9,4.3) {$E$};
  \draw[curveC] (0,2.83)--(5.66,0);
  \node[figlabel] at (5.35,3.0) {$E'$};
  \node[nodewhite] at (2,2) {};
  \node[figlabel,above left] at (1.9,2.1) {$h$};
  \node[nodewhite] at (2.83,1.41) {};
  \node[figlabel,below right] at (2.88,1.32) {$h'$};
  \draw[arrblue] (2.35,1.95) to[bend left=28] (2.62,1.55);
\end{tikzpicture}
\caption{Hicksian 需求沿同一条无差异曲线滑动}
\figexplan{固定目标效用 $\bar u$（实线）。虚线为价格比 $p_1/p_2=1$ 时的等支出线，切点 $h$；点线为 $p_1$ 下降（$p_1/p_2=1/2$）后变平的等支出线，切点 $h'$。切点从 $h$ 滑到 $h'$：商品 1 变便宜，补偿需求 $h_1$ 上升、$h_2$ 下降，且最小支出由 $E$ 降到 $E'$（价格下降使达到同一效用更便宜）。整条轨迹始终贴着同一条无差异曲线——这就是“保持效用不变”的替代效应。}
\label{fig:hicksian-slide}
\end{figure}

\subsection{对偶恒等式}

UMP 与 EMP 是同一枚硬币的两面：一个固定财富求最高效用，一个固定效用求最小财富。下列四条恒等式把两个问题焊接成闭环。以下均设 $u$ 连续、偏好严格凸且局部非饱和（保证 $x$、$h$ 单值），$p\gg0$；涉及 $\bar u$ 的恒等式要求 $\bar u$ 落在\textbf{相关区间}内：$u(0)<\bar u$ 且 $\bar u$ 可达。

\begin{theorem}[对偶恒等式]\label{thm:duality}
在上述假设下：
\begin{enumerate}
\item $e\bigl(p,\,v(p,w)\bigr)=w$；
\item $h\bigl(p,\,v(p,w)\bigr)=x(p,w)$；
\item $v\bigl(p,\,e(p,\bar u)\bigr)=\bar u$；
\item $x\bigl(p,\,e(p,\bar u)\bigr)=h(p,\bar u)$。
\end{enumerate}
\end{theorem}

\begin{proof}
\textbf{(1)} 记 $\bar v=v(p,w)$，$x^*=x(p,w)$。一方面，$u(x^*)=v(p,w)=\bar v\ge\bar v$，即 $x^*$ 是 EMP（目标 $\bar v$）的可行点，故
\[
 e(p,\bar v)\le p\cdot x^*=w,
\]
末一等号用了局部非饱和下的预算耗尽（Walras 法则）。另一方面，设 $z$ 为 EMP（目标 $\bar v$）的任一可行点：$u(z)\ge\bar v=v(p,w)$。断言 $p\cdot z\ge w$：若不然 $p\cdot z<w$，则 $z$ 落在预算集内部，由局部非饱和，存在 $z'$ 靠近 $z$ 使 $p\cdot z'\le w$ 且 $u(z')>u(z)\ge v(p,w)$，与 $v$ 是 $B(p,w)$ 上的最大效用矛盾。对所有可行 $z$ 取最小，得 $e(p,\bar v)\ge w$。两方向合并即 $e(p,v(p,w))=w$。

\textbf{(2)} 由 (1) 的证明，$x^*=x(p,w)$ 可行且成本恰为最小值 $e(p,\bar v)=w$，故 $x^*$ 本身就是 EMP 的最优解：$h(p,v(p,w))=x(p,w)$。

\textbf{(3)} 记 $w'=e(p,\bar u)$，$h^*=h(p,\bar u)$。\emph{先证 $v(p,w')\ge\bar u$}：$h^*$ 满足 $p\cdot h^*=w'$（EMP 约束绑定可如 Shephard 引理证明中那样由 $\bar u>u(0)$ 与连续性得到，或直接由 $p\cdot h^*=e=w'$），故 $h^*\in B(p,w')$，且 $u(h^*)\ge\bar u$，于是 $v(p,w')\ge u(h^*)\ge\bar u$。\emph{再证 $v(p,w')\le\bar u$}：反设 $v(p,w')>\bar u$，取 $x^*\in B(p,w')$ 使 $u(x^*)=v(p,w')>\bar u$。由 $u$ 连续，对充分接近 $1$ 的 $t<1$ 有 $u(t x^*)>\bar u$；而
\[
 p\cdot(t x^*)=t\,(p\cdot x^*)\le t\,w'<w'=e(p,\bar u),
\]
（$t<1$ 且 $p\cdot x^*>0$，后者因 $u(x^*)>\bar u>u(0)$ 蕴含 $x^*\ne0$、$p\gg0$）。于是 $t x^*$ 以严格低于 $e(p,\bar u)$ 的成本达到了 $\bar u$，与 $e$ 的最小性矛盾。故 $v(p,e(p,\bar u))=\bar u$。

\textbf{(4)} 在恒等式 (2) 中取财富为 $w'=e(p,\bar u)$，得
\[
 h\bigl(p,\,v(p,w')\bigr)=x(p,w').
\]
再由 (3)，$v(p,w')=\bar u$，代入即 $h(p,\bar u)=x\bigl(p,e(p,\bar u)\bigr)$。
\end{proof}

四条恒等式两两配对，构成一个闭环：

\begin{center}
\begin{tikzpicture}[>=Stealth,node distance=1.05cm]
  \node[boxsoft] (ump) {UMP：固定 $w$\\ $v(p,w)=\max\{u(x):p\cdot x\le w\}$};
  \node[boxsoft,right=2.2cm of ump] (emp) {EMP：固定 $\bar u$\\ $e(p,\bar u)=\min\{p\cdot x:u(x)\ge\bar u\}$};
  \draw[arrblue] (ump)--node[figlabel,above]{$\bar u=v(p,w)$} (emp);
  \draw[arrblue] (emp)--node[figlabel,above]{$w=e(p,\bar u)$} (ump);
  \node[figlabel,below=0.35cm of ump] {$x(p,w)\ \longleftrightarrow\ h(p,\bar u)$};
\end{tikzpicture}
\end{center}
\figexplan{上箭头：把 UMP 的最优效用 $\bar u=v(p,w)$ 喂给 EMP，EMP 还回同样的财富 $w$ 且 $h=x$（恒等式 1、2）。下箭头：把 EMP 的最小支出 $w=e(p,\bar u)$ 喂给 UMP，UMP 还回同样的效用 $\bar u$ 且 $x=h$（恒等式 3、4）。两个问题互为镜像，$(p,w)\mapsto(p,\bar u)$ 的换算是无损的。}

\begin{remark}[“相关区间”不可省略]
恒等式 (3)、(4) 在 $\bar u\le u(0)$ 时会失效：什么都不买就已达 $\bar u$，$e(p,\bar u)=0$，而 $v(p,0)\ge u(0)$ 可能严格超过 $\bar u$。此时“最小支出”是平凡的零，无法携带关于偏好的信息。因此对偶理论习惯上只在 $\bar u\in(u(0),\sup u)$ 的区间内展开；该区间正是每一点效用都“值得花钱去购买”的范围。
\end{remark}

\begin{remark}[对偶性的方法论意义]
四条恒等式意味着：知道 $(v,x)$ 就能复原 $(e,h)$，反之亦然；而复原的运算是求逆与代入，不需要重新解任何最优化问题。再配合 Roy identity（由 $v$ 得 $x$）与 Shephard 引理（由 $e$ 得 $h$），四个对象 $v,x,e,h$ 中\textbf{任何一个}都足以生成其余三个。这为需求分析提供了自由度：研究者可以挑一个性质最好刻画的对象（通常是 $e$ 或 $h$，因为它们享受凹性与对称性）作为建模起点，其余对象按表索取——这正是下一节 Slutsky 矩阵的路线。
\end{remark}

把四个对象的性质并排放在一起，便于总览与记忆：

\begin{center}
\begin{tabular}{lllll}
\toprule
对象 & 定义 & 齐次性 & 凸凹/单调 & 生成工具\\
\midrule
$v(p,w)$ & UMP 最优值 & $(p,w)$ 零次 & 对 $w$ 增、对 $p$ 减、拟凹 & $\to x$：Roy identity\\
$x(p,w)$ & Marshallian 需求 & $(p,w)$ 零次 & 满足预算耗尽 & $\to v$：代入 $u$\\
$e(p,\bar u)$ & EMP 最优值 & $p$ 一次 & 对 $p$ 凹、对 $\bar u$ 增 & $\to h$：Shephard 引理\\
$h(p,\bar u)$ & Hicksian 需求 & $p$ 零次 & $D_ph$ 对称、负半定 & $\to e$：代入 $p\cdot h$\\
\bottomrule
\end{tabular}
\end{center}
\figexplan{四对象速查表：每行第一格是对象，末格是“由它生成另一个对象”的引理；中间两格列出各自的结构性质。配合对偶恒等式（$(p,w)\leftrightarrow(p,\bar u)$ 的无损换算），从任一格出发都能走完整张表。注意凹性集中在 $e$ 一侧、拟凹性在 $v$ 一侧——这正是“参数进入目标函数”的 EMP 在结构上更易刻画的原因。}

\section{Hicksian 需求与 Slutsky 分解}
% ============================================================
% 第五讲 5.8：Slutsky 矩阵与 Slutsky 方程
% ============================================================
\subsection{Slutsky 矩阵}

对偶理论最大的红利是关于\textbf{价格效应}的一组定理。回忆 Hicksian 需求 $h(p,\bar u)$ 把效用钉死、只让价格起作用，因此它对价格的导数恰好是“纯粹的”替代效应。

\begin{definition}[Slutsky 矩阵]
设 $h(p,\bar u)$ 可微。\textbf{Slutsky 矩阵}（补偿价格效应矩阵）定义为
\[
 D_p h(p,\bar u)=\Bigl[\frac{\partial h_i}{\partial p_j}(p,\bar u)\Bigr]_{i,j=1}^{n},
\]
即第 $(i,j)$ 元为：商品 $j$ 价格上升一单位时，为保持效用 $\bar u$ 不变，商品 $i$ 的补偿需求变化多少。在财富水平 $w$ 处取 $\bar u=v(p,w)$ 时，记
\[
 S(p,w)=D_p h\bigl(p,v(p,w)\bigr).
\]
\end{definition}

\subsection{Slutsky 矩阵的三大性质}

\begin{theorem}[对称性、负半定性与零次齐次性]\label{thm:slutsky-props}
设 $u$ 连续可微、偏好严格凸且局部非饱和，$e$ 对 $p$ 二阶连续可微，$h$ 可微且为内点。则：
\begin{enumerate}
\item \textbf{对称性}：$\dfrac{\partial h_i}{\partial p_j}=\dfrac{\partial h_j}{\partial p_i}$（$S$ 为对称矩阵）；
\item \textbf{负半定性}：对任意向量 $\mathrm dp\in\R^n$，$\mathrm dp^{\top}S\,\mathrm dp\le0$；
\item \textbf{行和为零}：$D_p h(p,\bar u)\cdot p=0$，即 $\sum_j\frac{\partial h_i}{\partial p_j}p_j=0$ 对每个 $i$ 成立。
\end{enumerate}
\end{theorem}

\begin{proof}
(1) Shephard 引理给出 $h(p,\bar u)=\nabla_p e(p,\bar u)$（按分量 $h_i=\partial e/\partial p_i$）。于是
\[
 D_p h(p,\bar u)=\nabla^2_p e(p,\bar u),
\]
即 Slutsky 矩阵是支出函数关于价格的 \textbf{Hessian 矩阵}。$e$ 对 $p$ 二阶连续可微，由 Young 定理（混合偏导相等，亦称 Schwarz/Clairaut 定理），
\[
 \frac{\partial h_i}{\partial p_j}
 =\frac{\partial^2 e}{\partial p_j\,\partial p_i}
 =\frac{\partial^2 e}{\partial p_i\,\partial p_j}
 =\frac{\partial h_j}{\partial p_i}.
\]
(2) 命题~\ref{prop:e-props}（上节）证明过 $e$ 对 $p$ 凹。二阶连续可微的凹函数其 Hessian 处处负半定：对任意 $\mathrm dp$，
\[
 e(p+t\,\mathrm dp)\le e(p)+t\,\nabla_p e(p)\cdot\mathrm dp
 \quad\text{与}\quad
 e(p+t\,\mathrm dp)=e(p)+t\,\nabla_p e(p)\cdot\mathrm dp+\tfrac{t^2}{2}\,\mathrm dp^{\top}\nabla^2_p e\,\mathrm dp+o(t^2)
\]
联立并除以 $t^2/2$、令 $t\to0$，得 $\mathrm dp^{\top}\nabla^2_p e\,\mathrm dp\le0$。代入 $\nabla_p e=h$ 即 $\mathrm dp^{\top}D_p h\,\mathrm dp\le0$。
(3) $h$ 对 $p$ 零次齐次（命题~\ref{prop:e-props} 第 4 条）。对齐次可微函数用 Euler 定理：若 $h_i(tp,\bar u)=h_i(p,\bar u)$，两边在 $t=1$ 处对 $t$ 求导，得
\[
 \sum_{j=1}^{n}\frac{\partial h_i}{\partial p_j}(p,\bar u)\,p_j=0,
\]
即 $(D_p h\,p)_i=0$ 对每个 $i$ 成立，故 $D_p h\cdot p=0$。
\end{proof}

\begin{corollary}[补偿自身价格效应非正]\label{cor:diag}
负半定矩阵的对角元非正：取 $\mathrm dp=e_i$（第 $i$ 个单位向量），得
\[
 \frac{\partial h_i}{\partial p_i}(p,\bar u)\le0.
\]
\end{corollary}
\begin{proof}
$\mathrm dp^{\top}S\,\mathrm dp\le0$ 中取 $\mathrm dp=e_i$ 立得 $S_{ii}\le0$。
\end{proof}

这就是\textbf{补偿需求定律}：保持效用不变，商品自身涨价（其余价格不动），对该商品的补偿需求不增。注意它与（未补偿的）需求定律不同：后者可以因财富效应而失效（Giffen 品），前者是偏好凸性的代数必然，不可违反——只要偏好是凸的，补偿自身价格效应必然非正。负半定性还蕴含更一般的命题：对任意价格变化向量 $\mathrm dp$，补偿需求的变化 $\mathrm dx=D_p h\,\mathrm dp$ 与价格变化本身的方向“不顺行”：
\[
 \mathrm dp\cdot\mathrm dx=\mathrm dp^{\top}D_p h\,\mathrm dp\le0.
\]

\subsection{Slutsky 方程：总效应的分解}

现在把两个世界接起来。对偶恒等式 $h(p,\bar u)=x\bigl(p,e(p,\bar u)\bigr)$ 说明：补偿需求就是“财富被恰好补偿到能维持效用 $\bar u$”的普通需求。对第 $i$ 个分量
\[
 h_i(p,\bar u)=x_i\bigl(p,\,e(p,\bar u)\bigr)
\]
两边关于 $p_j$ 求偏导（$\bar u$ 固定），用链式法则展开右边——$p_j$ 既直接进入 $x_i$，又通过最小支出 $e(p,\bar u)$ 间接进入：
\[
 \frac{\partial h_i}{\partial p_j}
 =\frac{\partial x_i}{\partial p_j}
 +\frac{\partial x_i}{\partial w}\bigl(p,e(p,\bar u)\bigr)\cdot\frac{\partial e}{\partial p_j}(p,\bar u).
\]
再用 Shephard 引理 $\partial e/\partial p_j=h_j$，并在对偶点 $\bar u=v(p,w)$（此时 $e(p,\bar u)=w$、$h_j=x_j$）处取值，移项即得：

\begin{theorem}[Slutsky 方程]\label{thm:slutsky-eq}
在定理~\ref{thm:slutsky-props} 的假设下，对一切 $i,j$：
\begin{equation}
 \boxed{\;
 \frac{\partial x_i}{\partial p_j}(p,w)
 =\underbrace{\frac{\partial h_i}{\partial p_j}(p,\bar u)}_{\text{替代效应}}
 -\underbrace{x_j(p,w)\,\frac{\partial x_i}{\partial w}(p,w)}_{\text{财富效应}},\qquad \bar u=v(p,w).\;}
 \label{eq:slutsky}
\end{equation}
矩阵形式（$\nabla_p x$ 为 $x$ 对 $p$ 的 Jacobian，$\nabla_w x=(\partial x_i/\partial w)_i$ 为列向量）：
\[
 \nabla_p x(p,w)=S(p,w)-\bigl(\nabla_w x(p,w)\bigr)\,x(p,w)^{\top}.
\]
\end{theorem}

\begin{proof}
如上推导；形式地：对 $h(p,\bar u)=x(p,e(p,\bar u))$ 关于 $p$ 全微分，
\[
 D_p h=D_p x+\bigl(\nabla_w x\bigr)\bigl(\nabla_p e\bigr)^{\top}
 =D_p x+\bigl(\nabla_w x\bigr)h^{\top},
\]
第二等号用了 Shephard 引理 $\nabla_p e=h$。移项并在 $\bar u=v(p,w)$ 处取值（此时 $h=x$）即得矩阵形式。
\end{proof}

\begin{remark}[逐项经济解释]
把 $p_j$ 上升一元对 $x_i$ 的影响拆成两块：
\begin{itemize}
\item \textbf{替代效应 $\partial h_i/\partial p_j$}：假想消费者被“补偿”到恰好维持原效用 $\bar u$（财富被调整为 $e(p,\bar u)$），此时 $x_i$ 的变化纯由相对价格驱动。它满足对称与负半定——理论对它的符号有硬约束（自身价格情形必为非正，见推论~\ref{cor:diag}）。
\item \textbf{财富效应 $-x_j\,\partial x_i/\partial w$}：撤销假想的补偿，回到固定财富 $w$。$p_j$ 上升一元使维持原消费束贵了 $x_j$ 元，等价于财富缩水 $x_j$ 元；财富缩水对 $x_i$ 的影响是 $-x_j\cdot\partial x_i/\partial w$。它的大小取决于 $x_j$（商品 $j$ 在预算中的份额越大，缩水越狠）与 $x_i$ 对财富的敏感度（$x_i$ 是正常品还是劣等品）。
\end{itemize}
“总效应 $=$ 替代效应 $+$ 财富效应”由此不再是图形直觉，而是恒等式——它对任何偏好（凸、连续、局部非饱和）无条件成立。
\end{remark}

\begin{remark}[Hicks 补偿与 Slutsky 补偿]
“补偿”有两种操作定义，须加区分。\textbf{Hicks 补偿}（本节的 $h$）：把财富调整到恰好维持\emph{原效用水平}，即 $w'=e(p',\bar u)$——精确但不可观察（$\bar u$ 不可观测）。\textbf{Slutsky 补偿}：把财富调整到恰好还能买得起\emph{原消费束}，即 $w''=p'\cdot x(p,w)$——可观察（只用到价格与数量），Slutsky（1915）原始论文用的正是这种。两种补偿一阶等价：对 $p_j$ 的微分，Slutsky 补偿给出的分解项为 $\partial x_i/\partial p_j+x_j\,\partial x_i/\partial w$，与式~(\ref{eq:slutsky}) 右端的替代项在 $(p,w)$ 处取值后相差 $O(\mathrm dp^2)$ 的项。差异的二阶性质（Slutsky 补偿矩阵不对称、Hicks 补偿矩阵对称）是指数理论中 Living Cost Index 与 Laspeyres 型指数分歧的根源。
\end{remark}

\begin{remark}[自身价格情形的符号分析]
取 $i=j$：
\[
 \frac{\partial x_i}{\partial p_i}
 =\underbrace{\frac{\partial h_i}{\partial p_i}}_{\le\,0}-x_i\,\frac{\partial x_i}{\partial w}.
\]
\begin{itemize}
\item 若 $x_i$ 是\textbf{正常品}（$\partial x_i/\partial w>0$）：两项同为非正，故 $\partial x_i/\partial p_i<0$ 严格——正常品必是普通品（价格升、需求减）。
\item 若 $x_i$ 是\textbf{劣等品}（$\partial x_i/\partial w<0$）：第二项为正，与第一项对抗。仅当财富效应在数值上压倒替代效应（$x_i\,|\partial x_i/\partial w|>|\partial h_i/\partial p_i|$，要求商品既劣等又占预算大头）时才有 $\partial x_i/\partial p_i>0$，即 Giffen 品。Giffen 品因此是劣等品的特例，且是“预算份额很大的劣等品”这一相当苛刻的特例。
\end{itemize}
这一符号逻辑的完整展开与经典案例见下一节。
\end{remark}

\subsubsection*{交叉价格效应的符号分析}

自身价格效应的符号已如上述；交叉项 $\partial x_i/\partial p_j$（$i\ne j$）则\textbf{没有}普适符号。Slutsky 方程把它拆成两个方向都可能变的量：
\[
 \frac{\partial x_i}{\partial p_j}
 =\underbrace{S_{ij}}_{\text{符号自由（}n\ge3\text{ 时）}}
 -\underbrace{x_j\,x_{i,w}}_{\text{符号自由}} .
\]
逐情形讨论：
\begin{itemize}
\item \textbf{Hicksian 替代 $+$ $i$ 正常}（$S_{ij}>0$，$x_{i,w}>0$）：总效应必为正——$j$ 是 $i$ 的总替代品。这是最常见的组合（如茶对咖啡）。
\item \textbf{Hicksian 互补 $+$ $i$ 正常}（$S_{ij}<0$，$x_{i,w}>0$）：财富效应 $-x_jx_{i,w}<0$ 与替代效应同向为负，总效应必为负——$j$ 是 $i$ 的总互补品（如糖对咖啡：糖涨价，少喝咖啡）。
\item \textbf{Hicksian 替代 $+$ $i$ 劣等}（$S_{ij}>0$，$x_{i,w}<0$）或 \textbf{Hicksian 互补 $+$ $i$ 劣等}：两项方向相反，总效应符号不定，取决于 $x_j$（$j$ 的预算份额）与敏感度的乘积大小。
\end{itemize}
由此可见“总替代/总互补”的不对称性来源：$\partial x_i/\partial p_j$ 含 $x_j x_{i,w}$、$\partial x_j/\partial p_i$ 含 $x_i x_{j,w}$，两个财富项一般不等；而补偿项 $S_{ij}=S_{ji}$ 恒对称。分类学细节（两商品世界的必然替代、CD 的中性总效应）见下一节。

\begin{example}[完全替代：可微性失效的边界情形]\label{ex:perf-sub}
设 $u=x_1+x_2$。UMP 的解是角点：$p_1<p_2$ 时 $x=(w/p_1,0)$。EMP 同样极端：
\[
 e(p,\bar u)=\bar u\min\{p_1,p_2\},\qquad
 h(p,\bar u)=
 \begin{cases}
 (\bar u/p_1,0), & p_1<p_2,\\
 (0,\bar u/p_2), & p_1>p_2,
 \end{cases}
\]
而在 $p_1=p_2$ 处，任何 $x_1+x_2=\bar u$ 的束都是最优——\textbf{$h$ 多值、$e$ 有折点（kink）、不可微}。Shephard 引理与整个 Slutsky 矩阵在折点处失效：$D_p h$ 在 $p_1=p_2$ 两侧是两个不同的常数矩阵，中间发生跳跃（$h_1$ 从 $\bar u/p_1$ 跳到 $0$）。教训有二：其一，定理~\ref{thm:slutsky-props} 中“$e$ 二阶连续可微”不是装饰性假设——偏好凸性只保证 $e$ 凹与 $h$ 存在，可微性需要无差异曲线本身光滑（无折点）；其二，折点恰恰对应\emph{离散选择}（只买较便宜那种）——计量中处理角点与离散选择需求（如 logit 型）时，正是这一情形的现代推广。
\end{example}

图~\ref{fig:slutsky-decomp} 把三点分解画成几何：价格下降使预算线绕纵轴截距外旋；先沿原无差异曲线滑到补偿切点（替代效应），再平行移到新预算线上的最终点（财富效应）。

\begin{figure}[H]
\centering
\begin{tikzpicture}[scale=.88]
  \draw[axisline] (0,0)--(7.4,0) node[below]{$x_1$};
  \draw[axisline] (0,0)--(0,4.65) node[left]{$x_2$};
  % 原预算线 x2=4-x1
  \draw[thick] (0,4)--(4,0);
  \node[figlabel] at (3.62,4.28) {原预算线};
  % 新预算线 x2=4-x1/2
  \draw[thick,dashed] (0,4)--(7.2,0.4);
  \node[figlabel] at (6.05,4.4) {新预算线};
  % 补偿预算线 x2=2.83-x1/2
  \draw[thick,dotted] (0,2.83)--(5.66,0);
  \node[figlabel] at (5.15,3.28) {补偿预算线};
  % 无差异曲线 u1: xy=4
  \draw[curveA,domain=1.0:5.3,smooth] plot (\x,{4/\x});
  \node[figlabel] at (4.7,1.28) {$u_1$};
  % 无差异曲线 u2: xy=8
  \draw[curveB,domain=1.75:6.9,smooth] plot (\x,{8/\x});
  \node[figlabel] at (6.35,1.62) {$u_2$};
  % 三点
  \node[nodeblue] at (2,2) {};
  \node[figlabel,left] at (1.9,2.06) {$x$};
  \node[nodewhite] at (2.83,1.41) {};
  \node[figlabel,left] at (2.68,1.36) {$h$};
  \node[nodeblue] at (4,2) {};
  \node[figlabel,right] at (4.1,2.05) {$x'$};
  % 分解箭头
  \draw[arrblue] (2.25,1.9) to[bend left=30] (2.6,1.5);
  \node[figlabel] at (2.95,1.05) {替代效应};
  \draw[arrblue] (3.15,1.65) to[bend left=30] (3.7,1.9);
  \node[figlabel] at (3.8,1.25) {财富效应};
\end{tikzpicture}
\caption{价格下降的 Slutsky 三点分解}
\figexplan{商品 1 价格下降（预算线由实线外旋为虚线，纵截距不变）。$x$ 为原最优束；补偿预算线（点线，与新预算线平行、与原无差异曲线 $u_1$ 相切）给出补偿点 $h$；新预算线与更高无差异曲线 $u_2$ 的切点为最终束 $x'$。$x\to h$ 沿 $u_1$ 滑动，是替代效应（$x_1$ 增加）；$h\to x'$ 平行外移，是财富效应（两种商品都是正常品时 $x_1$ 进一步增加）。}
\label{fig:slutsky-decomp}
\end{figure}

\begin{remark}[Slutsky 方程是可检验命题]
注意方程两边都是\textbf{可观察}的量：$\partial x_i/\partial p_j$、$x_j$ 与 $\partial x_i/\partial w$ 都能由需求数据估出，而方程断言它们组合出的矩阵 $S=\nabla_p x+(\nabla_w x)x^{\top}$ 必然对称、负半定、行和为零。这是消费者理论对横截面与面板需求数据最核心的可检验限制（也是需求系统计量估计中约束设定的来源）。反过来，\emph{可积性问题}问：给定一个满足这些限制的矩阵 $S$，能否反解出产生它的偏好？答案是肯定的——这正是 $S=\nabla^2_p e$ 把 $S$ “积分”回支出函数的路线，对偶理论在这里与实证需求分析会师。
\end{remark}

\begin{remark}[副产品：真实生活成本指数]
Slutsky 机制还有一个重要的福利应用。要在保持效用 $\bar u$ 不变的意义下度量从价格 $p^0$ 到 $p^1$ 的生活成本变化，正确的指数是 \textbf{Konüs 真实生活成本指数}
\[
 \mathrm{COLI}(p^0,p^1;\bar u)=\frac{e(p^1,\bar u)}{e(p^0,\bar u)} :
\]
分子分母都是达到\emph{同一}效用水平的（最小）支出，比值即纯价格效应。统计机构用的 Laspeyres 指数 $P_L=p^1\cdot x(p^0,w)/w$ 以\emph{基期束}为权重，而 $e(p^1,\bar u)\le p^1\cdot h(p^0,\bar u)$（最优束成本不超过任一可行束成本，取 $\bar u=v(p^0,w)$ 即基期束）给出 $P_L\ge\mathrm{COLI}$：Laspeyres 系统性\emph{高估}生活成本（忽略替代）。对偶地可以证明 Paasche 指数系统性低估。这组不等式是支出函数凹性（确切说是定义式的最小性）的直接推论—— Slutsky 理论之外，对偶方法免费赠送的福利经济学结论。
\end{remark}

% ============================================================
% 第五讲 5.8：财富效应、商品分类与综合例题
% ============================================================
\subsection{财富效应与 Engel 曲线}

Slutsky 方程把价格效应拆成替代效应与财富效应，两块的符号由不同力量决定。本节先按\textbf{财富效应}的方向给商品分类，再按\textbf{自身价格效应}的方向分类，并用 Slutsky 方程把两套分类焊起来。

\begin{definition}[正常品与劣等品]
固定 $(p,w)$，设 $x_i(p,w)$ 可微。若 $\partial x_i/\partial w>0$（财富增加消费增加），称商品 $i$ 在 $(p,w)$ 处为\textbf{正常品}（normal good）；若 $\partial x_i/\partial w<0$，称为\textbf{劣等品}（inferior good）。若对所有 $w$ 都是正常品，称全局正常品。
\end{definition}

劣等品并非“质量差”，而是\textbf{随财富增长被排挤出消费束}的商品：便宜的主食、低档交通方式等——穷时多买，富时少买。正常与劣等是\emph{局部}性质：同一商品可能在低收入区间正常、高收入区间劣等。

财富 $w$ 变动（价格不动）时，最优束 $x(p,w)$ 在商品空间中描出的轨迹称为\textbf{收入--消费曲线}（income--consumption curve, ICC）；把其中第 $i$ 个分量对 $w$ 作图，即得第 $i$ 种商品的 \textbf{Engel 曲线} $w\mapsto x_i(p,w)$。

\begin{figure}[H]
\centering
\begin{tikzpicture}[scale=.9]
  \draw[axisline] (0,0)--(6.1,0) node[below]{$x_1$};
  \draw[axisline] (0,0)--(0,4.75) node[left]{$x_2$};
  % 三条预算线 x+y=w, w=2.4,3.6,4.8（p=(1,1)）
  \draw[curveC] (0,2.4)--(2.4,0);
  \draw[thick] (0,3.6)--(3.6,0);
  \draw[curveB] (0,4.8)--(4.8,0);
  \node[figlabel] at (1.05,2.62) {$w_1$};
  \node[figlabel] at (2.25,3.82) {$w_2$};
  \node[figlabel] at (3.55,5.02) {$w_3$};
  % 三条无差异曲线（u=x^{1/3}y^{2/3} 型，切点 (0.8,1.6),(1.2,2.4),(1.6,3.2)）
  \draw[curveA,domain=0.55:3.1,smooth] plot (\x,{1.431/sqrt(\x)});
  \draw[curveA,domain=0.75:4.3,smooth] plot (\x,{2.629/sqrt(\x)});
  \draw[curveA,domain=0.95:5.4,smooth] plot (\x,{4.048/sqrt(\x)});
  % 收入-消费曲线：斜率 2 的射线
  \draw[thick,-{Stealth}] (0.15,0.3)--(2.0,4.0);
  \node[figlabel,above left] at (1.95,3.95) {ICC};
  \node[nodewhite] at (0.8,1.6) {};
  \node[figlabel,right] at (0.87,1.55) {$x^{(1)}$};
  \node[nodewhite] at (1.2,2.4) {};
  \node[figlabel,right] at (1.27,2.35) {$x^{(2)}$};
  \node[nodewhite] at (1.6,3.2) {};
  \node[figlabel,right] at (1.67,3.15) {$x^{(3)}$};
\end{tikzpicture}
\caption{财富增加：收入--消费曲线与 Engel 曲线}
\figexplan{价格固定为 $p=(1,1)$，财富由 $w_1$ 增至 $w_3$，预算线（两条虚点线夹一条实线）平行外移，最优束沿切点 $x^{(1)}\to x^{(2)}\to x^{(3)}$ 移动，其轨迹即收入--消费曲线 ICC（带箭头粗线）。此处两种商品都是正常品：ICC 单调向右上。把 $x_1$ 分量对 $w$ 作图即第一种商品的 Engel 曲线；对位似偏好（本图效用为 $x_1^{1/3}x_2^{2/3}$），ICC 是过原点的射线，Engel 曲线是过原点的直线。}
\label{fig:engel}
\end{figure}

\begin{remark}[Engel 曲线的形状与位似偏好]
位似偏好（边际替代率只依赖消费比例）下 $x(p,w)=w\cdot\phi(p)$，Engel 曲线是过原点的直线；其中 Cobb--Douglas 给出 $x_i=\alpha_i w/p_i$，斜率即支出份额除以价格。恩格尔本人观察到食品支出份额随财富上升而下降（Engel 定律），即食品的财富弹性小于 $1$——效用函数层面常用食品份额内嵌的 PIGLOG/QUAIS 形式来拟合这一特征。
\end{remark}

\subsection{自身价格效应：普通品与 Giffen 品}

\begin{definition}[普通品与 Giffen 品]
若 $\partial x_i/\partial p_i<0$（自身涨价、需求减少），称商品 $i$ 为\textbf{普通品}（ordinary good）；若 $\partial x_i/\partial p_i>0$（自身涨价、需求反而增加），称为\textbf{Giffen 品}。
\end{definition}

两套分类通过 Slutsky 方程严格相连。

\begin{proposition}[正常品必为普通品]\label{prop:normal-ordinary}
若商品 $i$ 在 $(p,w)$ 处正常（$\partial x_i/\partial w>0$）且 $x_i(p,w)>0$，则 $\partial x_i/\partial p_i<0$，即 $i$ 是普通品。
\end{proposition}
\begin{proof}
由 Slutsky 方程（定理~\ref{thm:slutsky-eq}）取 $i=j$：
\[
 \frac{\partial x_i}{\partial p_i}
 =\frac{\partial h_i}{\partial p_i}-x_i\,\frac{\partial x_i}{\partial w}.
\]
右边第一项由补偿需求定律（推论~\ref{cor:diag}）非正；第二项中 $x_i>0$、$\partial x_i/\partial w>0$，故 $-x_i\,\partial x_i/\partial w<0$。两项之和严格为负。
\end{proof}

\begin{corollary}[Giffen 品必为劣等品]
若商品 $i$ 是 Giffen 品（$\partial x_i/\partial p_i>0$）且 $x_i>0$，则 $\partial x_i/\partial w<0$。
\end{corollary}
\begin{proof}
命题~\ref{prop:normal-ordinary} 的逆否命题：若 $\partial x_i/\partial w\ge0$ 且 $x_i>0$，则 $\partial x_i/\partial p_i<0$，不可能是 Giffen 品。
\end{proof}

因此 Giffen 品必须同时满足苛刻的三个条件：(i) 劣等品；(ii) 在预算中占\textbf{很大}份额（财富效应 $\propto x_i$ 才够大）；(iii) 替代余地小（补偿效应弱）。三者叠加使“涨价使你更穷 $\Rightarrow$ 更多依赖这种廉价主食”的机制压倒替代机制。

\begin{remark}[爱尔兰土豆故事的评述]
教科书的经典叙事是：$19$ 世纪爱尔兰饥荒中土豆价格上涨，贫穷家庭实际购买力锐减，被迫\emph{减少}肉与面包、\emph{增加}土豆消费，土豆遂成 Giffen 品。评述有三点。其一，该故事主要是 Marshall 的转述，Giffen 本人的著述中并无土豆的明确记载，作为史实并不可靠。其二，机制本身在逻辑上成立且已被实验支持：Jensen 与 Miller（2008）在中国湖南与甘肃对大米/面食的补贴实验中观察到了与 Giffen 行为一致的证据——对象恰是“收入不高、该主食占预算比重大、无更便宜的替代品”的家庭。其三，Giffen 行为是\textbf{局部}现象：同一商品只在贫富区间的一段内、且只在特定价格范围内表现如此；它违背的是马歇尔需求定律，而不是补偿需求定律（后者由偏好凸性保证永不可能被违反）。这再次说明把“需求定律”精确化为 Slutsky 分解的价值：可违反的从来只是财富效应那一项。
\end{remark}

\begin{figure}[H]
\centering
\begin{tikzpicture}[scale=.88]
  \draw[axisline] (0,0)--(7.2,0) node[below]{$x_1$};
  \draw[axisline] (0,0)--(0,4.75) node[left]{$x_2$};
  % 原预算线 (0,4)-(4,0)
  \draw[thick] (0,4)--(4,0);
  \node[figlabel] at (3.6,4.3) {原预算线};
  % 新预算线 (0,4)-(8,0)，画到 (7.4,0.3)
  \draw[thick,dashed] (0,4)--(7.4,0.3);
  \node[figlabel] at (6.0,4.42) {新预算线};
  % 补偿预算线 (0,2.83)-(5.66,0)
  \draw[thick,dotted] (0,2.83)--(5.66,0);
  \node[figlabel] at (5.1,3.3) {补偿预算线};
  % 无差异曲线 u1: xy=4
  \draw[curveA,domain=1.0:5.3,smooth] plot (\x,{4/\x});
  \node[figlabel] at (4.7,1.3) {$u_1$};
  % 无差异曲线 u2：在 (1.5,3.25) 与新预算线相切（分段抛物线拼接，相切于 x=1.5）
  \draw[curveA,domain=1.15:1.5,smooth] plot (\x,{4-0.5*\x+2*(\x-1.5)*(\x-1.5)});
  \draw[curveA,domain=1.5:4.8,smooth] plot (\x,{4-0.5*\x+0.15*(\x-1.5)*(\x-1.5)});
  \node[figlabel] at (1.02,4.35) {$u_2$};
  % 三点
  \node[nodeblue] at (2,2) {};
  \node[figlabel,below] at (2,1.85) {$x$};
  \node[nodewhite] at (2.83,1.41) {};
  \node[figlabel,right] at (2.93,1.36) {$h$};
  \node[nodeblue] at (1.5,3.25) {};
  \node[figlabel,above left] at (1.42,3.32) {$x'$};
  % 分解箭头
  \draw[arrblue] (2.3,1.95) to[bend left=32] (2.62,1.5);
  \node[figlabel] at (3.3,1.05) {替代效应（$x_1\uparrow$）};
  \draw[arrblue] (2.6,1.6) to[bend left=45] (1.75,2.9);
  \node[figlabel] at (3.35,2.35) {财富效应（$x_1\downarrow\downarrow$，占优）};
\end{tikzpicture}
\caption{Giffen 品：财富效应压倒替代效应}
\figexplan{商品 1 劣等且占预算大头。$p_1$ 下降使预算线外旋（实线 $\to$ 虚线）。补偿点 $h$ 在原无差异曲线 $u_1$ 上：替代效应使 $x_1$ 由 $2$ 增至 $2.83$。但真实财富效应把消费者推向更高的无差异曲线 $u_2$ 上的 $x'$，其横坐标 $1.5$ \emph{小于}原点 $x$ 的横坐标 $2$——降价反而少买，即 Giffen 行为。注意 $u_2$ 在左侧急剧抬升：劣等品的财富效应方向与替代效应相反，且数值上占优。}
\label{fig:giffen}
\end{figure}

\subsection{交叉价格效应：总替代与 Hicksian 替代}

自身价格之外的另一组分类看\textbf{交叉}价格效应 $\partial x_i/\partial p_j$（$i\ne j$），并区分“是否已补偿”。

\begin{definition}[总替代/总互补与 Hicksian 替代/互补]
商品 $j$ 是商品 $i$ 的\textbf{总替代品}（gross substitute），若 $\partial x_i/\partial p_j>0$；\textbf{总互补品}（gross complement），若 $\partial x_i/\partial p_j<0$（均指未补偿的 Marshallian 需求）。相应地，若补偿需求满足 $\partial h_i/\partial p_j>0$，称 $j$ 是 $i$ 的 \textbf{Hicksian 替代品}；若 $\partial h_i/\partial p_j<0$，称 \textbf{Hicksian 互补品}。
\end{definition}

两套概念的关键差别在\textbf{对称性}：

\begin{remark}[总概念不对称，补偿概念对称]
总替代不必对称：$\partial x_i/\partial p_j$ 与 $\partial x_j/\partial p_i$ 可以一正一负。原因在 Slutsky 方程的财富项：
\[
 \frac{\partial x_i}{\partial p_j}
 =\frac{\partial h_i}{\partial p_j}-x_j\frac{\partial x_i}{\partial w},\qquad
 \frac{\partial x_j}{\partial p_i}
 =\frac{\partial h_j}{\partial p_i}-x_i\frac{\partial x_j}{\partial w},
\]
两个财富项一般不相等（取决于各自的 $x$ 与 $\partial x/\partial w$）。而替代项由 Slutsky 矩阵的\textbf{对称性}（定理~\ref{thm:slutsky-props}）保证 $\partial h_i/\partial p_j=\partial h_j/\partial p_i$：补偿意义上“$i$ 的替代品”与“$j$ 的替代品”是同一关系，像“邻居”一样是相互的。这使 Hicksian 分类在理论上更干净；总分类则更贴近市场上可观察的未补偿反应，二者各有用途，但不应混用。
\end{remark}

\begin{example}[Cobb--Douglas：总分类为中性，补偿分类为替代]
$u=x_1^{\alpha}x_2^{1-\alpha}$ 给出 $x_1=\alpha w/p_1$：$\partial x_1/\partial p_2=0$，两商品\textbf{既非总替代也非总互补}；但补偿意义上 $\partial h_1/\partial p_2=\alpha(1-\alpha)w/(p_1p_2)>0$（见下节例~\ref{ex:cd-slutsky}），是严格的 Hicksian 替代品。同一对商品在两套分类下结论不同，直观原因：CD 下两种商品的支出份额固定，$p_2$ 变动不改变花在商品 $1$ 上的钱（总效应为零），但固定效用后相对价格的替代机制仍然存在。
\end{example}

\begin{proposition}[两商品世界必为 Hicksian 替代品]\label{prop:two-goods}
设 $n=2$，偏好凸、$h$ 可微且 $D_p h$ 负半定、行和为零。则 $\partial h_1/\partial p_2\ge0$ 且 $\partial h_2/\partial p_1\ge0$；若 $\partial h_1/\partial p_1<0$（自身补偿效应严格为负），则两者严格为正。
\end{proposition}
\begin{proof}
对第 $1$ 行写出行和为零（Euler）：$h_{11}p_1+h_{12}p_2=0$，故
\[
 h_{12}=-\frac{p_1}{p_2}\,h_{11}\;\ge\;0,
\]
末一不等号用了对角元非正（推论~\ref{cor:diag}）。再由对称性 $h_{21}=h_{12}\ge0$。严格性直接由 $h_{11}<0$ 与上式传递。
\end{proof}

直觉：只有两种商品时，保持效用不变、商品 $1$ 涨价，钱只能省在商品 $2$ 上（或反过来）——补偿世界里的“敌人”只有对方，必然互为替代。互补品（如咖啡与糖）至少需要三种商品：商品 $1$、商品 $2$，以及涨价后把两者一起挤出的“第三方”。

\subsection{Offer 曲线}

\begin{definition}[价格--提供曲线]
固定财富 $w$ 与 $p_{-i}$，让 $p_i$ 变动，最优束 $x(p,w)$ 在商品空间中描出的轨迹称为\textbf{价格--提供曲线}（price offer curve）。
\end{definition}

提供曲线是 Slutsky 分解的“原始底片”：它把价格变化引起的\emph{全部}总效应记录成一条路径，每一点的切向即 $\nabla_p x$ 沿该方向的列。提供曲线向哪边弯，直接显示商品是普通品还是 Giffen 品、总替代还是总互补。

\begin{figure}[H]
\centering
\begin{tikzpicture}[scale=.88]
  \draw[axisline] (0,0)--(7.5,0) node[below]{$x_1$};
  \draw[axisline] (0,0)--(0,4.35) node[left]{$x_2$};
  % 三条预算线均过 (0,4)：p1=2,1,0.5
  \draw[curveC] (0,4)--(2,0);
  \node[figlabel] at (1.05,4.25) {$p_1=2$};
  \draw[thick] (0,4)--(4,0);
  \node[figlabel] at (3.05,4.25) {$p_1=1$};
  \draw[curveB] (0,4)--(7.4,0.3);
  \node[figlabel] at (5.9,4.35) {$p_1=\tfrac12$};
  % 提供曲线 x2=4/(x1+1)
  \draw[curveA,domain=0.25:7.25,smooth] plot (\x,{4/(\x+1)});
  \node[figlabel] at (5.9,1.0) {offer 曲线};
  % 一条无差异曲线过 (3,1)：x1=3-ln(x2)
  \draw[curveB,domain=0.55:3.0,smooth,variable=\t] plot ({3-ln(\t)},{\t});
  \node[figlabel] at (1.6,2.5) {$u=3$};
  % 切点
  \node[nodewhite] at (1,2) {};
  \node[figlabel,above left] at (0.9,2.1) {$x^{(2)}$};
  \node[nodewhite] at (3,1) {};
  \node[figlabel,above left] at (2.9,1.1) {$x^{(1)}$};
  \node[nodewhite] at (7,0.5) {};
  \node[figlabel,left] at (6.9,0.55) {$x^{(1/2)}$};
\end{tikzpicture}
\caption{价格--提供曲线（拟线性效用 $u=x_1+\ln x_2$）}
\figexplan{财富固定为 $w=4$、$p_2=1$，$p_1$ 由 $2$ 降到 $1$ 再降到 $1/2$，预算线绕纵截距 $(0,4)$ 外旋（点线 $\to$ 实线 $\to$ 虚线）。切点 $x^{(2)}\to x^{(1)}\to x^{(1/2)}$ 的轨迹即提供曲线（实粗线），此处为双曲线 $x_2=4/(x_1+1)$。虚线无差异曲线 $u=3$ 与中间预算线相切于 $x^{(1)}=(3,1)$，可检验切点性质。提供曲线向右下弯：$p_1$ 下降时 $x_1$ 沿曲线增加、$x_2$ 减少，商品 $1$ 是普通品，且商品 $2$ 是商品 $1$ 的总互补品（$\partial x_2/\partial p_1<0$ 对称地看即 $\partial x_2$ 随 $p_1$ 下降而下降）。}
\label{fig:offer}
\end{figure}

\subsection{综合例题}

\begin{example}[拟线性效用的完整求解与 Slutsky 验证]\label{ex:quasilinear}
设 $u(x_1,x_2)=x_1+\ln x_2$，$p\gg0$。最大化 $x_1+\ln x_2$ 于 $p_1x_1+p_2x_2=w$。
\textbf{需求。}内点一阶条件：$1=\lambda p_1$，$1/x_2=\lambda p_2$，故
\[
 x_2=\frac{p_1}{p_2},\qquad
 x_1=\frac{w-p_1}{p_2}\cdot\frac{p_2}{p_1}=\frac{w}{p_1}-1 .
\]
内点解要求 $w>p_1$（否则 $x_1=0$，全部财富买 $x_2=w/p_2$——非负约束在 $w\le p_1$ 时绑定，是角点解）。
\textbf{间接效用。}
\[
 v(p,w)=\frac{w}{p_1}-1+\ln\frac{p_1}{p_2}.
\]
Roy identity 验证：$v_w=1/p_1=\lambda$；$v_{p_1}=-w/p_1^2+1/p_1$，$v_{p_2}=-1/p_2$，故
\[
 -\frac{v_{p_1}}{v_w}=\frac{w}{p_1}-1=x_1,\qquad
 -\frac{v_{p_2}}{v_w}=\frac{p_1}{p_2}=x_2 .\checkmark
\]
\textbf{支出函数。}由对偶恒等式，令 $v(p,w)=\bar u$ 反解 $w$：
\[
 e(p,\bar u)=p_1\Bigl(\bar u+1-\ln\frac{p_1}{p_2}\Bigr).
\]
\textbf{Hicksian 需求。}Shephard 引理：
\[
 h_1=\frac{\partial e}{\partial p_1}=\bar u-\ln\frac{p_1}{p_2},\qquad
 h_2=\frac{\partial e}{\partial p_2}=\frac{p_1}{p_2}.\checkmark
\]
（$h_2$ 与 Marshallian $x_2$ 相同：商品 $2$ 无财富效应，补偿与否无差别。）
\textbf{Slutsky 验证。}先算 Marshallian 导数：$\partial x_1/\partial p_1=-w/p_1^2$，$\partial x_1/\partial w=1/p_1$，$x_1=w/p_1-1$。补偿导数：由 $h_1=\bar u-\ln p_1+\ln p_2$ 得（在 $\bar u=v$ 处）$\partial h_1/\partial p_1=-1/p_1$。代入 Slutsky 方程右端：
\[
 \frac{\partial h_1}{\partial p_1}-x_1\frac{\partial x_1}{\partial w}
 =-\frac{1}{p_1}-\Bigl(\frac{w}{p_1}-1\Bigr)\frac{1}{p_1}
 =-\frac{w}{p_1^2}
 =\frac{\partial x_1}{\partial p_1} .\checkmark
\]
交叉项：$\partial x_1/\partial p_2=0$；$\partial h_1/\partial p_2=1/p_2$，$x_2=p_1/p_2$，$\partial x_1/\partial w=1/p_1$，故右端 $=1/p_2-(p_1/p_2)(1/p_1)=0$ $\checkmark$。对称性：$\partial h_1/\partial p_2=1/p_2=\partial h_2/\partial p_1$ $\checkmark$。矩阵性质：由 $h_2=p_1/p_2$ 补齐第二行，得
\[
 S=\begin{pmatrix}
 -1/p_1 & 1/p_2\\
 1/p_2 & -p_1/p_2^2
 \end{pmatrix},
\]
行列式 $=(1/p_1)(p_1/p_2^2)-1/p_2^2=0$、对角元均为负，故特征值为 $0$ 与负数，负半定 $\checkmark$；两行与 $p$ 内积均为零，$Sp=0$ $\checkmark$。
\end{example}

\begin{example}[Cobb--Douglas 的 Slutsky 矩阵]\label{ex:cd-slutsky}
设 $u=x_1^{\alpha}x_2^{1-\alpha}$，$0<\alpha<1$。已知
\[
 x_1=\frac{\alpha w}{p_1},\qquad x_2=\frac{(1-\alpha)w}{p_2},\qquad
 e(p,\bar u)=\frac{\bar u\,p_1^{\alpha}p_2^{1-\alpha}}{\alpha^{\alpha}(1-\alpha)^{1-\alpha}} .
\]
\textbf{补偿需求。}$h_1=\partial e/\partial p_1=\alpha e/p_1$，$h_2=(1-\alpha)e/p_2$。
\textbf{Slutsky 矩阵。}在 $\bar u=v(p,w)$（此时 $e=w$）处：
\[
 \frac{\partial h_1}{\partial p_1}
 =-\alpha(1-\alpha)\frac{e}{p_1^2}\Big|_{e=w}
 =-\frac{\alpha(1-\alpha)w}{p_1^2},\qquad
 \frac{\partial h_1}{\partial p_2}
 =\frac{\alpha(1-\alpha)w}{p_1p_2},
\]
由对称结构（$\alpha\leftrightarrow1-\alpha$、$p_1\leftrightarrow p_2$）：
\[
 \frac{\partial h_2}{\partial p_1}=\frac{\alpha(1-\alpha)w}{p_1p_2},\qquad
 \frac{\partial h_2}{\partial p_2}=-\frac{\alpha(1-\alpha)w}{p_2^2}.
\]
于是
\[
 S(p,w)=-\alpha(1-\alpha)w
 \begin{pmatrix}
 1/p_1^2 & -1/(p_1p_2)\\
 -1/(p_1p_2) & 1/p_2^2
 \end{pmatrix}
 =-\alpha(1-\alpha)w\;q\,q^{\top},\qquad
 q=\Bigl(\frac1{p_1},-\frac1{p_2}\Bigr)^{\top}.
\]
\textbf{三大性质验证。}
对称性：$S_{12}=S_{21}$，显然。负半定：对任意 $z\in\R^2$，
\[
 z^{\top}Sz=-\alpha(1-\alpha)w\,(q^{\top}z)^2\le0,
\]
且 $S$ 秩为 $1$，特征值为 $0$ 与 $-\alpha(1-\alpha)w\|q\|^2<0$。行和为零：$q^{\top}p=p_1/p_1-p_2/p_2=0$，故 $Sp=-\alpha(1-\alpha)w\,q\,(q^{\top}p)=0$。
\textbf{Slutsky 方程验证。}自身价格：
\[
 \frac{\partial h_1}{\partial p_1}-x_1\frac{\partial x_1}{\partial w}
 =-\frac{\alpha(1-\alpha)w}{p_1^2}-\frac{\alpha w}{p_1}\cdot\frac{\alpha}{p_1}
 =-\frac{\alpha w}{p_1^2}=\frac{\partial x_1}{\partial p_1}\;\checkmark
\]
交叉价格：$\partial x_1/\partial p_2=0$，而
\[
 \frac{\partial h_1}{\partial p_2}-x_2\frac{\partial x_1}{\partial w}
 =\frac{\alpha(1-\alpha)w}{p_1p_2}-\frac{(1-\alpha)w}{p_2}\cdot\frac{\alpha}{p_1}=0\;\checkmark.
\]
注意结论：CD 的 $S$ 负半定但不满秩（一行乘比例即另一行），体现两商品世界中补偿替代关系的秩约束；而 Marshallian 交叉效应为零，再次对照出总分类与补偿分类的差别。
\end{example}

% ============================================================
% 第五讲章末小结：完整逻辑链、一页检查清单与常见陷阱
% ============================================================
\section{本章小结}

本章从偏好出发，经由表示、最优化、对偶与分解，走完了消费者需求理论的全部主干。
把九个小节串成一条逻辑链：

\begin{enumerate}
\item \textbf{表示}（\S 5.1）：Debreu 定理保证完备、传递、连续的偏好可用连续效用函数表示；
      字典序偏好证明连续性不可省略——它完备、传递、局部非饱和，却不存在任何实值表示。
\item \textbf{概念与存在性}（\S 5.2）：Marshallian 需求 $x(p,w)$ 固定财富问“选什么”，
      Hicksian 需求 $h(p,\bar u)$ 固定效用问“最少花多少”；紧预算集加连续效用
      （Weierstrass 定理）保证最大值存在——$\sup$ 不必达到，$\max$ 必须达到。
\item \textbf{值函数}（\S 5.3）：间接效用 $v(p,w)$ 把最优值封装为价格与财富的函数，
      零次齐次与拟凸性是预算线性结构的赠品，对 $w$ 严格递增由局部非饱和买到，
      连续性由 Berge 最大值定理买到。
\item \textbf{求解}（\S 5.4）：Kuhn--Tucker 四组条件——可行性、驻点、乘子非负、互补松弛——
      逐条排除错误候选；凸问题下升级为充分必要；等式拉格朗日法只是其“先验内点、先验绑定”的特例。
\item \textbf{边际条件}（\S 5.5）：MRS 是序数概念，内点最优处 MRS 等于价格比；
      可分性把 MRS 的依赖限制在组内，换来实证上的两阶段预算。
\item \textbf{包络}（\S 5.6）：值函数对参数的导数只需对目标直接求偏导（间接项被一阶条件消去），
      由此得 $v_w=\lambda$、$v_{p_i}=-\lambda x_i$，两式相除即 Roy identity。
\item \textbf{对偶}（\S 5.7）：支出最小化与成本最小化同构，支出函数对价格一次齐次、凹、对效用非减，
      Shephard 引理 $h_i=\partial e/\partial p_i$ 与四条对偶恒等式把两个问题互为镜像。
\item \textbf{分解}（\S 5.8）：Slutsky 方程把价格总效应拆为替代效应与财富效应；
      Slutsky 矩阵对称、负半定、行和为零，是需求数据对理论的核心可检验限制。
\end{enumerate}

\noindent 从偏好到 Slutsky 分解的完整逻辑链，按顺序自检：

\begin{enumerate}
\item \textbf{偏好与表示}：完备、传递、连续 $\Rightarrow$ 连续效用表示存在（Debreu）；局部非饱和 $\Rightarrow$ 预算绑定 $p\cdot x=w$；严格凸 $\Rightarrow$ 需求单值。
\item \textbf{存在性与求解}：$p\gg0$、$u$ 连续 $\Rightarrow$ $B(p,w)$ 紧、解存在；内点解满足 $u_i/u_j=p_i/p_j$（MRS $=$ 相对价格）；角点解必须用 Kuhn--Tucker（四条件 $+$ 互补松弛）。
\item \textbf{规模性质}：$x(tp,tw)=x(p,w)$、$v(tp,tw)=v(p,w)$（零次齐次、无货币幻觉）。
\item \textbf{两个问题}：UMP 固定 $w$ 最大化 $u$，得 $x(p,w)$、$v(p,w)$；EMP 固定 $\bar u$ 最小化 $p\cdot x$，得 $h(p,\bar u)$、$e(p,\bar u)$。EMP 与成本最小化同构。
\item \textbf{包络定理}：值函数对参数的导数 $=$ 拉格朗日函数直接对参数求偏导（间接项被一阶条件消去）。由此：$v_w=\lambda$（财富的边际效用）、$v_{p_i}=-\lambda x_i$（涨价一元 $\equiv$ 损失 $x_i$ 元财富）。
\item \textbf{两个引理}：Roy identity $x_i=-v_{p_i}/v_w$（由 $v$ 得 $x$）；Shephard 引理 $h_i=\partial e/\partial p_i$（由 $e$ 得 $h$）。
\item \textbf{对偶恒等式}：$e(p,v)=w$、$h(p,v)=x$、$v(p,e)=\bar u$、$x(p,e)=h$——两个问题互为镜像，知一得四。
\item \textbf{支出函数性质}：对 $p$ 一次齐次、凹、对 $\bar u$ 非减；$h$ 对 $p$ 零次齐次。
\item \textbf{Slutsky 矩阵} $S=D_ph=\nabla^2_pe$：对称（Young 定理）、负半定（$e$ 凹）、行和为零（Euler 定理）；对角元 $\le0$ 即补偿需求定律。
\item \textbf{Slutsky 方程}：$\partial x_i/\partial p_j=\partial h_i/\partial p_j-x_j\,\partial x_i/\partial w$；总效应 $=$ 替代效应 $+$ 财富效应。推论：正常品必为普通品；Giffen 品必为劣等品（且须占预算大份额）。
\item \textbf{分类口径}：总替代/互补（Marshallian，不对称）对 Hicksian 替代/互补（补偿，对称）；两商品世界必为 Hicksian 替代品。
\item \textbf{可检验性}：由可观察的 $\nabla_px$、$\nabla_wx$、$x$ 组装出的 $S$ 必须对称、负半定、行和为零——这是需求数据对理论的核心检验（可积性问题是其逆）。
\end{enumerate}

\medskip
\noindent\textbf{三个常见陷阱。}
其一，$v(p,w)$ 对 $(p,w)$ 是\emph{拟凸}而非拟凹——上水平集一般不凸，
初学阶段最常见的混淆即在此（参见 \S 5.3 的反例）。
其二，Kuhn--Tucker 条件只是必要的——满足四条件的候选点可能是局部极小或鞍点，
只有凸问题（约束集凸、目标适当凹/拟凹）下才升级为充分必要。
其三，等式拉格朗日法在角点情形系统性失效——使用它必须先验内点与先验约束绑定，
否则必须回到 Kuhn--Tucker 并由 $\mu_i>0$ 解释零消费。

\end{document}
